\documentclass[11pt,letterpaper]{article}

\usepackage[T1]{fontenc}
\usepackage[utf8]{inputenc}
\usepackage{lmodern}
\usepackage{microtype}
\usepackage[margin=0.8in]{geometry}
\usepackage{amsmath,amssymb}
\usepackage{array,booktabs,longtable}
\usepackage{enumitem}
\usepackage{xcolor}
\usepackage{listings}
\usepackage[hidelinks]{hyperref}

% Fix-design fragments cite v1.0 manuscript labels with \cref; in this
% document those labels do not exist, so render them literally as label
% names (the honest reading: they are addresses into the tagged source).
% \spectt is a line-breakable teletype (breaks at : - / .) via url.
\usepackage{url}
\DeclareUrlCommand\spectt{\urlstyle{tt}}
\newcommand{\cref}[1]{\spectt{#1}}
\newcommand{\Cref}[1]{\spectt{#1}}

\lstdefinestyle{spec}{
  basicstyle=\ttfamily\small,
  columns=fullflexible,
  keepspaces=true,
  frame=single,
  rulecolor=\color{black!25},
  breaklines=true
}
\lstset{style=spec}

\setlist{nosep}
\raggedbottom
\emergencystretch=3em
\setcounter{secnumdepth}{3}
\setcounter{tocdepth}{2}

\title{\bfseries Gap Analysis and Completion Specification\\
\Large for the \emph{Handbook of Formal Logic}}
\author{Campaign document --- generated from the audit registers}
\date{10 August 2026 \textperiodcentered\ registers at \texttt{spec/registers/}}

\begin{document}
\maketitle

\begin{abstract}
\noindent
This document is the completion specification for transforming
\emph{Formal Logic: Language, Semantics, Deduction, and Metatheory}
(v1.0, 9 August 2026) into a comprehensive CRC-style handbook of logic.
It records every verified defect and gap from three exhaustive full-text
audits of the 14{,}780-line manuscript, fixes the target architecture and
editorial standards, and lays out the phased campaign.  The registers in
\texttt{spec/registers/} are the single source of truth; every table in
this document is generated from them by \texttt{tools/csv2tex.py} and is
regenerated, never hand-edited.
\end{abstract}

\tableofcontents

\section{Baseline and how to use this document}
\label{sec:baseline}

\paragraph{Source identification.}
The audited manuscript is
\texttt{source/formal-logic-reference.tex} at git tag
\texttt{v1.0-audit} (commit \texttt{676d170}), 14{,}780 lines,
SHA-256
\texttt{81bf3c0a\-b12863d7\-d89e1b44\-648ccc1d\-63832dc6\-8ab8e32b\-84aa9a07\-00cd6de1}.
\emph{Every line number in this document and in the registers refers to
that revision.}  Line references do not track the working tree; when a
register row is executed, the executing commit is recorded in the row's
status cell, and the frozen reference remains valid via the tag.

\paragraph{Audit provenance.}
Three independent full-text audits were performed on 10~August~2026,
covering lines 1--4995 (Parts I--III), 4995--10622 (Parts IV--IX), and
10622--14780 (Parts X--XI and back matter).  Each audit read its entire
segment, re-derived displayed calculations and countermodels, verified
the hypotheses of the major theorems against the literature, and ran
mechanical checks over the whole file (labels, references, citations).
The consolidated findings were reviewed and approved before this
specification was drawn up.

\paragraph{Identifier scheme.}
Every finding carries a permanent identifier.  Line numbers rot as soon
as editing begins; identifiers plus the frozen tag do not.

\begin{center}\small
\begin{tabular}{@{}llp{3.4in}@{}}
\toprule
\textbf{Prefix} & \textbf{Register} & \textbf{Class} \\
\midrule
R-\emph{nnn} & \texttt{defects.csv} & rigor defects (mathematical / proof content) \\
S-\emph{nnn} & \texttt{defects.csv} & standards defects (status policy, notation, citation discipline) \\
L-\emph{nnn} & \texttt{defects.csv} & \LaTeX{} / production defects \\
D-P\emph{n}-\emph{nnn} & \texttt{depth-gaps.csv} & depth gaps within existing part $n$ \\
N-\emph{d}-\emph{nnn} & \texttt{new-content.csv} & new-domain items, $d\in\{a,b,c,d\}$; N-x = recorded exclusions \\
A-\emph{nnn} & \texttt{apparatus.csv} & apparatus items (index, registries, sheets, tables, figures, checks) \\
\bottomrule
\end{tabular}
\end{center}

\paragraph{Dispositions.}
Each R/S/L row carries one of four dispositions, which is what prevents
wasted work given the restructure:
\textbf{FIX-SOURCE} (mathematical repair applied to the v1.0 text in
Phase~2, because migrated text inherits it);
\textbf{RESOLVE-BY-MIGRATION} (fixed for free when the text is rewritten
into the new tree);
\textbf{OBSOLETE} (mooted by the restructure; no action ever);
\textbf{CARRY} (a requirement transferred to the new tree's
specification).

\paragraph{Status vocabulary of this document.}
Every claim about the manuscript in this document is
\emph{verified-by-audit} unless marked \emph{spot-check pending}.  The
handbook's own three-term result-status policy (complete proof / proof
sketch / cited theorem) is quoted throughout and becomes
compile-enforced in the new tree (§\ref{sec:rulings}).

\section{Executive dashboard}
\label{sec:dashboard}

\paragraph{The three fixed decisions.}
These were made by the author and define ``done'':
\begin{enumerate}[label=\textbf{D\arabic*.}]
  \item \textbf{Comprehensiveness} spans four domains: (a)~deep
        mathematical logic, (b)~philosophical logic, (c)~logic in
        computer science, (d)~algebraic and categorical logic.
  \item \textbf{Format}: restructure toward a pure CRC-style handbook ---
        topical lookup entries and tables; proofs compressed into a
        proof annex; the dependency-ordered v1.0 manuscript becomes
        verified source material and is not printed.
  \item \textbf{First deliverable}: this specification; the writing
        campaign runs against it in later sessions.
\end{enumerate}

\paragraph{Headline findings.}
(i)~The manuscript's mathematics is largely sound --- every re-derived
countermodel and calculation checked out --- but roughly a dozen proofs
have genuine defects, the worst being that the book does not prove
$\mathrm{Con}(\mathsf{ZFC})\Rightarrow\mathrm{Con}(\mathsf{ZFC}+\neg\mathrm{CH})$
(R-001).
(ii)~The book's signature status discipline is applied in only about a
quarter of the text (48 of 60 status labels sit in lines 5000--7999).
(iii)~Whole fields required by decision~D1 are absent (registers
N-a\ldots N-d).
(iv)~The reference apparatus is far below handbook grade: the index is
$\sim$308 auto-generated theorem titles, there are no figures in 14{,}780
lines, rule sheets cover 3 of $\sim$10 calculi, and the notation
register covers none of Part~X.
(v)~The back matter's Exercise Feedback chapter is the exception ---
mechanically verified complete (112 = 112) --- and its discipline is
institutionalized by ruling~\ref{ruling:annex}.

\medskip
\noindent Register totals: 88 defects, 141 depth gaps, 185 new-content rows (170 in-scope, 15 recorded exclusions), 201 apparatus items --- 615 register rows in all.

\medskip

\noindent\begingroup\scriptsize
\begin{tabular}[t]{@{}lr@{}}
\multicolumn{2}{@{}l@{}}{\textbf{Defects by class}}\\
\midrule
production & 20 \\
rigor & 42 \\
standards & 26 \\
\midrule total & 88 \\
\end{tabular}
\quad
\begin{tabular}[t]{@{}lr@{}}
\multicolumn{2}{@{}l@{}}{\textbf{Defects by priority}}\\
\midrule
P0 & 13 \\
P1 & 54 \\
P2 & 21 \\
\midrule total & 88 \\
\end{tabular}
\quad
\begin{tabular}[t]{@{}lr@{}}
\multicolumn{2}{@{}l@{}}{\textbf{Defects by disposition}}\\
\midrule
CARRY & 13 \\
FIX-SOURCE & 37 \\
OBSOLETE & 11 \\
RESOLVE-BY-MIGRATION & 27 \\
\midrule total & 88 \\
\end{tabular}

\medskip

\begin{tabular}[t]{@{}lr@{}}
\multicolumn{2}{@{}l@{}}{\textbf{Depth gaps by treatment}}\\
\midrule
E & 56 \\
P & 6 \\
S & 68 \\
T & 11 \\
\midrule total & 141 \\
\end{tabular}
\quad
\begin{tabular}[t]{@{}lr@{}}
\multicolumn{2}{@{}l@{}}{\textbf{Depth gaps by priority}}\\
\midrule
could & 16 \\
must & 58 \\
should & 67 \\
\midrule total & 141 \\
\end{tabular}
\quad
\begin{tabular}[t]{@{}lr@{}}
\multicolumn{2}{@{}l@{}}{\textbf{New content by domain}}\\
\midrule
a & 71 \\
b & 35 \\
c & 40 \\
d & 24 \\
x & 15 \\
\midrule total & 185 \\
\end{tabular}

\medskip

\begin{tabular}[t]{@{}lr@{}}
\multicolumn{2}{@{}l@{}}{\textbf{New content by treatment}}\\
\midrule
E & 40 \\
P & 12 \\
S & 75 \\
T & 43 \\
\midrule total & 170 \\
\end{tabular}
\quad
\begin{tabular}[t]{@{}lr@{}}
\multicolumn{2}{@{}l@{}}{\textbf{Apparatus by category}}\\
\midrule
bibliography & 3 \\
check & 11 \\
environment & 7 \\
figure & 40 \\
index & 6 \\
navigation & 7 \\
rule-sheet & 23 \\
symbol-registry & 13 \\
table & 91 \\
\midrule total & 201 \\
\end{tabular}
\endgroup


\section{Defect register}
\label{sec:defects}

\subsection{Priorities}
\textbf{P0} --- mathematical falsity or unsoundness: the statement as
printed is wrong, the proof as printed is invalid, or the source fails
to compile.  \textbf{P1} --- violates the book's own declared policy
(the three-term status policy, the one-symbol-one-role rule, the
primary-source rule, the rule-display requirements).  \textbf{P2} ---
polish.  All P0 items and load-bearing P1 proof gaps are dispositioned
FIX-SOURCE and repaired in Phase~2; fix designs for them follow in
§\ref{sec:fixdesigns}.

\subsection{The register}
Projection of \texttt{registers/defects.csv} (full columns --- part,
acceptance criterion, effort --- live in the CSV).

\begingroup\scriptsize\setlength{\tabcolsep}{3pt}
\begin{longtable}{@{}p{0.45in}p{0.25in}p{0.7in}p{2.2in}p{1.75in}p{0.45in}p{0.16in}p{0.34in}@{}}
\toprule
\textbf{ID} & \textbf{Pri} & \textbf{v1.0 lines} & \textbf{Finding} & \textbf{Required fix} & \textbf{Disp} & \textbf{Ph} & \textbf{St} \\
\midrule
\endfirsthead
\toprule
\textbf{ID} & \textbf{Pri} & \textbf{v1.0 lines} & \textbf{Finding} & \textbf{Required fix} & \textbf{Disp} & \textbf{Ph} & \textbf{St} \\
\midrule
\endhead
\midrule\multicolumn{8}{r}{\emph{continued}}\\\endfoot
\bottomrule\endlastfoot
R-001 & P0 & 9062--9133 & CH-independence treatment does not prove \(\mathsf{Con}(\mathrm{ZFC})\Rightarrow\mathsf{Con}(\mathrm{ZFC}+\neg\mathrm{CH})\): the sketch uses the countable-transitive-model presentation that the adjacent Warning (9107) says does not suffice, and the nonderivability step invokes soundness over two structures never produced, so the theorem statement (9066) overclaims. & Rewrite per the book's own Warning: either present relative consistency with honest cited status or give the correct syntactic route (forcing over finite fragments / ctm with the inaccessibility caveat), deleting the two-structures soundness justification. & FIX & 2 & fixed (4201fbf) \\
R-002 & P0 & 7011--7031 vs 6567--6585, 7440--7454 & Coding inconsistency: the text asserts iterated Cantor pairing is not the coding used for bounded verification while the fixed machine numbering is built from exactly that pairing, so \(\mathsf{Comp}(e,x,s)\) lacks the bounded extraction from \(e\) it needs (the beta-function is applied only to histories), and the gap propagates to the \(\Sigma^0_1\) status of \(\mathsf{Prov}_T\). & Adopt a single coding with a proved bounded-extraction lemma covering machine codes as well as histories, and re-derive the \(\Sigma^0_1\) form of \(\mathsf{Comp}\) and \(\mathsf{Prov}_T\) from it. & FIX & 2 & fixed (563da95) \\
R-003 & P0 & 7666--7699 & Rosser sketch uses the object symbol \(\leq\), never defined or registered anywhere (only \(<\) at 6924), and silently uses two unstated Q-lemmas (bounded-case enumeration and the dichotomy \(y\leq\bar n\lor\bar n\leq y\)) that carry the proof's only nontrivial content. & Define and register object \(\leq\) (or eliminate it in favor of \(<\)) and state and prove or cite both Q-lemmas as labeled results before the sketch uses them. & FIX & 2 & fixed (563da95) \\
R-004 & P0 & 7154--7165 & The jump is defined via halting although the numbering makes \(\varphi_e(e)\) undefined on halt-without-valid-output, mismatching \(K\), and the jump theorem is asserted in unlabeled prose without proof or citation while load-bearing for hierarchy strictness (7180--7189). & Align the halting predicate with the numbering's convergence notion (or adjust the numbering) and promote the jump theorem to a labeled result with status and proof or citation. & FIX & 2 & fixed (563da95) \\
R-005 & P0 & 6564--6593 & The enumeration theorem for the program numbering (every partial computable function occurs in it) is asserted in unlabeled prose without proof while load-bearing for Rice, the recursion theorem, and Post. & State the enumeration theorem as a labeled result with a proof (or precise citation) tied to the constructed numbering. & FIX & 2 & fixed (563da95) \\
R-006 & P0 & 6952--6971 & The closure-class definition of the arithmetical hierarchy (with bounded-quantifier closure) is never proved equivalent to the prefix-normal-form characterization; the required bounded-collection lemma is missing. & State and prove the bounded-collection lemma and derive the closure-class/prenex equivalence as a labeled theorem. & FIX & 2 & fixed (563da95) \\
R-007 & P0 & 4420--4446 & Prenex normal form theorem rests on an unstated and unproved first-order NNF existence theorem, unproved existence of rectified alpha-variants, and an unstated rectification-preservation invariant; the propositional NNF theorem (1899) does not transfer. & Prove first-order NNF existence and rectified-variant existence as labeled lemmas, state the preservation invariant, and run the prenex induction on them. & FIX & 2 & fixed (fe1f17d) \\
R-008 & P0 & 4480--4507 & Skolemization equisatisfiability is labeled complete but is a sketch: no induction hypothesis is formulated, the unregistered term quantifier game is used (no game semantics exists in the book), and AC is over-assumed for the choice-free direction. & Write the induction with an explicit hypothesis and cases, remove or define the game terminology, and separate the AC-using direction from the choice-free one. & FIX & 2 & fixed (fe1f17d) \\
R-009 & P0 & 3934--3985 & The Henkin--Lindenbaum maximality step silently reuses the omega-recursion argument of the propositional Lindenbaum lemma where a transfinite recursion with limit stages and a chain argument is needed (the nearby witness stage does invoke transfinite recursion properly). & Redo the maximality step by transfinite recursion with explicit successor and limit stages and the union/chain argument. & FIX & 2 & fixed (fe1f17d) \\
R-010 & P0 & 3006--3016 & The substitution/alpha lemma applies structural induction to \(\psi[z/y]\), which is not a subformula of the original formula, so the induction principle invoked does not apply. & Recast as strong induction on formula complexity via the substitution-complexity lemma, the fix the book itself uses at 2951--2958 and 3195--3199. & FIX & 2 & fixed (fe1f17d) \\
R-011 & P0 & 1958--1965 & The CNF half of the propositional normal-form theorem is procedural with no induction or termination argument and fails degenerate cases (valid or variable-free formulas) against the book's rigid CNF definition; the DNF half is exemplary. & Give the CNF half a genuine induction (or dualize the DNF argument) with the degenerate cases handled explicitly. & FIX & 2 & fixed (2465859) \\
R-012 & P0 & 9513--9548, 9555, 9575 & The cut-elimination sketch invokes a mix/multicut rule that is never displayed anywhere, states the termination measure for Cut rather than mix, and two corollaries with full proof environments rest on the sketch, violating the release checklist's own rule (14346). & Display and label the mix rule, restate the measure for the rule actually eliminated, and either complete the induction or downgrade the corollaries' status to match the sketch. & FIX & 2 & fixed (2465859) \\
R-013 & P1 & 2020--2021, 2374, 2446--2449 & ND repetition is permitted in prose but is not among the twelve displayed rules, and vacuous discharge is never licensed although the completeness proof needs it; the first-order analogue (3804--3806) handles both explicitly. & Add repetition to the displayed rule inventory (or derive it) and license vacuous discharge explicitly, matching the first-order treatment. & FIX & 2 & fixed (2465859) \\
R-014 & P1 & 2488 vs 2295; 4165 vs 3674 & LK completeness is stated only for single-formula succedents, strictly weaker than LK soundness for arbitrary sequents, and the book never closes the loop for arbitrary sequents. & State and prove LK completeness for arbitrary sequents, or scope both metatheorems to the same sequent class with an explicit remark. & FIX & 2 & fixed (fe1f17d) \\
R-015 & P1 & 5994--6015 & Model completeness of DLO is defined by quantification over all embeddings but the proof treats only substructure inclusions; the reduction step is missing. & Add the reduction from arbitrary embeddings to substructure inclusions or restate the definition to match the proof. & FIX & 2 & fixed (2465859) \\
R-016 & P1 & 6068--6110 & The back-and-forth corollary uses an extends-the-original-map clause that is available from the proof but absent from the theorem statement, which concludes only isomorphism. & Add the extension clause to the theorem statement. & FIX & 2 & fixed (2465859) \\
R-017 & P1 & 678--702 & Mixed-status theorem thm:infinite-cardinal-finite-coding: one label covers clauses cited to Jech, a clause unlocatable in the cited source, and a clause proved inline, against the status policy at 313--319. & Split into separately labeled results with per-clause status and locate or reprove the unlocatable clause. & FIX & 2 & fixed (4201fbf) \\
R-018 & P1 & 1145--1152 & The associativity-of-concatenation proof is a plan rather than a proof: the announced induction is not carried out. & Write the induction on string length in full. & FIX & 2 & fixed (89cfaa8) \\
R-019 & P1 & 1239 & The complexity-recursion example precedes the unique-readability theorem it depends on. & Reorder so unique readability precedes the example, or add an explicit forward-dependency remark. & MIGR & 4 & open \\
R-020 & P1 & 1282 & The path characterization of the transitive closure \(R^*\) is used without proof. & State and prove the path characterization as a labeled lemma or mark its status honestly. & FIX & 2 & fixed (89cfaa8) \\
R-021 & P1 & 1819, 1823--1824 & The controlled-prose translation table and its commentary on only and unless contain translation errors. & Correct the offending table rows and the accompanying only/unless commentary. & FIX & 2 & fixed (89cfaa8) \\
R-022 & P1 & 2299--2317 & The LK soundness wording slides between valid and satisfied. & Rewrite the soundness argument keeping the two registered notions distinct throughout. & FIX & 2 & fixed (58ef447) \\
R-023 & P1 & 2547 & Davis--Putnam elimination is unattributed and stated for finite clause sets only, with no remark relating the finite restriction to compactness. & Add the primary-source attribution and an explicit finiteness/compactness scope remark. & FIX & 2 & fixed (58ef447) \\
R-024 & P1 & 3415--3418 & Conservativity of definitional extensions is asserted in an example without proof although Beth definability later leans on it. & State conservativity of definitional extensions as a labeled result with proof or citation. & FIX & 2 & fixed (58ef447) \\
R-025 & P1 & 3683--3685 & The LK universal-quantifier-right soundness case is hand-waved. & Write the case in full with the eigenvariable condition doing explicit semantic work. & FIX & 2 & fixed (58ef447) \\
R-026 & P1 & 4199--4215 & The compactness theorem statement is garbled (a dangling equivalently-clause mixes the model form with the consequence form) and the proof closes the loop redundantly. & Restate the equivalent forms unambiguously and remove or repair the redundant converse step. & FIX & 2 & fixed (58ef447) \\
R-027 & P1 & 546--548, 4127 & Broken forward promise: the front matter promises countable-language versions of the completeness theorems, but no countable or choice-free completeness theorem exists anywhere, and the AC hypothesis on completeness (4127) is thereby over-stated for the countable case. & Add a countable-language choice-free completeness theorem (or delete the promise) and scope the AC hypothesis accordingly. & FIX & 2 & fixed (58ef447) \\
R-028 & P1 & 5456--5472 & The type-space compactness proof omits the complete-type/Lindenbaum step while presented as complete. & Supply the missing step or downgrade the result's status. & FIX & 2 & fixed (b3eb76d) \\
R-029 & P1 & 5573--5585 & The saturation-implies-homogeneity proof omits its compactness appeal while presented as complete. & Supply the compactness appeal explicitly or downgrade the result's status. & FIX & 2 & fixed (b3eb76d) \\
R-030 & P1 & 6119, 13692 & The DLO \(\aleph_0\)-categoricity proof suppresses the four interval cases, which are supplied only by the exercise solution at 13692 -- the appendix doing the theorem's work. & Promote the four-case analysis from the solution into the theorem's own proof during migration (annex). & MIGR & 4 & open \\
R-031 & P1 & 10092 & The injectivity step in the second-order categoricity proof is hand-waved. & Write the injectivity argument in full. & FIX & 2 & fixed (13c1cea) \\
R-032 & P1 & 10142 & \(\mathrm{PA}_2\) is used but never defined. & Define \(\mathrm{PA}_2\) (language, axioms, semantics) before first use in the migrated part. & MIGR & 4 & open \\
R-033 & P1 & 10162--10168 & Failure of downward L\"owenheim--Skolem for full second-order semantics is justified by an undisplayed, unproved, uncited categoricity of the second-order reals. & State the second-order categoricity of the reals as a labeled result with status and citation, and cite it in the justification. & FIX & 2 & fixed (13c1cea) \\
R-034 & P1 & 8389--8394 & The ordinal trichotomy proof idea (foundation applied to the symmetric difference) is not a recognizable argument. & Replace with the standard argument via \(\alpha\cap\beta\). & FIX & 2 & fixed (13c1cea) \\
R-035 & P1 & 6340--6349 & The definition of partial computable is malformed by a double iff. & Restate the definition as a single well-formed biconditional. & FIX & 2 & fixed (563da95) \\
R-036 & P1 & 5690--5703 & Nonprincipal-ultrafilter existence is unlabeled prose whose induction is not carried out. & Promote to a labeled result and carry out the chain/Zorn argument or cite it. & FIX & 2 & fixed (b3eb76d) \\
R-037 & P1 & 11273 & Finite-matrix tableau completeness wears a proof environment over a sketch (the induction is asserted, not performed) and uses the term saturated, never defined anywhere in the book. & Define saturation for the tableau system and either perform the induction or downgrade the result to sketch status. & FIX & 2 & fixed (13c1cea) \\
R-038 & P1 & 11911 & The DPLL correctness proof is one sentence over a two-specification claim, violating the book's own template at 11791 and style rule 12354. & Supply the two-part correctness argument or downgrade to sketch status under the migration's compile-enforced status scheme. & MIGR & 4 & open \\
R-039 & P1 & 11179--11191 & Common knowledge dangles: \(C_G\) receives language and semantics but no axioms or induction rule anywhere, the non-compactness/weak-completeness caveat stated for LTL (11090) is omitted for \(C_G\), and an exercise then uses it. & Supply the \(C_G\) axiomatization (mix and induction rules) with the standard caveat as new-tree content, or rescope the exercise. & CARRY & 4 & open \\
R-040 & P1 & 10806 & The G\"odel--Gentzen translation recursion is not literally well-defined: whether \(\bot\) counts as atomic (priority of the \(\bot\)-clause versus the double-negation clause for atoms) is unstated. & Fix the clause priority explicitly in the definition. & FIX & 2 & fixed (13c1cea) \\
R-041 & P1 & 10735 & The prime-theory description is ambiguous between \(\vdash A\lor B\) and \(T\vdash A\lor B\). & Disambiguate the defining condition. & FIX & 2 & fixed (13c1cea) \\
R-042 & P1 & 13049 vs 9513 & The rule sheet declares the equality rules part of the complete active first-order sequent calculus while the cut-elimination sketch never treats equality rules, so two passages refer to different calculi under one name; the sheet's symmetry-derivability claim is correct but silently requires cut. & Name the calculi apart under the canonical-LK ruling (LK versus LK-with-equality), scope the cut-elimination claim to the calculus it covers, and state the cut-dependence of the symmetry derivation. & CARRY & 4 & open \\
S-001 & P1 & 313--319 vs 5007 & The book's three-term status policy is scoped to Part IV onward by the convention at 5007, leaving 56 of 66 numbered results in Parts I--III silently presented as complete proofs with no status labels. & Extend the status convention book-wide and label every Part I--III result, via the new tree's compile-enforced status keys. & MIGR & 4 & open \\
S-002 & P1 & 550, 3028, 4317, 4386 & Notation collision against the front-matter policy (239--241): \(|\cdot|\) denotes both cardinality (550) and structure domain (3028), yielding nested-bars expressions at 4317 and 4386. & Ratify and enforce the ruling: \(|\cdot|\) is cardinality only and domains are written \(\mathrm{dom}\,M\). & CARRY & 4 & open \\
S-003 & P1 & 1468, 1770, 3256, 4371 & Notation collision: \(\equiv\) carries four roles (metalinguistic sign, propositional equivalence, first-order equivalence, elementary equivalence -- the last registered mid-theorem at 4371). & Ratify and enforce the ruling: \(\equiv\) is elementary equivalence only, formula equivalence is \(\Leftrightarrow\), and the definitional sign is \(:\equiv\). & CARRY & 4 & open \\
S-004 & P1 & 4256--4264, 4324 & Parameter notation \(\varphi(\bar x)\) lacks the side condition \(\mathrm{FV}(\varphi)\subseteq\{\bar x\}\) and conflicts with the variable-indexing convention at 4324. & State the side condition in the notation's registration and reconcile it with the variable indexing. & CARRY & 4 & open \\
S-005 & P1 & 704--724 & The Ultrafilter Lemma, an existence theorem, is housed in a definition environment (hence unlabeled and unsourced) and contains an uncited independence claim (strictly weaker than AC over ZF). & Rehouse as a labeled theorem with status and add the Halpern--L\'evy citation. & MIGR & 4 & open \\
S-006 & P1 & 4016, 4225 & Results are cited at whole-chapter level because the intended targets are unlabeled examples. & Label the target results and repoint the citations at result level. & MIGR & 4 & open \\
S-007 & P1 & 4995--10622 & Status coverage collapses after Parts IV--VI: 26 of 77 theorem/lemma/corollary environments in the segment carry no status (all of Part VII zero, Part VIII two). & Apply the compile-enforced status keys of the new tree to every result in the segment. & MIGR & 4 & open \\
S-008 & P1 & 5007--7710 vs 7710--10622 & Three inconsistent status mechanisms are in use (Status line with proofplan environment, bare paragraph Proof sketch, paragraph Cited theorem), split by page range -- the same defect the production audit records as two sketch mechanisms (proofplan 10 times before 7710, paragraph 20 times after). & Adopt exactly one compile-enforced status mechanism book-wide. & MIGR & 4 & open \\
S-009 & P1 & 8970, 9839, 10224, 10493 & The four most important imported results -- the forcing theorem, Gentzen's consistency theorem, and both Henkin completeness theorems -- are unnumbered, unlabeled paragraphs with no TOC or index presence. & Promote all four to numbered, labeled, statused theorem environments. & MIGR & 4 & open \\
S-010 & P1 & 8733, 9065 & Misclassification: results labeled sketch are by the book's own three-way definition cited theorems. & Relabel with the correct status. & MIGR & 4 & open \\
S-011 & P1 & 6137--6145 & Uncountable categoricity of \(\mathrm{ACF}_p\) is packaged as an uncited example (Steinitz). & Promote to a labeled statused result with the Steinitz citation. & MIGR & 4 & open \\
S-012 & P1 & 2248--2695, 3440--3780 vs 9135--9948 & Name clash: LK denotes two different calculi -- Parts II/III's presentation with primitive negation rules and contexted initial sequents (reproduced by the back-matter rule sheets) versus Part VIII's variant with defined negation and context-free identity -- with no reconciliation, and the cut-elimination sketch is by its own admission incomplete for the former (9258--9260). & Ratify the canonical-LK ruling: the Parts II/III presentation is canonical and the Part VIII variant is documented as a named variant with a labeled equivalence theorem. & CARRY & 4 & open \\
S-013 & P1 & 9492 & Analytic derivation appears as a chapter title but is never defined. & Define the term at first use or retitle the chapter. & MIGR & 4 & open \\
S-014 & P1 & 10622--14780 & Status policy is nearly absent in Parts X--XI and the back matter: 3 Status lines in 4,159 lines, with 48 of the book's 60 status labels sitting in lines 5000--7999. & Apply the compile-enforced status keys to every result in the segment. & MIGR & 4 & open \\
S-015 & P1 & 10926 & The modal theorem titled Frame correspondence and completeness contains no correspondence claim -- the book's own warned-against failure mode (12357). & Retitle the theorem or add the correspondence statement it advertises. & MIGR & 4 & open \\
S-016 & P1 & 10744 & The IPL completeness sketch omits its inline citation (the controlling sources) where the modal, LTL, and epistemic sketches all cite inline. & Add the inline citation. & MIGR & 4 & open \\
S-017 & P1 & 10625--11683, 12568, 14354 & Primary-source rule violated: Glivenko 1929, G\"odel 1933, and Gentzen 1936 are absent from the bibliography (their results cited only to Troelstra volumes) and the many-valued chapter has zero citations (no Kleene at K3, no Priest 1979 at LP), against the book's own rules at 12568 and 14354. & Add the primary sources, cite them at point of use, and enforce the no-zero-citation-chapter check as a standing control. & CARRY & 4 & open \\
S-018 & P1 & 11628, 12024 & Rule displays violate the book's own naming requirement (12518): the four exponential rules appear as bare unnamed fractions, and the resolution rule's soundness-critical side conditions (renaming apart, mgu) live only in prose. & Name and label every rule and display its side conditions as part of the rule, sourced from the calculus registry. & CARRY & 4 & open \\
S-019 & P1 & 12117, 12126, 12053 & Three imported theorems live in prose with no status or label: SAT NP-completeness, FOL validity computably enumerable and undecidable (with no citation at all), and refutational completeness. & Promote each to a labeled, statused, cited statement. & MIGR & 4 & open \\
S-020 & P1 & 10628, 11566 & Same-glyph collisions within the segment: \(A\) as formula metavariable and monoid carrier in the same chapter, \(v\) as world versus registered valuation, \(u,v,w\) registered as strings but used as worlds, \(M\) as matrix versus structure, and \(V\) and \(D\) carrying three and four roles. & Apply the glyph-role rulings and the symbol-registry lint of the new tree. & CARRY & 4 & open \\
S-021 & P1 & 10677 vs 14071 & The Kripke forcing relation is typeset with two different macros in the body (10677) and the Exercise Feedback chapter (14071), and the second macro is unregistered. & Ratify one registered forcing macro and use it everywhere. & CARRY & 4 & open \\
S-022 & P1 & 14324, 14326, 14343, 14354 & Release-checklist self-test fails: the book fails four items of its own editorial checklist (undefined saturated, glyph roles, inductions without objects and cases, primary sources) while passing 14339 and 14366--14367. & Make the release checklist a mechanically checked phase gate and re-run it at release. & CARRY & 6 & open \\
S-023 & P2 & 8784, 8858, 8907 & The five-condition definition template is applied only sporadically in the set-theory chapters (poset filter and genericity, names and extension, semantic forcing), leaving sibling definitions unstructured. & Apply the new tree's uniform definition/entry template to every definition during migration. & MIGR & 4 & open \\
S-024 & P2 & 1--4995 & Chapters 10 and 12 contain no exercises while every other chapter in the segment carries them (40 exercises in Parts I--III). & None: exercises are retired from print to verification assets by the format decision, so per-chapter exercise coverage is moot. & OBS & 4 & open \\
S-025 & P2 & 13072--13220 & The qualifying-examination chapter poses ten problems but supplies solutions for only five; the final five are listed without solutions. & None: qualifying problems are retired from print to verification assets by the format decision. & OBS & 4 & open \\
L-001 & P0 & 9683 & Compile breaker: amsmath parses the bracketed discharge annotation \({[A]}\) immediately after a matrix opening as the matrix's optional alignment argument, producing an illegal-preamble-token error. & Brace-protect the bracketed annotation (applied in Phase 0, commit aa34611). & FIX & 0 & fixed \\
L-002 & P2 & 76--85 & Auto-index wrapper defects: accent commands in indexed titles are escaping hazards under makeindex (the quote character in umlauted names will be mangled) and the wrapper emits its index entry after the sectioning head, risking page drift. & Retire the wrapper; the new tree indexes inline from birth with sort keys. & OBS & 4 & open \\
L-003 & P2 & preamble & Dead macros (BV, Theory, eqdef, HOL, coNP, CNF, DNF, bracket, paren, id -- the complexity-class macros suggesting planned-never-written content -- plus the later-retired proofplan environment). & Drop them from the new preamble and add a defined-but-unused macro lint. & OBS & 4 & open \\
L-004 & P2 & 2082--2104, 3537--3561 & Array rows switch inconsistently between the cr primitive and the double-backslash line terminator. & Normalize the line terminator during migration retypesetting. & MIGR & 4 & open \\
L-005 & P2 & preamble & The rule-typesetting macro produces mixed small-caps metrics and overwide multi-premise displays (disjunction elimination, existential elimination). & Retire the macro in favor of registry-generated rule typesetting. & OBS & 4 & open \\
L-006 & P1 & 2082--2104, 3537--3561 vs 12823--13071 & Deduction rules are typeset in raw arrays -- unciteable -- and duplicated by the back-matter rule sheets, creating divergence risk. & Introduce the single calculus registry that body, entry cards, and rule sheets all input, making divergence impossible. & CARRY & 4 & open \\
L-007 & P2 & 363, 11684 & Two part headings are unlabeled while the other nine are labeled. & Label every sectioning unit in the new tree. & MIGR & 4 & open \\
L-008 & P2 & preamble & The listings package is configured but unused in Parts I--III. & Let configuration follow use in the new preamble. & OBS & 4 & open \\
L-009 & P2 & global & The shared theorem counter makes citation numbers jump (Lemma 12.7 style); whether this is intended is unverified. & Decide and document the numbering scheme in the new tree's spec, deliberately. & MIGR & 4 & open \\
L-010 & P2 & 1--14780 & The book has no figures at all in 14,780 lines while tikz and three libraries are loaded: no syntax trees, derivation trees, Kripke diagrams, or hierarchy pictures. & Execute the apparatus program's figure quota (roughly 30--40 figures) per part during migration. & CARRY & 4 & open \\
L-011 & P2 & preamble & Unused packages (mathrsfs, bm, multirow) and the unused color LogicRed. & Drop them from the new preamble. & OBS & 4 & open \\
L-012 & P2 & 8678, 7069, 10049 & Non-starred set-builder delimiters around tall bodies yield undersized braces. & Use auto-sized (starred) delimiters during migration retypesetting. & MIGR & 4 & open \\
L-013 & P2 & 7069 & Prose is set inside set-builder notation as math. & Move the prose out of math or set it with text commands. & MIGR & 4 & open \\
L-014 & P2 & 8970, 9839, 10224, 10493 and 8 uses in X--XI & paragraph sectioning commands do semantic work (sketch markers, cited theorems, internal proof boundaries) at least eight times: results at paragraph level break numbering, TOC, and referenceability and use the wrong spacing model. & Replace with numbered environments (statused results, a numbered proofboundary environment) in the new tree. & MIGR & 4 & open \\
L-015 & P2 & 76--85 & The index wrapper passes math and accent commands into the idx file with no sort keys. & Retire with the wrapper; inline index entries carry explicit sort keys. & OBS & 4 & open \\
L-016 & P2 & 10682+ & Negated forcing is typeset as a not-slash composition instead of the dedicated negated symbol, and no negated-forcing macro exists, against the book's own semantic-macro doctrine. & Register a negated-forcing macro and use it at every occurrence. & MIGR & 4 & open \\
L-017 & P2 & 12465 & A label inside an lstlisting environment is not a real label and cannot be referenced. & Move the label outside the listing or use the listings label key. & MIGR & 4 & open \\
L-018 & P2 & 12678--12774 and other longtables & Longtables lack endfirsthead, endfoot, and captions, so the 34-row register breaks across pages with no continuation marker. & Regenerate register tables from the CSV registers with proper longtable head and foot in the new apparatus. & OBS & 4 & open \\
L-019 & P2 & 12678--12821 & Register table column widths are near-overfull. & Size regenerated tables by template with compile warnings treated as errors. & OBS & 4 & open \\
L-020 & P2 & 13056--13058 & Typo in the modal rule sheet: the calculus is said to contain substitution instances of classical classically valid propositional formulas (doubled word). & None in print: rule sheets are regenerated from the calculus registry, which retires the typo with the sheet. & OBS & 4 & open \\
S-026 & P1 & 13749--13759 & Exercise solutions use marked quantifiers with \(>\) (\(\exists x>b\), \(\forall x>b\)) registered nowhere in the book; same family as R-003's \(\le\) defect (found during adversarial verification of the G2 fix designs) & Register \(>\)-marked bounded quantifiers alongside \(<\) and \(\le\), or rewrite the displays using registered comparisons & MIGR & 4 & open \\
\end{longtable}
\endgroup


\subsection{Fix designs for the Phase-2 repairs}
\label{sec:fixdesigns}

Each design below is a work order: it states the defect, the exact
mathematical repair (statements, proof routes, landing sites), and
acceptance criteria.  The designs were drafted against the v1.0 text and
adversarially verified.

% Within the fragments every \texttt argument is a plain label/file token,
% so the breakable teletype is safe there (it is NOT safe document-wide:
% \texttt elsewhere carries \- discretionaries).
\begingroup\let\texttt\spectt
% Phase-2 fix designs, group G1 (set theory): R-001, R-017.
% Fragment: compiles inside the gap-analysis document under amsmath/amssymb.

\subsection*{R-001 --- Honest CH independence (v1.0 lines 9062--9133; chapter 8665--9133)}

\paragraph{Defect.}
The theorem at 9065 claims ``the constructions using \(L\) and forcing establish'' \(\mathsf{Con}(\mathsf{ZFC})\Rightarrow\mathsf{Con}(\mathsf{ZFC}+\neg\mathsf{CH})\), but the only forcing presented is over a countable transitive model, which Warning 9107--9117 says ``assumes more than the bare arithmetical statement \(\mathsf{Con}(\mathsf{ZFC})\)'' and must be discharged; it never is.  The sketch (9092--9094) derives nonderivability by soundness from ``the two structures satisfying the displayed theories,'' but neither is a set-sized structure: \(L\) is a proper class and \(M\) is hypothesized (so too the completeness appeal at 8759--8760).

\paragraph{Repair.}
\textbf{(1) 8733--8762.}  Keep ``\(L\) is transitive and contains every ordinal'' (\emph{Status: proof sketch}).  Delete the \(\mathsf{Con}\) displays (8739--8750) and the completeness sentence (8758--8760).  Replace the satisfaction display by a separately stated scheme (\emph{Status: cited theorem}; Jech 2003, Ch.~13, pp.~175--200 --- reflection, definable well-order, condensation, as \(\mathsf{ZF}\)-theorems about the definable \(L\)): \(\mathsf{ZF}\vdash\sigma^{L}\) for every axiom \(\sigma\) of \(\mathsf{ZFC}+\mathsf{GCH}+(V{=}L)\), with \((\cdot)^{L}\) the convention of 8710--8731.
\textbf{(2) 8764--8782.}  Retitle the section ``Forcing over a countable transitive set: a pedagogical device.''  After the convention insert a Remark (\emph{Status: proof sketch}; omitted verifications: the translation convention, the absoluteness computation, and the interpreted derivability-condition bookkeeping).  \(\mathsf{Con}(\mathsf{ZFC})\) (as at 7808) is a sentence of arithmetic, and the manuscript fixes no translation of arithmetic into the membership language, so the Remark must first state one: numbers as finite ordinals under the canonical interpretation, writing \(\sigma^{\in}\) for the membership-language translation of an arithmetic sentence \(\sigma\).  Route: if such an \(M\) exists, soundness for the set-sized \((M,\in)\) gives metatheoretic consistency of \(\mathsf{ZFC}\); correctness of the fixed proof predicate (Warning at 7818--7824) makes \(\mathsf{Con}(\mathsf{ZFC})\) true in the standard structure; absoluteness of arithmetic over transitive models containing \(\omega\) (Jech, Ch.~12, pp.~155--174) then gives \((M,\in)\models\mathsf{ZFC}+\mathsf{Con}(\mathsf{ZFC})^{\in}\), and soundness again gives consistency of \(\mathsf{ZFC}+\mathsf{Con}(\mathsf{ZFC})^{\in}\).  That consistency sentence is unprovable in \(\mathsf{ZFC}+\mathsf{Con}(\mathsf{ZFC})^{\in}\) itself --- not by \texttt{thm:godel-second} verbatim, which covers computably enumerable \(T\supseteq\mathsf Q\) in the language of arithmetic and so does not apply to a membership-language theory, but by that theorem transported across the canonical interpretation, under which the fixed provability predicate still satisfies the derivability conditions (this bookkeeping is the third omitted verification); so the hypothesis strictly exceeds \(\mathsf{Con}(\mathsf{ZFC})\); other section results stay unchanged.
\textbf{(3) 9062--9096.}  Delete the theorem and sketch (9065--9096); insert:
(a) \emph{Lemma (Relativization of derivations)} [\emph{complete proof}]: for membership-language formulas \(\varphi_1,\dots,\varphi_k,\psi\) with free variables among \(v_1,\dots,v_m\), if \(\{\varphi_1,\dots,\varphi_k\}\vdash\psi\) in the fixed natural-deduction calculus, then \(\mathsf{ZF}\vdash\forall v_1\cdots\forall v_m\,(\bigwedge_{j\le m} v_j\in L\wedge\bigwedge_{i\le k}\varphi_i^{L}\to\psi^{L})\).
Route (induction on the derivation): \((\cdot)^{L}\) commutes with connectives and \(\bot^{L}=\bot\), so propositional and classical rules lift verbatim; terms are variables, so \(\forall\)-elimination and \(\exists\)-introduction consume a guard \(u\in L\) from the antecedent, while \(\forall\)-introduction and \(\exists\)-elimination form the guarded clauses of 8718--8723 (eigenvariable condition unchanged; \(x\in L\) is parameter-free); a free variable vanishing from the conclusion is discharged via \(\mathsf{ZF}\vdash\exists u\,(u\in L)\), witness \(\varnothing\in L_1\) (justify inside \(\mathsf{ZF}\) by the comprehension reading of \(\operatorname{Def}\): \(\{x\in L_0:\ldots\}=\varnothing\) by separation whatever the satisfaction clause returns on the empty domain excluded at v1.0 3028; do not route through the proof-sketch result that \(L\) contains every ordinal); equality rules act on unchanged atoms; the explicit \(\alpha\) rule (v1.0 3475--3487) maps to an \(\alpha\) step, since relativization inserts the same parameter-free guard under alpha-variant binders.  For \(m=0\) or \(k=0\) the corresponding conjunction is absent from the displayed conclusion.
(b) \emph{Theorem (Relative consistency of \(\mathsf{CH}\))} [\emph{complete proof}]: if \(\mathsf{ZF}\) is consistent, so is \(\mathsf{ZFC}+\mathsf{GCH}+(V{=}L)\); hence if \(\mathsf{ZFC}\) is consistent, so is \(\mathsf{ZFC}+\mathsf{CH}\).  Proof: a derivation of \(\bot\) from axioms \(\sigma_1,\dots,\sigma_k\), the Lemma (\(m=0\), \(\psi=\bot\)), and the item-1 scheme give \(\mathsf{ZF}\vdash\bot\); \(\mathsf{ZF}\)-derivations are \(\mathsf{ZFC}\)-derivations, and \(\mathsf{ZFC}+\mathsf{GCH}+(V{=}L)\) derives \(\mathsf{CH}\) (the \(\alpha=0\) instance of 8633), so inconsistency of \(\mathsf{ZFC}+\mathsf{CH}\) lifts to it.
(c) \emph{Cited theorem (Relative consistency of \(\neg\mathsf{CH}\))} [\emph{cited theorem}; Jech, Ch.~14, pp.~201--224]: if \(\mathsf{ZFC}\) is consistent, so is \(\mathsf{ZFC}+\neg\mathsf{CH}\).  Imported as \(\mathsf{ZFC}\)-theorem schemes about the definable class \(V^{B}\), \(B\) the completion of the poset of finite partial functions \(\omega_2\times\omega\to 2\): (i) \(\lVert\sigma\rVert_B=1\) for each \(\mathsf{ZFC}\) axiom; (ii) Boolean soundness: \(\{\sigma_1,\dots,\sigma_k\}\vdash\psi\) implies \(\mathsf{ZFC}\vdash\lVert\sigma_1\rVert\wedge\cdots\wedge\lVert\sigma_k\rVert\le\lVert\psi\rVert\), with \(\lVert\bot\rVert=0\) and \(B\) provably nontrivial; (iii) \(B\) is ccc and ccc forcing preserves cardinals: \(\lVert\check\aleph_1=\aleph_1\wedge\check\aleph_2=\aleph_2\rVert=1\); (iv) the distinctness density argument (9048--9055, per \(\omega_2\)-section) gives \(\lVert 2^{\aleph_0}\ge\check\aleph_2\rVert=1\), hence \(\lVert\neg\mathsf{CH}\rVert_B=1\).  A \(\bot\)-derivation from \(\mathsf{ZFC}\cup\{\neg\mathsf{CH}\}\) then yields \(\mathsf{ZFC}\vdash 1=0\) in \(B\): contradiction.  This is the ``Boolean-valued'' discharge Warning 9107 demands --- no countable transitive model; consistency is calculus-invariant via soundness and strong completeness.
(d) \emph{Remark (Arithmetized form)} [\emph{proof sketch}]: the two derivation transformations are primitive recursive, so \(\mathsf{Con}(\mathsf{ZFC})\to\mathsf{Con}(\mathsf{ZFC}+\mathsf{CH})\) and \(\mathsf{Con}(\mathsf{ZFC})\to\mathsf{Con}(\mathsf{ZFC}+\neg\mathsf{CH})\) are \(\mathsf{PA}\)-theorems; omitted: the arithmetization, including a primitive recursive translation between the fixed natural-deduction calculus and the source calculus of (c) --- the semantic calculus-invariance argument in (c) does not supply this; the deduction-theorem translation does.
(e) \emph{Theorem (Nonderivability)} [\emph{complete proof}]: if \(\mathsf{ZFC}\) is consistent then \(\mathsf{ZFC}\nvdash\mathsf{CH}\) and \(\mathsf{ZFC}\nvdash\neg\mathsf{CH}\).  Proof: \(\mathsf{ZFC}\vdash\mathsf{CH}\) with monotonicity and the assumption rule gives \(\mathsf{ZFC}+\neg\mathsf{CH}\vdash\mathsf{CH},\neg\mathsf{CH}\), negation elimination gives \(\vdash\bot\), and the contrapositive of (c) refutes consistency of \(\mathsf{ZFC}\); symmetrically \(\mathsf{ZFC}\vdash\neg\mathsf{CH}\) makes \(\mathsf{ZFC}+\mathsf{CH}\vdash\bot\), contradicting (b).
\textbf{(4) 9098--9117.}  Retain both Warnings, appending to the second that (c) supplies the Boolean-valued discharge; Exercise 9126 and the sketch at 13941--13952 remain correct.

\paragraph{Acceptance.}
(1) The section replacing 9065--9096 contains exactly the five results of item (3) with the stated statuses; the phrases ``the constructions using \(L\) and forcing establish'' and ``the two structures satisfying the displayed theories'' and the completeness sentence at 8759--8760 no longer occur in 8665--9133.
(2) No proof in the exact-status section uses the countable-transitive-model convention, the generic-filter lemma, or \(L\models\); nonderivability uses only monotonicity, assumption, negation elimination, and contrapositives of (b), (c).
(3) Jech citations carry pages 175--200, 201--224, 155--174; the retitled document builds cleanly.

\subsection*{R-017 --- Split the five-clause finite-coding theorem by status (v1.0 lines 678--702)}

\paragraph{Defect.}
The five-clause theorem \texttt{thm:infinite-cardinal-finite-coding} carries the single label ``Status: cited theorem,'' but the source is invoked only ``for complete proofs of the cardinal products and union bounds'' (clauses 1--4); clause 5 rests on the unlabeled inline sentence ``The final item follows by assigning each generated expression its finite syntax string and applying the preceding items'' (699--700).  A result the manuscript itself proves is hidden inside a cited status, violating the policy at 313--319 that each result carries exactly one of the three status terms and that in-document dependencies be stated separately.

\paragraph{Repair.}
\begin{enumerate}
\item Restate the theorem with clauses 1--4 only, statements unchanged, label kept, ``\emph{Status: cited theorem}''; cite Jech 2003, Theorem~3.5 (\(\aleph_\alpha\cdot\aleph_\alpha=\aleph_\alpha\); with the already-cited well-ordering theorem every infinite cardinal is an aleph) and Ch.~5 (cardinal sums and the union bound \(\lvert\bigcup_{i\in I}A_i\rvert\le\lvert I\rvert\cdot\sup_{i\in I}\lvert A_i\rvert\) under \(\mathsf{AC}\)).  The postamble 698--702 shrinks to this citation sentence plus the retained no-infinite-products caveat; the inline clause-5 sentence is deleted.
\item Immediately after (landing at v1.0 line 697) add \texttt{cor:syntax-counting}, ``\emph{Status: complete proof}'': \emph{Assume the Axiom of Choice and let \(\kappa\) be an infinite cardinal.  (i) Every set \(A\) with \(\lvert A\rvert\le\kappa\) satisfies \(\lvert A^{<\omega}\rvert\le\kappa\), where \(A^{<\omega}\) is the set of finite sequences from \(A\).  (ii) If \(\Sigma\) is an alphabet with \(\lvert\Sigma\rvert\le\kappa\), then the finite strings over \(\Sigma\), the finite tuples of such strings, and the formulas generated by any fixed finitary grammar whose alphabet extends \(\Sigma\) by at most countably many further symbols each form a set of cardinality at most \(\kappa\).}
Proof (promoted from the inline argument): (i) fix an injection \(f\colon A\to\kappa\) (the meaning of \(\lvert A\rvert\le\kappa\)); then \((a_1,\dots,a_n)\mapsto(f(a_1),\dots,f(a_n))\) is a length-preserving injection of \(A^{<\omega}\) into \(\kappa^{<\omega}\), and clause 2 gives \(\lvert\kappa^{<\omega}\rvert=\kappa\).  (ii) Strings over \(\Sigma\) lie in \(\Sigma^{<\omega}\): apply (i) with \(A=\Sigma\).  Tuples of strings lie in \((\Sigma^{<\omega})^{<\omega}\): apply (i) with \(A=\Sigma^{<\omega}\), of cardinality \(\le\kappa\) by the previous step.  Every generated expression is a finite string over the grammar's full alphabet (the manuscript's own fact, v1.0 line 1126); writing \(X\) for the at most countable set of adjoined symbols hypothesized in the statement (logical symbols, punctuation, and variables), clause 4 (applied to \(\Sigma,X,\varnothing,\varnothing,\dots\)) gives \(\lvert\Sigma\cup X\rvert\le\kappa\), so the formula set injects into \((\Sigma\cup X)^{<\omega}\) and (i) applies.
\item Retarget the clause-5 uses: v1.0 lines 3949 (Henkin--Lindenbaum, ``at most \(\kappa\) formulas'') and 4347 (downward L\"owenheim--Skolem, ``at most \(\mu\) function symbols, formulas, and finite tuples'') now cite \texttt{cor:syntax-counting}; line 3969 (clause 3) keeps its explicit reference to \texttt{thm:infinite-cardinal-finite-coding}; at 4350, the anaphor ``The same theorem gives \(\lvert N\rvert\le\mu\)'' (clause 4) currently resolves through the citation at 4347 and would silently retarget with it, so it must be rewritten as an explicit cross-reference to \texttt{thm:infinite-cardinal-finite-coding}.
\end{enumerate}

\paragraph{Acceptance.}
(1) The region replacing 678--702 contains exactly two labeled results: the four-clause theorem (``Status: cited theorem,'' Jech citation naming the product and union-bound facts) and \texttt{cor:syntax-counting} (``Status: complete proof'' with a proof environment containing the injection argument); the sentence beginning ``The final item follows'' no longer occurs.
(2) The corollary's proof appeals only to clauses 2 and 4, the fixed injection, and the formulas-are-strings fact, with no citation to Jech inside the proof.
(3) References formerly at 3949 and 4347 resolve to \texttt{cor:syntax-counting}, those at 3969 and 4350 to \texttt{thm:infinite-cardinal-finite-coding}; the document compiles with no undefined references.

% Phase-2 fix designs, group G2 (computability / arithmetization), defects R-002..R-006.
% Fragment: plain LaTeX, compilable under amsmath/amssymb when wrapped in a document.
% Line anchors refer to the v1.0 manuscript source/formal-logic-reference.tex (14780 lines).

\subsection*{R-002 --- Bounded \(\mathsf{Comp}\) versus the Cantor-pair numbering (v1.0 lines 7004--7031, 6564--6593)}

\paragraph{Defect.} The theorem at 7015--7028 asserts a bounded \(\mathsf{Comp}(e,x,s)\) recognizing accepting computations in the fixed machine numbering (6567--6585), while 7011--7013 concedes that numbering's iterated Cantor pairing ``comes with no bounded extraction formula and is not the coding used here.''  \(\mathsf{Comp}\) must still decode \(e\) and \(x\) from iterated Cantor codes; the citations (7030--7031) arithmetize other numberings.

\paragraph{Repair.} Option (ii): keep the numbering; make the history coding carry \(\beta\)-extractable \emph{decoding certificates}, so \(\mathsf{Comp}\) becomes bounded.  Replace 7004--7013 by the material below and keep 7015--7027 plus one new clause; 7030--7031 stays.

Lemma A (pairing graph; lands before 7015).  \(\mathsf{P}(x,y,z):\equiv(x+y)\times S(x+y)+(y+y)=z+z\) is quantifier-free, defines the graph of \(\operatorname{pair}_{\mathbb N}\), and \(a,b\le\operatorname{pair}_{\mathbb N}(a,b)\) (\(b\le z\); \(a\le a+b\le(a+b)(a+b+1)/2\le z\)).  Status: complete proof.  Lemma B (G\"odel \(\beta\)).  With the bounded formula
\[
  \mathsf{B}(w,i,v):\equiv\exists t<Sw\,\exists m<Sw\,\exists q<St\,
  \bigl(\mathsf{P}(t,m,w)\land t=d\times q+v\land\exists u<d\,(v+Su=d)\bigr),
\]
writing \(d\) for the term \(S((Si)\times m)\) (\(u\) fresh, \(u\notin\mathrm{Var}(d)\)).  The final conjunct is the \emph{marked-bounded} form of \(v<d\), \(\mathcal N\)-equivalent to the plain abbreviation of 6924--6926 because any witness of \(v+Su=d\) satisfies \(Su\le d\), hence \(u<d\); the marked form is required because the book counts only marked bounded quantifiers as bounded (6936--6937, 6939--6951), while a bare comparison \(r<t\) expands to an unbounded existential.  Claims: (a) for all \(w,i\) exactly one \(v\le w\) satisfies \(\mathsf{B}\) in \(\mathcal N\); write \(\beta(w,i)=v\); (b) every finite sequence \(a_0,\dots,a_{k-1}\) is \(\beta(w,\cdot)\upharpoonright k\) for some \(w\).  Proof of (a): \(\operatorname{pair}_{\mathbb N}\) is a bijection, so \(w\) determines \(t,m\) uniquely, with \(t,m\le w\) by Lemma A, matching the displayed bounds; Euclidean division of \(t\) by \(d\ge1\) determines \(q,v\) uniquely, with \(q\le t<St\) and \(v\le t\le w\).  Proof of (b): put \(m:=N!\) with \(N:=\max(k,a_0,\dots,a_{k-1})+1\); the moduli \(1+(i+1)m\) are pairwise coprime (a common prime divides \((j-i)m\) but not \(m\), hence divides \(j-i\); yet \(0<|j-i|<k\le N\) forces it to divide \(m\)); each \(a_i<m\); Chinese remaindering gives \(t\equiv a_i\pmod{1+(i+1)m}\); put \(w:=\operatorname{pair}_{\mathbb N}(t,m)\).  G\"odel 1931 and Boolos--Burgess--Jeffrey 2007, ch.~16, are parallel reading.  Status: complete proof.

Definition (history coding; replaces 7004--7013).  \(s\) is a \emph{history code} for \((e,x)\) iff, under \(\beta(s,\cdot)\) with a fixed index scheme (families at \(\operatorname{pair}_{\mathbb N}(t,j)\)), it contains: (h1) the binary digits of \(x\) with certified expansion \(v_0=d_0\), \(v_{j+1}=2v_j+d_{j+1}\), \(v_{L-1}=x\), \(v_j\le x\), \(d_j\le1\); (h2) the table of \(e\): a chain \(\mathsf P(k,c_0,e)\), \(\mathsf P(r_i,c_{i+1},c_i)\) (\(i<k-1\)), \(r_{k-1}=c_{k-1}\), rows unpacked by fixed-depth \(\mathsf P\)-steps, plus the validity checks of 6567--6576 under the designated-slot normalization of the R-005 repair; (h3) configurations \(C_0,\dots,C_T\), tapes coded in the tape-alphabet base with certified expansions, the first from (h1), each step matching the decoded table by bounded digit comparisons; (h4) final state accepting, tape prefix before the first blank a valid binary numeral (6327--6338).  Every auxiliary value is an entry of \(s\) (Lemma B(a)) or a Cantor-code component (Lemma A), so all quantifiers are bounded by terms in \(e,x,s\); and every comparison conjunct \(r<t\) or \(r\leq t\) occurring in \(\mathsf{Comp}\) is written in its marked-bounded form \(\exists u<t\,(r+Su=t)\), respectively \(\exists u\leq t\,(u+r=t)\) with \(\leq\) as registered by the R-003 repair (\(u\) fresh), each \(\mathcal N\)-equivalent to the plain abbreviation, so the matrix is bounded under the counting rule of 6936.  Repaired theorem: \(\mathsf{Comp}(e,x,s)\) is bounded; \(\mathcal N\models\mathsf{Comp}(\overline e,\overline x,\overline s)\) iff \(s\) is a history code for \((e,x)\); and \(\varphi_e(x)\) is defined iff \(\mathcal N\models\exists s\,\mathsf{Comp}(\overline e,\overline x,s)\) (existence via Lemma B(b); invalid \(e\) fail (h2), matching divergence).  Status: theorem proof sketch with plan (h1)--(h4); Lemmas A, B complete proof.

Finally, add at 7449--7454 one sentence noting that the coded-computation-sequence apparatus behind \cref{thm:representability}'s \(\Sigma^0_1\) clause (the construction paragraph at 7351--7357) is the \(\beta\) coding of Lemma B, so a single in-document coding serves machine tables, histories, and proof sequences; this discharges the register's one-coding acceptance for \(\mathsf{Prov}_T\).

\paragraph{Acceptance.} (1) 7011--7013 deleted; history coding in-document.  (2) Every displayed quantifier bounded by a term in \(e,x,s\); every comparison conjunct in marked-bounded form; monotonicity proved.  (3) 7040--7051, 7068--7074, 13735--13760 type-check unchanged.  (4) The \(\Sigma^0_1\) presentations of \(\mathsf{Comp}\) and of \(\mathsf{Prov}_T\) (7440--7454) both point at Lemma B's apparatus.

\subsection*{R-003 --- The two \(\mathsf{Q}\)-lemmas for Rosser's bounded comparison (v1.0 lines 7663--7758, 6924--6937, 7210--7222)}

\paragraph{Defect.} The Rosser proof plan (7682--7700) rests on ``evaluating the bounded quantifier in \(\mathsf Q\)'' and on the exhibited proof witnessing ``the bounded existential for every larger candidate.''  Neither is backed: \(\leq\) is never introduced (only \(<\), 6924--6926), and the two needed \(\mathsf Q\)-facts --- bounded-case enumeration and dichotomy --- are neither stated nor cited.

\paragraph{Repair.} (1) Register, next to \(<\) at 6924--6937: \(r\leq t:\equiv\exists z\,(z+r=t)\), \(z\) least fresh; \(\exists x\!\leq\! t\,\varphi:\equiv\exists x(x\leq t\land\varphi)\), \(\forall x\!\leq\! t\,\varphi:\equiv\forall x(x\leq t\to\varphi)\) for \(x\notin\mathrm{Var}(t)\), counted as bounded, like \(<\).  Record: \(\mathcal N\models x\leq t\leftrightarrow x<St\), \emph{not} claimed in \(\mathsf Q\); the witness sits \emph{left} of \(+\) so the lemmas need only \(\mathsf Q\)'s right-recursion axioms.

(2) Two numbered lemmas after \(\mathsf Q\)'s definition (7222), proved by metainduction on \(n\); each instance is a \(\mathsf Q\)-derivation from \(Sx\neq0\), \(S\)-injectivity, \(x\neq0\to\exists y\,x=Sy\), \(x+0=x\), \(x+Sy=S(x+y)\).

Lemma Q1 (bounded-case enumeration) and Lemma Q2 (dichotomy).  For every \(n\in\mathbb N\):
\[
  \mathsf Q\vdash\forall y\,\bigl(y\leq\overline n\to(y=\overline0\lor y=\overline1\lor\dots\lor y=\overline n)\bigr),
  \qquad
  \mathsf Q\vdash\forall y\,\bigl(y\leq\overline n\lor\overline n\leq y\bigr).
\]
Proof of Q1: \(n=0\): from \(z+y=0\), if \(y=Sw\) then \(S(z+w)=z+Sw=0\) contradicts \(Sx\neq0\); so \(y=0\).  Step: from \(z+y=S\overline n\): \(y=0\) gives the first disjunct; \(y=Sw\) gives \(S(z+w)=S\overline n\), \(S\)-injectivity gives \(z+w=\overline n\), so \(w\leq\overline n\); the induction hypothesis gives \(w=\overline0\lor\dots\lor w=\overline n\); apply \(S\).  Proof of Q2: \(n=0\): \(\overline0\leq y\), witness \(z:=y\), by \(y+0=y\).  Step: fix \(y\); if \(y=0\), \(y\leq S\overline n\) (witness \(S\overline n\)); if \(y=Sw\), the hypothesis on \(w\) gives either \(z+w=\overline n\), whence \(z+Sw=S(z+w)=S\overline n\) and \(y\leq S\overline n\), or \(u+\overline n=w\), whence \(u+S\overline n=S(u+\overline n)=Sw=y\) and \(S\overline n\leq y\).  Status of both: complete proof.

(3) Rewrite 7682--7700 to invoke the lemmas by label, with fixed point \(\mathsf Q\vdash R_T\leftrightarrow\forall y\,(\mathsf{Proof}_T(y,\overline{\ulcorner R_T\urcorner})\to\exists z\!\leq\! y\,\mathsf{Proof}_T(z,\overline{\ulcorner\neg R_T\urcorner}))\).  Branch \(T\vdash R_T\), least proof code \(p\): representability gives \(T\vdash\mathsf{Proof}_T(\overline p,\overline{\ulcorner R_T\urcorner})\), hence \(T\vdash\exists z\!\leq\!\overline p\,\mathsf{Proof}_T(z,\overline{\ulcorner\neg R_T\urcorner})\); consistency gives \(T\nvdash\neg R_T\), so numeralwise \(T\vdash\neg\mathsf{Proof}_T(\overline q,\overline{\ulcorner\neg R_T\urcorner})\) for \(q\leq p\); Lemma Q1 turns the bounded existential into a finite disjunction: inconsistency.  Branch \(T\vdash\neg R_T\), least code \(q\): reason in \(T\) from \(\mathsf{Proof}_T(y,\overline{\ulcorner R_T\urcorner})\); Lemma Q2 gives \(y\leq\overline q\lor\overline q\leq y\); first case: Lemma Q1 plus \(T\vdash\neg\mathsf{Proof}_T(\overline k,\overline{\ulcorner R_T\urcorner})\) for \(k\leq q\) (from \(T\nvdash R_T\)) yields a contradiction; second case: witness \(z:=\overline q\).  Hence \(T\vdash R_T\), contradicting consistency.  Status: Rosser stays proof sketch, now citing Q1/Q2 by label; the solution at 13830--13838 names both lemmas.

\paragraph{Acceptance.} (1) \(\leq\) is registered with its expansion, the bounded-counting rule, and the non-claim of \(\mathsf Q\)-equivalence with \(x<St\); no arithmetic object-formula uses \(\leq\) earlier.  (2) Q1 and Q2 appear exactly as displayed, proved by metainduction from the \(\mathsf Q\)-axioms, labeled complete proof.  (3) Each Rosser branch names which lemma discharges which step.

\subsection*{R-004 --- Jump definition and the partiality convention (v1.0 lines 7134--7189, 6576--6582)}

\paragraph{Defect.} The jump (7154--7161) uses ``the \(e\)-th oracle machine with oracle \(B\) \emph{halts} on input \(e\),'' while \(K\) (6678--6681) uses ``\(\varphi_e(e)\) is defined,'' which demands an accepting halt with valid numerical output (6576--6582); halting also covers rejection and invalid output.  So the claim at 7163--7165 that the \(K\) diagonal ``relativizes'' fails literally and is unlabeled prose; 7142--7144 also never fixes the oracle I/O convention.

\paragraph{Repair.} (1) Replace 7154--7161.  Fix an effective numbering of oracle machines analogous to 6567--6585 (same normalization, query rows included, invalid codes diverge everywhere).  For valid \(e\) and \(B\subseteq\mathbb N\), let \(\varphi^B_e(x)\) be the numerical output when the decoded machine with oracle \(B\), on \(\mathrm{bin}(x)\), halts accepting with a valid numerical output, and undefined when it rejects, diverges, or accepts without one --- 6576--6582 verbatim, relativized.  Define \(B':=\{e\in\mathbb N:\varphi^B_e(e)\text{ is defined}\}\), \(B^{(0)}:=B\), \(B^{(n+1)}:=(B^{(n)})'\).  Amend 7142--7145: oracle machines ``compute'' by this discipline; ``computably enumerable relative to \(B\)'' means: domain of some \(\varphi^B_e\).

(2) New lemma (relativized adequacy).  Every partial function computed with oracle \(B\) occurs as \(\varphi^B_e\); a relativized universal machine gives \(U^B(e,x)\simeq\varphi^B_e(x)\); and a relativized parameter theorem holds: for every \(m,n\geq1\) there is a total computable \(\mathsf{par}_n^m\) with \(\varphi^B_{\mathsf{par}_n^m(e,\bar a)}(\bar x)\simeq\varphi^B_e(\bar a,\bar x)\) for every oracle \(B\) simultaneously.  Route: R-005's adequacy proof (renaming commutes with the step relation, queries included), the universal simulation of 6608--6618 with queries passed through, and the compiler construction of 6632--6641 with query rows copied unchanged.  Status: proof sketch.

(3) Promote 7163--7165 to a numbered theorem.  Jump theorem: for every \(B\subseteq\mathbb N\), (a) \(B'\) is computably enumerable relative to \(B\); (b) \(B'\not\leq_T B\).  Proof (relativized diagonal): (a) \(B'\) is the domain of \(e\mapsto U^B(e,e)\).  (b) If an oracle machine with oracle \(B\) decides \(B'\) via total \(h^B\), define \(d^B(e):=0\) if \(h^B(e)=0\), undefined otherwise; \(d^B\) is oracle-computable (run the decider; write \(0\) and accept, or loop); fix \(b\) with \(\varphi^B_b\simeq d^B\) by the adequacy lemma; then \(b\in B'\) iff \(\varphi^B_b(b)\) is defined iff \(h^B(b)=0\) iff \(b\notin B'\).  Status: proof sketch --- the displayed diagonal calculation is complete, but clause (a) rests on the relativized universal machine and clause (b) on the occurrence clause of the relativized adequacy lemma, both proof sketches; the status policy (313--319) and the release-checklist rule (14346) forbid complete-proof status over a sketch dependency.

(4) New remark after 7189, keeping citations sound.  (i) \(K\equiv_m\varnothing'\): an ordinary table embeds as a query-free oracle table; an oracle-\(\varnothing\) machine is simulated by answering queries negatively.  (ii) The halting-based jump of the cited sources is \(\equiv_m\) this \(B'\): \(e\mapsto e^*\) simulating \(e\) and, on any halt, writing \(0\) and accepting (halting \(\leq_m\) defined); \(e\mapsto e^{**}\) diverging unless the halt is accepting with valid output (defined \(\leq_m\) halting).  All four maps are computable table transformations, uniform via the parameter theorem (\cref{thm:parameter}).  Jump monotonicity --- \(A\leq_T B\) implies \(A'\leq_m B'\) --- is stated and used explicitly: replace each \(A\)-query of an oracle-\(A\) machine by a run of a fixed oracle-\(B\) decider for \(A\), then align self-application indices by the relativized parameter theorem of the adequacy lemma with padding.  Hence Post's theorem (7167--7175) is insensitive to the convention; the citations stand.  Status: proof sketch --- the four table transformations are elementary, but the monotonicity step runs through the relativized parameter theorem, part of the sketch-status adequacy lemma.

\paragraph{Acceptance.} (1) \(B'\) is defined via ``\(\varphi^B_e(e)\) is defined,'' three-way undefinedness relativized verbatim; 7142--7145 fixes the oracle I/O convention.  (2) The jump theorem is a numbered statused result argued by the relativized diagonal invoking the labeled adequacy lemma, with a status label matching its sketch dependencies.  (3) The remark states \(K\equiv_m\varnothing'\), the \(\equiv_m\)-bridge to the halting convention, and the jump-monotonicity step it uses, each with honest status.

\subsection*{R-005 --- The enumeration claim for the fixed numbering (v1.0 lines 6564--6593, decisive at 6582--6585)}

\paragraph{Defect.} The claim that ``every unary partial computable function occurs in this numbering'' and, via the tuple convention, every \(k\)-ary one (6582--6585), is unlabeled prose with no proof, status, or citation.  Every construction that picks an index for an informally described machine (6692--6699, 6724, 6795--6816, 13735--13760) depends on it, and a bare citation cannot discharge it: the numbering is bespoke (normalization, validity constraint \(\{0,1\}\subseteq\mathcal A_{\mathrm{in}}\), wrapper).

\paragraph{Repair.} (1) Amend 6567--6569 to fix the normalization completely: a description is \emph{normalized} iff its tape symbols are the numbers \(0,\dots,|\mathcal A_{\mathrm{tape}}|-1\) with the designated assignments \(\sqcup\mapsto0\), digit \(0\mapsto1\), digit \(1\mapsto2\), and its states are the numbers \(0,\dots,|Q|-1\) with \(q_0\mapsto0\), \(q_{\mathrm{acc}}\mapsto1\), \(q_{\mathrm{rej}}\mapsto2\); the remaining numbers are assigned arbitrarily, and the row list is the lexicographic listing by numbered key \((q,a)\) over the non-halting domain.  (An ``order of first appearance'' clause would be circular: the lexicographic row order presupposes the numbering.)  No canonical machine-to-code map is needed, since indices are not unique (6587); valid codes thus have pairwise distinct designated states, and machines whose start state is halting (permitted by 6233--6234) are captured up to behavior by the surrogate-start step in the theorem's proof below.  Without the designated slots ``valid code'' is not well defined, and the validity clause \(\{0,1\}\subseteq\mathcal A_{\mathrm{in}}\) together with the input and output conventions (6296--6338) is to be read at the designated digit numbers \(1,2\).

(2) New lemma (renaming invariance; lands after 6593).  If \(\iota=(\iota_Q,\iota_A)\) bijectively renames states and tape symbols preserving the roles \(q_0,q_{\mathrm{acc}},q_{\mathrm{rej}},\sqcup,0,1\), then \(C\to_{\mathcal T}C'\) iff \(C^\iota\to_{\mathcal T^\iota}(C')^\iota\); hence \(\mathcal T\) and \(\mathcal T^\iota\) accept, reject, or diverge together on every input, with corresponding tapes.  Proof: one-step induction from 6248--6294.  Status: complete proof.

(3) Replace 6582--6585 by a numbered theorem.  Adequacy of the numbering (enumeration theorem): (a) every \(\varphi_e\) is partial computable; (b) for every \(k\geq1\) and partial computable \(f:\mathbb N^k\rightharpoonup\mathbb N\) there is \(e\) with \(\varphi_e(\bar x)\simeq f(\bar x)\) for all \(\bar x\in\mathbb N^k\), under the arity-\(k\) convention of 6588--6593.  Proof route: (a) is the wrapper argument at 6579--6582, folded in.  (b) Let \(\mathcal T\) witness \(f\) per 6340--6349.  First, \(\{0,1\}\subseteq\mathcal A_{\mathrm{in}}\) automatically: \(\operatorname{code}_k\) is surjective for every \(k\geq1\) (iterated pairing bijections), so the inputs \(\mathrm{bin}(\operatorname{code}_k(\bar x))\) exhaust all binary numerals; initial configurations exist only for words in \(\mathcal A_{\mathrm{in}}^*\) (6280--6288), and 6340--6349 demands an accepting or infinite computation on each such input; so both digits lie in \(\mathcal A_{\mathrm{in}}\), and no alphabet extension is needed.  Second, since 6233--6234 requires only \(q_{\mathrm{acc}}\neq q_{\mathrm{rej}}\), the start state may coincide with a halting state; in that case first replace it by a fresh start state \(q_0'\) with \(\delta(q_0',a)=(q_0,a,R)\) for every tape symbol \(a\) --- the one extra step leaves the tape unchanged, and the output convention ignores the final head position and off-prefix cells (6331--6338), so the accepted outputs and divergence behavior are unchanged.  Then the three designated states and three designated symbols are pairwise distinct; choose any role-preserving bijective renaming onto the designated-slot numbering, obtaining a normalized \(\mathcal T^*\) that behaves identically by the renaming lemma; list its rows lexicographically by numbered key (the row list is the full graph of \(\delta\), total on non-halting pairs by 6220--6246), encode as \(\operatorname{pair}_{\mathbb N}(k,\operatorname{code}_k(\text{rows}))\); the code \(e\) is valid by construction, and by the lemma \(e\)'s behavior on every \(\mathrm{bin}(\operatorname{code}_k(\bar x))\) is \(\mathcal T\)'s, i.e.\ \(\varphi_e(\bar x)\simeq f(\bar x)\).  Retain the Kleene 1952 citation (enumeration and normal form over standard numberings) explicitly as parallel reading; the in-document proof is primary because the numbering is bespoke.  Status afterwards: complete proof.

\paragraph{Acceptance.} (1) No unlabeled adequacy prose remains at 6582--6585; the claim exists only as the numbered theorem, clause (b) quantifying over all \(k\geq1\) with \(\varphi_e(\bar x)\simeq f(\bar x)\).  (2) The normalization convention (designated numbers for \(\sqcup,0,1,q_0,q_{\mathrm{acc}},q_{\mathrm{rej}}\)) is fixed in the numbering paragraph and used by the proof.  (3) The diagonal construction (6692--6699) and the Rice proof (6803--6816) cite the theorem by label when picking indices.

\subsection*{R-006 --- Closure-class versus prefix hierarchy (v1.0 lines 6939--7002, 7030--7189)}

\paragraph{Defect.} The hierarchy definition (6952--6971) builds \(\Sigma^0_{n+1}\), \(\Pi^0_{n+1}\) as \emph{closure classes} (under \(\land,\lor\), both bounded quantifiers, one unbounded quantifier), while every hierarchy fact discharged by citation (7030--7031, 7078--7080, 7082--7097, 7167--7189) concerns the literature's \emph{prefix} hierarchy.  Nothing identifies the two, so the citations do not apply; the missing ingredient is bounded-quantifier exchange.

\paragraph{Repair.} Insert after 7002:

Lemma C (contraction).  Let \(\mathsf P(x,y,z)\) be the quantifier-free pairing-graph formula of the R-002 repair's Lemma A, with its monotonicity bound \(a,b\le\operatorname{pair}_{\mathbb N}(a,b)\); the two designs land at adjacent sites (before 7015, after 7002), so the manuscript states \(\mathsf P\) and its lemma exactly once and Lemma C cites it.  Then, for \(z\) fresh (\(z\) not free in \(\theta\), \(z\notin\{x,y\}\)): \(\mathcal N\models\exists x\exists y\,\theta\leftrightarrow\exists z\,\exists x\!<\!Sz\,\exists y\!<\!Sz\,(\mathsf P(x,y,z)\land\theta)\), dually for \(\forall\forall\) with \(\to\).  Status: complete proof (via Lemma A).

Lemma D (bounded-quantifier exchange).  For every formula \(\theta\) and term \(t\) with \(y,w\notin\mathrm{Var}(t)\), \(w\) fresh:
\[
  \mathcal N\models\forall y\!<\!t\,\exists x\,\theta\leftrightarrow\exists w\,\forall y\!<\!t\,\exists x\!<\!w\,\theta,
  \quad
  \mathcal N\models\exists y\!<\!t\,\forall x\,\theta\leftrightarrow\forall w\,\exists y\!<\!t\,\forall x\!<\!w\,\theta .
\]
Proof: right-to-left immediate; left-to-right: finitely many \(y<\mathrm{val}(t)\), choose witnesses \(x_y\), put \(w:=\max_y x_y+1\); the second display is dual.  Equivalences \emph{in \(\mathcal N\)}, for boldface classes only; no \(\mathsf Q\)- or \(\mathsf{PA}\)-provability is claimed.  Status: complete proof.

Theorem E (prefix normal form).  For every \(n\geq0\), relations defined by prefix-form \(\Sigma^0_n\) formulas (\(n\) alternating unbounded blocks starting \(\exists\), bounded matrix) are closed, up to \(\mathcal N\)-equivalence of the defining formulas, under \(\land\), \(\lor\), bounded \(\exists\), bounded \(\forall\), unbounded \(\exists\); dually for \(\Pi^0_n\).  Hence, by induction on the clauses of 6952--6971, every closure-class formula is \(\mathcal N\)-equivalent to a prefix-form formula of the same class; since conversely every prefix-form formula is generated by those clauses, the boldface classes of 6974--7002 coincide with the cited sources'.  Proof route (execute literally): induction on \(n\), both classes at once, \(\Pi\) dual.  \(n=0\): all operations except unbounded \(\exists\) stay bounded syntactically.  \(n\geq1\), relation \(\exists x\,B\), \(B\) prefix-\(\Pi^0_{n-1}\): (i) \(\land,\lor\): rename, pull existentials out, contract by Lemma C; the residual bounded \(\exists\)'s land on prefix-\(\Pi^0_{n-1}\) formulas, absorbed at level \(n-1\).  (ii) bounded \(\exists\): \(\exists y\!<\!t\,\exists x\,B\equiv\exists x\,\exists y\!<\!t\,B\) by pure logic; absorb at level \(n-1\).  (iii) bounded \(\forall\) (collection): by Lemma D, \(\forall y\!<\!t\,\exists x\,B\) is \(\mathcal N\)-equivalent to \(\exists w\,\forall y\!<\!t\,\exists x\!<\!w\,B\); absorb both bounded quantifiers at level \(n-1\), prepend \(\exists w\).  (iv) unbounded \(\exists\): contract by Lemma C.  Base: a level-\((n-1)\) formula converts by hypothesis and is padded with dummy quantifiers.  Status afterwards: complete proof.  Cite Theorem E at 7078--7080 for the \(\mathsf{TOT}\) bound and at 7030--7031, 7097, 7178--7189 as license for the prefix-based citations.

\paragraph{Acceptance.} (1) Lemma D appears exactly as displayed, proved by the finite-witness maximum, \(\mathcal N\)-equivalence only.  (2) Theorem E's induction covers all five operations for both classes at every level, bounded-\(\forall\)-over-\(\Sigma\) via Lemma D, contractions via Lemma C.  (3) All four citation sites reference Theorem E.

% Phase-2 fix designs, group G3 (first-order chain): R-007, R-008, R-009, R-010, R-014.
% Line numbers and typewriter labels are anchors into source/formal-logic-reference.tex
% (v1.0, 14780 lines).  Plain LaTeX; compilable under amsmath/amssymb.
% Set in \small so that each design occupies at most one page.
\begingroup\small

\subsection*{R-007 --- Prenex normal form: missing lemma chain (v1.0 lines 4396--4446)}

\paragraph{Defect.}
The prenex theorem's proof (4424--4446) opens with unproved global operations:
``move negations inward using De Morgan equivalences and quantifier duality,
and alpha-rename bound variables so that the formula is rectified''
(4425--4427).  No first-order negation-normal-form lemma exists (the
propositional theorem, 1899--1917, covers no quantifiers and its termination
measure counts only connectives below a negation); no rectified-variant
existence result exists (no renaming recursion is ever defined); and
``quantifier duality'' resolves to \texttt{lem:quantifier-duality} (3814), a
derivability lemma of \(\mathsf{ND}_{\mathrm{FOL}}\), where a semantic
equivalence is required.  The extraction induction then rests on
``Rectification guarantees the side condition'' (4441) without any stated
invariant.

\paragraph{Repair.}
Insert Lemmas A--C between the definitions (4399--4418) and the theorem
(4420), rewrite the theorem's proof as their composition, and attach the
missing label \texttt{cor:alpha-invariance} at 3203; \(c\) is the complexity
of 2875 (counting quantifiers, so the measure covers quantifier prefixes) and
\(\equiv\) is logical equivalence (3256).

\emph{Lemma A (\texttt{lem:fol-nnf}, negation normal form).}  Define
\(\operatorname{nnf},\operatorname{nnf}':\mathrm{Fm}(\mathcal{L})\to\mathrm{Fm}(\mathcal{L})\)
by simultaneous recursion: \(\operatorname{nnf}\) fixes atomic formulas and
\(\bot\), commutes with \(\wedge,\vee,\forall,\exists\), and satisfies
\(\operatorname{nnf}(\neg\varphi)=\operatorname{nnf}'(\varphi)\),
\(\operatorname{nnf}(\varphi\to\psi)=\operatorname{nnf}'(\varphi)\vee\operatorname{nnf}(\psi)\);
and, for \(\alpha\) atomic,
\(\operatorname{nnf}'(\alpha)=\neg\alpha\),
\(\operatorname{nnf}'(\bot)=\neg\bot\),
\(\operatorname{nnf}'(\neg\varphi)=\operatorname{nnf}(\varphi)\),
\(\operatorname{nnf}'(\varphi\wedge\psi)=\operatorname{nnf}'(\varphi)\vee\operatorname{nnf}'(\psi)\),
\(\operatorname{nnf}'(\varphi\vee\psi)=\operatorname{nnf}'(\varphi)\wedge\operatorname{nnf}'(\psi)\),
\(\operatorname{nnf}'(\varphi\to\psi)=\operatorname{nnf}(\varphi)\wedge\operatorname{nnf}'(\psi)\),
\(\operatorname{nnf}'(\forall x\,\varphi)=\exists x\,\operatorname{nnf}'(\varphi)\),
\(\operatorname{nnf}'(\exists x\,\varphi)=\forall x\,\operatorname{nnf}'(\varphi)\)
(every call is on an immediate subformula, so \(c\) strictly decreases and the
recursion is legitimate).  Claims, by simultaneous strong induction on
\(c(\varphi)\): (i) both values are in negation normal form (definition
4399--4408; \(\neg\bot\) is admitted there); (ii)
\(\operatorname{nnf}(\varphi)\equiv\varphi\) and
\(\operatorname{nnf}'(\varphi)\equiv\neg\varphi\), each case checked against
the satisfaction clauses (3111--3145): the \(\operatorname{nnf}'\) quantifier
cases are the pointwise computations
\(\mathcal{M},s\models\neg\forall x\,\varphi\) iff some \(a\) has
\(\mathcal{M},s[x\mapsto a]\models\neg\varphi\) iff
\(\mathcal{M},s\models\exists x\,\neg\varphi\), dually for \(\exists\) ---
not citations of \texttt{lem:quantifier-duality}; (iii) \(\mathrm{FV}\) and
\(\mathrm{Var}\) of both values equal those of \(\varphi\).

\emph{Lemma B (\texttt{lem:rectified-variant}, rectified variant).}  For
\(\varphi\) and finite \(V\supseteq\mathrm{Var}(\varphi)\) define
\(r(\varphi,V)\): atomic and \(\bot\) unchanged; \(r(\neg\varphi,V)=\neg
r(\varphi,V)\); \(r(\varphi\circ\psi,V)=r(\varphi,V)\circ r(\psi,\,V\cup
\mathrm{Var}(r(\varphi,V)))\) for \(\circ\in\{\wedge,\vee,\to\}\);
\(r(Qx\,\varphi,V)=Qz\,r(\varphi[z/x],\,V\cup\{z\})\) with \(z\) least outside
\(V\).  Well-foundedness is the 2951--2958 pattern:
\(z\notin\mathrm{Var}(\varphi)\cup\{x\}\), so no binder of \(\varphi\) equals
\(z\), the substitution runs on the first two quantifier clauses only (textual
replacement, so
\(\mathrm{Var}(\varphi[z/x])\subseteq\mathrm{Var}(\varphi)\cup\{z\}\)), and
\(c(\varphi[z/x])=c(\varphi)<c(Qx\,\varphi)\) by
\texttt{lem:substitution-complexity} (2965).  Claims, by strong induction on
\(c(\varphi)\) under the invariant \(\mathrm{Var}(\varphi)\subseteq V\):
(i) \(r(\varphi,V)\equiv_{\alpha}\varphi\) --- connective cases by congruence,
quantifier case by the generating clause \(Qx\,\varphi\equiv_\alpha
Qz\,(\varphi[z/x])\) (2984, licensed by \(z\notin\mathrm{Var}(\varphi)\)), the
\(Q\)-congruence, and transitivity --- hence \(r(\varphi,V)\equiv\varphi\) by
\texttt{cor:alpha-invariance}; (ii) \(r(\varphi,V)\) is rectified: its binders
are pairwise distinct and outside \(V\) (the binary clause forbids the left
output's variables to the right recursion), while
\(\mathrm{FV}(r(\varphi,V))=\mathrm{FV}(\varphi)\subseteq V\) by (iii);
(iii) alpha-equivalent formulas have equal free-variable sets: the
equal-\(\mathrm{FV}\) relation is a congruence equivalence containing each
generating pair, since both cases of clause~(b) of
\texttt{lem:substitution-fv} (added by R-010) give
\(\mathrm{FV}(Qy\,(\varphi[y/x]))=\mathrm{FV}(Qx\,\varphi)\), so leastness of
\(\equiv_\alpha\) applies; (iv) \(r\) preserves the constructor tree, hence
negation normal form.

\emph{Lemma C (\texttt{lem:prenex-extraction}, extraction with invariant).}
For every rectified \(\varphi\) in negation normal form there is
\(p(\varphi)\) with: (i) \(p(\varphi)\equiv\varphi\); (ii) \(p(\varphi)\)
rectified and prenex with quantifier-free negation-normal matrix; (iii)
\emph{invariant:} \(\mathrm{FV}\) and \(\mathrm{Var}\) unchanged, and the
prefix consists of exactly the quantifier occurrences of \(\varphi\) (same
quantifier, same variable) --- no new free or bound variables.  Route:
structural induction on the negation-normal generation (subformulas of
rectified formulas are rectified); \(p\) is the identity on base formulas and
commutes with \(Qx\).  Binary case \(\alpha\circ\beta\),
\(\circ\in\{\wedge,\vee\}\): with
\(p(\alpha)=Q_1y_1\cdots Q_ky_k\,\alpha_0\),
\(p(\beta)=R_1z_1\cdots R_mz_m\,\beta_0\), set
\(p(\alpha\circ\beta)=Q_1y_1\cdots Q_ky_k\,R_1z_1\cdots R_mz_m\,
(\alpha_0\circ\beta_0)\), proved \(\equiv\alpha\circ\beta\) by an inner
induction on \(k+m\) from two recorded facts: (C1)
\((Qy\,\pi)\circ\chi\equiv Qy\,(\pi\circ\chi)\) and
\(\chi\circ(Qy\,\pi)\equiv Qy\,(\chi\circ\pi)\) whenever
\(y\notin\mathrm{FV}(\chi)\), proved from the satisfaction clauses and
\texttt{lem:fol-coincidence} (3148), displaying the \(\forall\)/\(\vee\) case
(the classical distribution of \(\forall\) over a disjunct constant in the
quantified variable); (C2) \(\equiv\) is a congruence for \(\circ\) and
\(Qx\) (immediate from the satisfaction clauses).  The side conditions hold
because \(y_i\notin\mathrm{FV}(p(\beta))=\mathrm{FV}(\beta)\) (\(y_i\) is
bound in the rectified \(\alpha\circ\beta\)) and
\(z_j\notin\mathrm{FV}(\alpha_0)\subseteq
\mathrm{FV}(\alpha)\cup\{y_1,\ldots,y_k\}\) (rectifiedness, binder
distinctness); the output is rectified and satisfies the invariant since
nothing is renamed or introduced.

\emph{Rewritten theorem proof} (replaces 4424--4446; the statement gains label
\texttt{thm:prenex-normal-form}): by Lemmas A--C,
\(p\bigl(r(\operatorname{nnf}(\varphi),\mathrm{Var}(\operatorname{nnf}(\varphi)))\bigr)
\equiv\varphi\) is rectified, prenex, and in the same signature; all four
results then carry \textit{Status: complete proof}.

\paragraph{Acceptance.}
(1) The rewritten proof cites only Lemmas A--C and
\texttt{cor:alpha-invariance}; ``quantifier duality'' no longer occurs in
4396--4446.  (2) Each lemma names its induction measure; Lemma B's renaming
recursion cites \texttt{lem:substitution-complexity} for its non-structural
call on \(\varphi[z/x]\).  (3) Lemma C states the no-new-variables invariant and every extraction
step's side condition is discharged from it plus rectifiedness.

\subsection*{R-008 --- Skolemization: explicit prefix induction replacing the ``quantifier game'' (v1.0 lines 4480--4512)}

\paragraph{Defect.}
Both directions of \texttt{thm:skolem-equisatisfiable} (4492--4507) are prose:
``Induction along the prefix then shows\ldots'' states no induction hypothesis,
and the converse appeals to ``the original quantifier game'' (4501) --- a notion
defined nowhere in the document.  The Axiom of Choice is hypothesized for the
whole theorem (4482) although only the expansion direction uses it.

\paragraph{Repair.}
Fix \(\sigma=Q_1x_1\cdots Q_nx_n\,\theta\) rectified prenex.  For \(0\le i\le
n\) let \(\sigma_i=Q_{i+1}x_{i+1}\cdots Q_nx_n\,\theta\) (so
\(\sigma_0=\sigma\), \(\sigma_n=\theta\), and
\(\mathrm{FV}(\sigma_i)\subseteq\{x_1,\ldots,x_i\}\)); let \(q(i)\) count the
universal positions among the first \(i\), let \(\bar u(j)\) list the universal
variables before position \(j\), and for existential \(j\) let
\(t_j=f_j(\bar u(j))\) be its Skolem term, so
\(\operatorname{Sk}(\sigma)=\forall u_1\cdots\forall u_p\,\theta^{*}\) with
\(\theta^{*}=\theta[t_j/x_j]\) simultaneously for all existential \(j\)
(\(\theta\) is quantifier-free, so this is plain replacement; since no \(t_j\)
contains an existential variable, the simultaneous substitution equals the
composition of the single substitutions in any order).  Given a structure and a
tuple \(\bar a\) of values for the universal variables among
\(x_1,\ldots,x_i\), define the \emph{Skolem-term assignment} \(s_{\bar a}\):
universal \(x_j\) gets its entry of \(\bar a\); existential \(x_j\) gets the
value \(v_j\) of \(t_j\) computed from the entries for \(\bar u(j)\); values of
other variables are irrelevant by \texttt{lem:fol-coincidence}.

\emph{Restated theorem} (replaces 4480--4491): let \(\sigma\) be a rectified
prenex \(\mathcal{L}\)-sentence.  (1) \emph{(Axiom of Choice.)}  Every
\(\mathcal{L}\)-structure satisfying \(\sigma\) has an expansion to the Skolem
signature satisfying \(\operatorname{Sk}(\sigma)\).  (2) \emph{(No choice.)}
The \(\mathcal{L}\)-reduct of every model of \(\operatorname{Sk}(\sigma)\)
satisfies \(\sigma\).  Hence, assuming the Axiom of Choice, \(\sigma\) has a
model iff \(\operatorname{Sk}(\sigma)\) has one; the right-to-left implication
holds without choice.

\emph{Proof of (1).}  Let \(\mathcal{M}\models\sigma\).  By finite recursion on
\(i\) choose set-functions \(F_j\) (for existential \(j\le i\)) while proving
\[
 \mathrm{Inv}(i):\quad
 \mathcal{M},s_{\bar a}\models\sigma_i
 \quad\text{for every }\bar a\in|\mathcal{M}|^{\,q(i)},
\]
where \(s_{\bar a}\) computes \(v_j=F_j(\text{entries of }\bar a\text{ for }
\bar u(j))\).  \(\mathrm{Inv}(0)\) is the hypothesis.  Universal step: from
\(\mathcal{M},s_{\bar a}\models\forall x_{i+1}\,\sigma_{i+1}\), the
satisfaction clause gives \(\mathcal{M},s_{\bar a}[x_{i+1}\mapsto b]\models
\sigma_{i+1}\) for every \(b\), which is \(\mathrm{Inv}(i+1)\) for
\((\bar a,b)\).  Existential step: by \(\mathrm{Inv}(i)\) each witness set
\(W(\bar a)=\{b:\mathcal{M},s_{\bar a}[x_{i+1}\mapsto b]\models\sigma_{i+1}\}\)
is nonempty; apply the Axiom of Choice (definition at 651) to
\(\{\{\bar a\}\times W(\bar a):\bar a\in|\mathcal{M}|^{\,q(i)}\}\) to obtain
\(F_{i+1}\) with \(F_{i+1}(\bar a)\in W(\bar a)\) --- one application of choice
per existential quantifier, and nowhere else.  Define the expansion
\(\mathcal{M}^{*}\) by \(f_j^{\mathcal{M}^{*}}=F_j\).  From
\(\mathrm{Inv}(n)\): for every \(\bar a\in|\mathcal{M}|^{p}\), the term-value
clauses give \(t_j\) the value \(v_j\) under \([\bar u\mapsto\bar a]\) in
\(\mathcal{M}^{*}\), so iterating \texttt{thm:fol-substitution} once per
existential position turns \(\mathcal{M},s_{\bar a}\models\theta\) (equal to
\(\mathcal{M}^{*},s_{\bar a}\models\theta\) by \texttt{lem:reduct-invariance})
into \(\mathcal{M}^{*},[\bar u\mapsto\bar a]\models\theta^{*}\); the
\(p\)-fold universal clause yields
\(\mathcal{M}^{*}\models\operatorname{Sk}(\sigma)\).

\emph{Proof of (2).}  Let \(\mathcal{N}\models\operatorname{Sk}(\sigma)\) in
the expanded signature; now \(v_j\) is the value of the term \(t_j\) in
\(\mathcal{N}\).  By downward induction from \(i=n\) to \(i=0\) prove
\(\mathrm{Inv}'(i)\): \(\mathcal{N},s_{\bar a}\models\sigma_i\) for every
\(\bar a\in|\mathcal{N}|^{\,q(i)}\).  Base \(i=n\): instantiate the universal
prefix of \(\operatorname{Sk}(\sigma)\) at \(\bar a\) and reverse the same
iterated substitution-lemma computation as in (1).  Universal step
(\(Q_{i+1}=\forall\)): for every \(b\), \(\mathrm{Inv}'(i+1)\) at
\((\bar a,b)\) gives \(\mathcal{N},s_{\bar a}[x_{i+1}\mapsto b]\models
\sigma_{i+1}\) (the \(v_j\), \(j\le i\), depend only on \(\bar a\)); apply the
\(\forall\)-clause.  Existential step: \(\mathrm{Inv}'(i+1)\) at \(\bar a\)
exhibits the witness \(v_{i+1}\), and the \(\exists\)-clause gives
\(\mathrm{Inv}'(i)\).  Then \(\mathrm{Inv}'(0)\) and
\texttt{lem:reduct-invariance} yield
\(\mathcal{N}\mathbin{\upharpoonright}\mathcal{L}\models\sigma\).  No choice
principle is used.  The remark at 4509--4512 stays.  The theorem carries
\textit{Status: complete proof} afterwards.

\paragraph{Acceptance.}
(1) The string ``quantifier game'' occurs nowhere in the manuscript.  (2) Each
direction displays an induction statement indexed by prefix position with the
Skolem-term assignment defined, and every step is discharged by a named
satisfaction clause, \texttt{thm:fol-substitution},
\texttt{lem:fol-coincidence}, or \texttt{lem:reduct-invariance}.  (3) The Axiom
of Choice appears only in clause (1), invoked exactly once per existential
quantifier on a displayed family of witness sets; clause (2)'s proof invokes no
choice principle.

\subsection*{R-009 --- Transfinite Lindenbaum lemma for the Henkin construction (v1.0 lines 3934--3985)}

\paragraph{Defect.}
The second paragraph of the Henkin--Lindenbaum proof (3978--3985) says
``Well-order all formulas of \(\mathcal{L}^{H}\) and repeat the Lindenbaum
construction'' and claims maximal consistency ``by the same finite-derivation
argument as in \texttt{lem:prop-lindenbaum}.''  That lemma (2381--2405) is an
\(\omega\)-indexed recursion over a countable enumeration; for
\(\kappa\) possibly uncountable the construction needs transfinite recursion
with limit stages, which the propositional argument does not contain, and the
family whose union is taken is never defined.  The witness stage two paragraphs
earlier (3961--3976) does all of this explicitly via
\texttt{thm:transfinite-recursion}; the two halves of one proof sit at
different rigor levels.

\paragraph{Repair.}
Insert between 3932 and 3934:

\emph{Lemma (\texttt{lem:transfinite-lindenbaum}).}  Let \(\mathcal{L}'\) be a
signature, \(\lambda\) an ordinal, and
\(\langle\varphi_\alpha:\alpha<\lambda\rangle\) an enumeration of
\(\mathrm{Fm}(\mathcal{L}')\) in which every formula occurs at least once.
Then every set \(\Sigma\subseteq\mathrm{Fm}(\mathcal{L}')\) consistent for
\(\mathsf{ND}_{\mathrm{FOL}}\) is contained in a maximal consistent
\(\Delta\subseteq\mathrm{Fm}(\mathcal{L}')\).  (Choice-free given the
enumeration; the Axiom of Choice enters only when producing one.)

\emph{Proof route.}  By \texttt{thm:transfinite-recursion} (632) with domain
\(\lambda+1\) define \(\langle\Gamma_\alpha:\alpha\le\lambda\rangle\):
\[
 \Gamma_0=\Sigma,\qquad
 \Gamma_{\alpha+1}=
 \begin{cases}
   \Gamma_\alpha\cup\{\varphi_\alpha\},&\text{if this set is consistent},\\
   \Gamma_\alpha,&\text{otherwise},
 \end{cases}
 \qquad
 \Gamma_\beta=\bigcup_{\alpha<\beta}\Gamma_\alpha\ \text{ for limit }\beta.
\]
\emph{Chain claim:} \(\alpha\le\beta\le\lambda\) implies
\(\Gamma_\alpha\subseteq\Gamma_\beta\) (transfinite induction on \(\beta\)).
\emph{Consistency:} transfinite induction on \(\alpha\); zero by hypothesis;
successor by the case split; at a limit \(\beta\), a derivation of \(\bot\)
from \(\Gamma_\beta\) is finite, so its open assumptions form a finite
\(F\subseteq\Gamma_\beta\); each \(\psi\in F\) lies in some
\(\Gamma_{\alpha_\psi}\) with \(\alpha_\psi<\beta\), and with
\(\alpha^{*}=\max_{\psi\in F}\alpha_\psi<\beta\) if \(F\ne\varnothing\) (a
finite nonempty set of ordinals has a maximum) and \(\alpha^{*}=0<\beta\) if
\(F=\varnothing\), the chain claim locates \(F\subseteq\Gamma_{\alpha^{*}}\),
contradicting the induction hypothesis at \(\alpha^{*}\).  \emph{Maximality of}
\(\Delta=\Gamma_\lambda\): if \(\varphi\notin\Delta\), pick \(\alpha\) with
\(\varphi_\alpha=\varphi\); rejection at stage \(\alpha+1\) means
\(\Gamma_\alpha\cup\{\varphi\}\) derives \(\bot\), and since a derivation from
a set is a derivation from any superset (monotonicity, immediate from the
definition of derivability from assumptions),
\(\Delta\cup\{\varphi\}\) derives \(\bot\).

\emph{Rewrite of 3978--3985:} ``The set \(\mathrm{Fm}(\mathcal{L}^{H})\) has
cardinality exactly \(\kappa\), as established above from
\texttt{thm:infinite-cardinal-finite-coding} and the formulas
\(\exists x_0\,(x_0=d_\xi)\) (the count at 3949--3952 is retained); by the
cardinal-assignment corollary (662, under the Axiom of Choice already assumed)
fix a bijective enumeration
\(\langle\varphi_\alpha:\alpha<\kappa\rangle\) of it.  Apply
\texttt{lem:transfinite-lindenbaum} to \(\Sigma^{H}\) to obtain a maximal
consistent \(\Delta\supseteq\Sigma^{H}\supseteq\Sigma\).''  The existing
witness-property sentences (3981--3984) follow unchanged.  Status labels
afterwards: the new lemma and \texttt{lem:henkin-lindenbaum} both
\textit{Status: complete proof.}

\paragraph{Acceptance.}
(1) A lemma between \texttt{lem:new-constant-elimination} and
\texttt{lem:henkin-lindenbaum} displays the recursion with zero, successor, and
limit clauses and cites \texttt{thm:transfinite-recursion}.  (2) Its limit case
contains the chain claim and locates an arbitrary finite subset of a limit
stage inside a single earlier stage via the maximum of finitely many ordinals.
(3) Lines 3978--3981 no longer contain ``repeat the Lindenbaum construction''
or the appeal to \texttt{lem:prop-lindenbaum}; the paragraph applies the new
lemma to \(\Sigma^{H}\) under an explicitly fixed enumeration of order type
\(\kappa\).

\subsection*{R-010 --- Vacuous substitution is alpha-invariant: strong induction on complexity (v1.0 lines 3006--3016)}

\paragraph{Defect.}
The unnamed lemma at 3006 (``If \(x\notin\mathrm{FV}(\varphi)\), then
\(\varphi[t/x]\equiv_\alpha\varphi\)'') is proved ``by structural induction,''
but in the renaming case the recursion visits \(\psi[z/y]\), which is not a
subformula of \(Qy\,\psi\), so plain structural induction does not license the
step; the case is discharged by the bare assertion that renaming ``changes the
formula only within its alpha-equivalence class.''  The book's own template for
this situation is the simultaneous strong induction on complexity at 2951--2958
and its use at 3193--3199 of \texttt{lem:substitution-complexity}.

\paragraph{Repair.}
First insert, after the example ending at 3004 and before the lemma at 3006:

\emph{Lemma (\texttt{lem:substitution-fv}, free variables under
substitution).}  For every term \(r\): \(x\notin\mathrm{Var}(r)\) implies
\(r[t/x]=r\), and
\(\mathrm{Var}(r[t/x])\subseteq(\mathrm{Var}(r)\setminus\{x\})\cup\mathrm{Var}(t)\)
(induction on \(r\)).  For every formula \(\varphi\), term \(t\), and variables
\(x,y\):
\begin{itemize}
\item[(a)] \(\mathrm{FV}(\varphi[t/x])\subseteq
      (\mathrm{FV}(\varphi)\setminus\{x\})\cup\operatorname{Var}(t)\)
      (the book defines no \(\mathrm{FV}\) on terms; a term's variables
      are its \(\operatorname{Var}\), as in the \(\mathrm{FV}\) clauses for
      atomic formulas at 2862--2863);
\item[(b)] if \(y\notin\mathrm{Var}(\varphi)\), then
      \(\mathrm{FV}(\varphi[y/x])=(\mathrm{FV}(\varphi)\setminus\{x\})\cup\{y\}\)
      when \(x\in\mathrm{FV}(\varphi)\), and
      \(\mathrm{FV}(\varphi[y/x])=\mathrm{FV}(\varphi)\) when
      \(x\notin\mathrm{FV}(\varphi)\).
\end{itemize}
Route: strong induction on \(c(\varphi)\).  For (b) the renaming clause of the
substitution definition never fires (a binder \(w\) of a subformula lies in
\(\mathrm{Var}(\varphi)\), so \(w\ne y\)), and the substitution is textual
replacement.  For (a)'s renaming case
\((Qy\,\psi)[t/x]=Qz\,((\psi[z/y])[t/x])\): apply the induction hypothesis to
\(\psi\) with \([z/y]\) and to \(\psi[z/y]\) with \([t/x]\) --- the latter at
equal complexity \(c(\psi[z/y])=c(\psi)<c(Qy\,\psi)\), licensed by
\texttt{lem:substitution-complexity} (2965) --- then compute, using
\(z\notin\mathrm{Var}(t)\cup\{x\}\), that
\(\mathrm{FV}(Qz\,((\psi[z/y])[t/x]))\subseteq
(\mathrm{FV}(Qy\,\psi)\setminus\{x\})\cup\operatorname{Var}(t)\).

Then replace the proof at 3010--3016 by: \emph{Proceed by strong induction on
formula complexity \(c\) (2875).}  Atomic case: \(x\) occurs in no term, so
every term is unchanged by the term clause of \texttt{lem:substitution-fv} and
\(\varphi[t/x]=\varphi\); \(\bot\) likewise.  Connective cases:
\(x\notin\mathrm{FV}\) of each immediate subformula; induction hypotheses and
the congruence property of \(\equiv_\alpha\).  Quantifier \(Qy\,\psi\) with
\(y=x\): substitution stops, reflexivity.  With \(y\ne x\),
\(y\notin\mathrm{Var}(t)\): then \(x\notin\mathrm{FV}(\psi)\); the induction
hypothesis gives \(\psi[t/x]\equiv_\alpha\psi\), and the \(Q\)-congruence
concludes.  With \(y\ne x\), \(y\in\mathrm{Var}(t)\): the definition returns
\(Qz\,((\psi[z/y])[t/x])\) with \(z\) least outside
\(\mathrm{Var}(\psi)\cup\mathrm{Var}(t)\cup\{x,y\}\).  By clause (b) of
\texttt{lem:substitution-fv}, \(x\notin\mathrm{FV}(\psi[z/y])\) (as
\(x\notin\mathrm{FV}(\psi)\) and \(x\ne z\)).  By
\texttt{lem:substitution-complexity}, \(c(\psi[z/y])=c(\psi)<c(Qy\,\psi)\), so
the \emph{strong} induction hypothesis applies to \(\psi[z/y]\) and gives
\((\psi[z/y])[t/x]\equiv_\alpha\psi[z/y]\); the \(Q\)-congruence gives
\(Qz\,((\psi[z/y])[t/x])\equiv_\alpha Qz\,(\psi[z/y])\); the generating clause
(licensed since \(z\notin\mathrm{Var}(\psi)\)) gives
\(Qz\,(\psi[z/y])\equiv_\alpha Qy\,\psi\); transitivity concludes.  Landing:
new lemma at approximately 3005; rewritten proof in place at 3010--3016.  Both
results carry \textit{Status: complete proof} afterwards.  (Clause (b) is also
consumed by R-007's Lemma B.)

\paragraph{Acceptance.}
(1) The proof at 3010--3016 opens ``by strong induction on formula
complexity'' and cites \texttt{lem:substitution-complexity} exactly where the
induction hypothesis is applied to \(\psi[z/y]\).  (2) A stated and proved
\texttt{lem:substitution-fv} precedes the lemma and is cited for
\(x\notin\mathrm{FV}(\psi[z/y])\).  (3) No case applies an induction
hypothesis to a formula that is neither of strictly smaller complexity nor of
equal complexity justified by \texttt{lem:substitution-complexity}.

\subsection*{R-014 --- \(\mathsf{LK}\) completeness for arbitrary sequents (v1.0 lines 4165--4185, 2488--2507; soundness 3674--3687, 2295--2318)}

\paragraph{Defect.}
Soundness (3674) covers every derivable sequent: the calculus is
multi-succedent (rules at 2178--2234 and 3583--3618, satisfaction convention at
3579--3581), and the theorem quantifies over sequents
\(\Gamma\Rightarrow\Delta\).  Completeness (4165) is stated only ``for every
finite multiset \(\Gamma\) and formula \(\varphi\)'' --- single-succedent.  No
completeness statement covers sequents with empty or multi-formula succedents,
so the two theorems do not compose to a derivability-iff-validity equivalence
for the objects the calculus manipulates.  The propositional pair exhibits the
identical asymmetry --- \texttt{thm:prop-lk-soundness} (2295--2318) covers
every derivable \(\mathsf{LK}_{\mathrm{CPL}}\) sequent while completeness
(2488--2491) is again single-succedent --- and register row R-014 spans both
pairs (2488 vs 2295; 4165 vs 3674), so this design must repair both.

\paragraph{Repair.}
Attach \texttt{thm:fol-lk-soundness} to the theorem at 3674.  Replace the
statement at 4165--4170 by:

\emph{Theorem (\texttt{thm:fol-lk-completeness}).}  Assume the Axiom of Choice
in the ambient set theory.  For all finite multisets \(\Gamma,\Delta\) of
\(\mathcal{L}\)-formulas, the sequent \(\Gamma\Rightarrow\Delta\) is derivable
in \(\mathsf{LK}_{\mathrm{FOL}}\) iff for every \(\mathcal{L}\)-structure
\(\mathcal{M}\) and assignment \(s\): if \(\mathcal{M},s\models\gamma\) for
every \(\gamma\in\Gamma\), then \(\mathcal{M},s\models\delta\) for some
\(\delta\in\Delta\).

(The roadmap's propositional phrasing ``every valuation'' is this schema
instantiated for \(\mathsf{LK}_{\mathrm{CPL}}\); at first order the semantic
data are structure plus assignment, matching 3579--3581.)  Left to right is
\texttt{thm:fol-lk-soundness}.  Right to left, by reduction to
\texttt{thm:fol-completeness} (4128) plus the existing translation:
(1)~let \(\Delta=\delta_1,\ldots,\delta_m\) (\(m\ge 0\)) and let \(S\) be the
set of members of \(\Gamma\) together with
\(\neg\delta_1,\ldots,\neg\delta_m\); validity of the sequent makes \(S\)
unsatisfiable, i.e.\ \(S\models\bot\), since a structure and assignment
satisfying all of \(S\) would satisfy some \(\delta_j\) and its negation.
(2)~\texttt{thm:fol-completeness} gives
\(S\vdash_{\mathsf{ND}_{\mathrm{FOL}}}\bot\); the derivation's open assumptions
form a finite \(S_0\subseteq S\), say \(\gamma_1,\ldots,\gamma_k\in\Gamma\)
and \(\neg\delta_{j_1},\ldots,\neg\delta_{j_r}\); the translation of the
existing sketch (4171--4185) turns it into an
\(\mathsf{LK}_{\mathrm{FOL}}\)-derivation of
\(\gamma_1,\ldots,\gamma_k,\neg\delta_{j_1},\ldots,\neg\delta_{j_r}
\Rightarrow\bot\).
(3)~Cut with the initial \(\bot\)-sequent \(\bot\Rightarrow\;\) (instance of
\(\Gamma,\bot\Rightarrow\Delta\) with empty contexts, 2180) empties the
succedent.
(4)~For each \(j\), the sequent \(\;\Rightarrow\delta_j,\neg\delta_j\) is
derivable (identity sequent \(\delta_j\Rightarrow\delta_j\), then \(\neg R\));
cutting it against the current sequent, with cut formula \(\neg\delta_j\),
moves \(\delta_j\) into the succedent:
\[
 \frac{\;\Rightarrow\delta_j,\neg\delta_j
 \qquad \Gamma',\neg\delta_j\Rightarrow\Pi}
 {\Gamma'\Rightarrow\delta_j,\Pi}
 \qquad\text{yielding, after }r\text{ cuts,}\qquad
 \gamma_1,\ldots,\gamma_k\Rightarrow\delta_{j_1},\ldots,\delta_{j_r}.
\]
(5)~Weakening (\(W_L,W_R\)) adds the unused members of \(\Gamma,\Delta\) and
contraction (\(C_L,C_R\)) merges duplicates, giving exactly
\(\Gamma\Rightarrow\Delta\) as multisets.  The case \(m=0\) runs with no
step-(4) cuts and asserts that no structure and assignment satisfies all of
\(\Gamma\); empty \(\Gamma\) is likewise covered.  The route uses cut, which is
primitive in the displayed calculus; no cut-free claim is made, and extending
the sketch's translation induction to multi-succedent targets is unnecessary
since the reduction reuses it verbatim.  Landing: statement replaced at
4165--4170; steps (1)--(5) appended to the proof at 4171--4185, whose
translation subroutine is retained.  Status afterwards: \textit{Status: proof
sketch} while the rule-by-rule translation trees remain omitted; the theorem
inherits \textit{Status: complete proof} only if a separate work order supplies
them.

\emph{Propositional half.}  Attach \texttt{thm:prop-lk-completeness} to the
theorem at 2488 and replace its statement (2488--2491) by: for all finite
multisets \(\Gamma,\Delta\) of propositional formulas, the sequent
\(\Gamma\Rightarrow\Delta\) is derivable in \(\mathsf{LK}_{\mathrm{CPL}}\)
iff every valuation making every formula of \(\operatorname{supp}(\Gamma)\)
true makes some formula of \(\operatorname{supp}(\Delta)\) true ---
word-for-word the sequent-satisfaction convention of 2156--2161.  Left to
right is \texttt{thm:prop-lk-soundness} (2295).  Right to left is the
reduction (1)--(5) verbatim with \texttt{thm:prop-completeness} (2433) in
place of \texttt{thm:fol-completeness}, the translation sketch at 2493--2507
as the retained subroutine, and the same initial sequents and structural
rules (2178--2234); no choice principle enters (the propositional Lindenbaum
lemma runs on a countable enumeration).  Landing: statement replaced at
2488--2491; steps (1)--(5) appended to the proof at 2493--2507.  Status
afterwards: likewise \textit{Status: proof sketch} while its translation
trees remain omitted.

\paragraph{Acceptance.}
(1) The completeness theorem at 4165 quantifies over sequents
\(\Gamma\Rightarrow\Delta\) with both contexts arbitrary finite multisets, and
its semantic condition is word-for-word the satisfaction convention of
3579--3581, matching the scope of \texttt{thm:fol-lk-soundness}.  (2) The proof
displays the reduction chain of steps (1)--(5), including the
\(\;\Rightarrow\delta_j,\neg\delta_j\) construction and the empty-succedent
case.  (3) The status label matches the evidence actually present: sketch
while the translation trees are omitted, complete only if they are supplied.
(4) The propositional theorem at 2488 likewise quantifies over sequents with
both contexts arbitrary finite multisets, matches the scope of
\texttt{thm:prop-lk-soundness}, is proved by the same reduction without any
choice principle, and carries the same status discipline; after both
restatements, soundness and completeness are paired over the same sequent
class in each calculus, discharging both line pairs of the register row.
\endgroup

% Phase-2 fix designs for defects R-011, R-012, R-013, R-015, R-016.
% Section 2 fragment of the gap-analysis document.  Plain LaTeX; compiles
% under amsmath/amssymb; each design fits one page.  Line anchors refer to
% the v1.0 manuscript source/formal-logic-reference.tex (14780 lines).

\subsection*{R-011 --- CNF half of the normal-form theorem (v1.0 lines 1919--1966; defective half 1958--1965)}

\paragraph{Defect.} The CNF half of the theorem at 1919--1923 reads ``Apply
the disjunctive construction to \(\neg\varphi\), negate the result, and use
De Morgan equivalences'' (1958--1959).  Against the string-valued CNF
definition at 1878--1897 (a CNF formula is \(\top\), or the right-associated
conjunction \(N_1\) of clause interpretations \(\operatorname{form}(C_i)\),
each clause listed in strictly increasing code order), this fails: for valid
\(\varphi\) the DNF of \(\neg\varphi\) is \(\bot\), and \(\neg\bot\) is not a
CNF string; for variable-free false \(\varphi\) it yields \(\neg\top\)
instead of \(\bot=\operatorname{form}(\varnothing)\); and ``use De Morgan''
silently performs double-negation elimination and re-association that the
definition does not supply.

\paragraph{Repair.} Replace 1958--1965 by a dual direct construction
mirroring 1924--1956, plus one auxiliary lemma inserted after 1897.
\emph{Auxiliary lemma} (complete proof; two inductions on the definitional
recursions at 1855--1870 and 1883--1890): (i) a valuation \(w\) satisfies
\(\operatorname{form}(C)\) iff \(w\) satisfies some literal in \(C\)
(induction on \(|C|\) along \(D_k=L_k\), \(D_i=(L_i\lor D_{i+1})\); the case
\(C=\varnothing\) is the falsity of \(\bot\)); (ii) \(w\) satisfies \(N_1\)
iff \(w\) satisfies \(\operatorname{form}(C_j)\) for every \(j\le m\)
(induction on \(m\) along the \(N_i\) recursion).
\emph{Construction.}  Let \(p_1,\ldots,p_n\) be the variables of
\(\varphi\) in increasing index order; reuse the DNF half's valuations on
this finite set and the coincidence lemma (1925--1930).  For a valuation
\(v\) on \(\{p_1,\ldots,p_n\}\), let \(\bar\ell^{\,v}_i\) be \(\neg p_i\) if
\(v(p_i)=\mathsf{T}\) and \(p_i\) if \(v(p_i)=\mathsf{F}\), and put
\(K_v:=\{\bar\ell^{\,v}_1,\ldots,\bar\ell^{\,v}_n\}\).  The coding at
1849--1854 gives the variable of index \(i\) the code \(2i\) and its negation
the code \(2i+1\) --- both strictly increasing in \(i\) --- and the
enumeration \(p_1,\ldots,p_n\) is in increasing index order (the indices need
not form an initial segment), so any literal on \(p_j\) has smaller code than
any literal on \(p_k\) for \(j<k\): if their indices are \(i_j<i_k\), then
\(2i_j+1<2i_k\).  The definitional listing of \(K_v\) in increasing code
order is therefore exactly \(\bar\ell^{\,v}_1,\ldots,\bar\ell^{\,v}_n\):
\(\gamma_v:=\operatorname{form}(K_v)\) needs no re-ordering.  By lemma (i),
\(w\) falsifies \(\gamma_v\) iff \(w(p_i)=v(p_i)\) for all \(i\le n\).  Let
\(W_\varphi\) be the set of valuations on \(\{p_1,\ldots,p_n\}\) falsifying
\(\varphi\), enumerated lexicographically (\(\mathsf{F}<\mathsf{T}\), as at
1945--1948) as \(v_1,\ldots,v_m\).  Define \(C_\varphi:=\top\) if \(m=0\),
and otherwise \(C_\varphi:=N_1\) built from the clauses \(C_j:=K_{v_j}\) by
the recursion at 1883--1890.  Then \(C_\varphi\) is literally a CNF string
of 1878--1897, including the degenerate cases: valid \(\varphi\) gives
\(m=0\), the canonical \(\top\); variable-free true \(\varphi\) gives
\(n=m=0\), again \(\top\); variable-free false \(\varphi\) gives \(n=0\),
\(m=1\), \(K_{v_1}=\varnothing\), so
\(C_\varphi=\operatorname{form}(\varnothing)=\bot\), the CNF whose single
clause is empty.  \emph{Equivalence:} for a total valuation \(w\) with
restriction \(w_0\) to \(\{p_1,\ldots,p_n\}\) and \(m>0\): by lemma (ii)
and (i), \(w\) satisfies \(C_\varphi\) iff for every \(j\le m\) it satisfies
some literal of \(K_{v_j}\), iff \(w_0\ne v_j\) for all \(j\), iff
\(w_0\notin W_\varphi\), iff \(w_0\) satisfies \(\varphi\), iff
(coincidence) \(w\) satisfies \(\varphi\).  If \(m=0\), every valuation on
the variable set satisfies \(\varphi\), so by coincidence every \(w\)
satisfies both \(\varphi\) and \(\top\).  The theorem keeps its plain proof
environment; the DNF half is unchanged except that it may cite lemma (i).

\paragraph{Acceptance.} (1) For the three degenerate inputs (valid;
variable-free true; variable-free false) the constructed strings are
\(\top\), \(\top\), \(\operatorname{form}(\varnothing)\), each satisfying
the definition at 1878--1897 verbatim.  (2) The CNF half invokes only the
definitional recursions, the auxiliary lemma, and the coincidence lemma ---
no De Morgan, double negation, idempotence, or associativity.  (3) The
auxiliary lemma carries a complete proof by induction on both recursions.

\subsection*{R-012 --- Mix rule, measure, and honest corollary statuses (v1.0 lines 9505--9548, 9555--9565, 9575--9586)}

\paragraph{Defect.} Step 3 of the cut-elimination sketch (9528--9532)
appeals to ``the multicut (mix) rule'' and to elimination ``by the same
lexicographic measure,'' but no mix rule is displayed anywhere, the measure
at 9509--9511 is defined only for cut, and the equipollence of the cut and
mix calculi is unstated.  The corollaries at 9555 and 9575 carry unlabeled
proof environments, implying complete proofs, yet both consume the sketched
theorem, contravening the status policy at 313--319.

\paragraph{Repair.} (1) \emph{Display mix} after 9511, in the notation of
9200--9224, where \(A^{k}\) abbreviates \(k\) comma-separated occurrences
of \(A\).  For \(\mathsf{LK}\), with \(m,n\ge 1\) and \(\Delta,\Sigma\)
permitted to contain further occurrences of \(A\):
\[
  \frac{\Gamma\Rightarrow\Delta,A^{m}\qquad A^{n},\Pi\Rightarrow\Sigma}
       {\Gamma,\Pi\Rightarrow\Delta,\Sigma}\;\textsc{Mix}.
\]
For \(\mathsf{LJ}\) (single succedent, so the left premise's succedent
multiplicity \(m\) is fixed at one): from \(\Gamma\Rightarrow A\) and
\(A^{n},\Pi\Rightarrow C\) with \(n\ge1\),
conclude \(\Gamma,\Pi\Rightarrow C\).
(2) \emph{Equipollence lemma} (new; status: complete proof).  Cut is the
instance \(m=n=1\).  Conversely, mix is derived from the displayed
primitives: \(m-1\) applications of \textsc{CR} to the left premise,
\(n-1\) of \textsc{CL} to the right, then one cut.  Hence the cut and mix
formulations derive the same sequents (\(\mathsf{LJ}\): \textsc{CL} only),
so eliminating mix eliminates cut.
(3) \emph{Restate the measure} (replace 9505--9511): a mix instance with
mix formula \(A\) and immediate subderivations \(D_1,D_2\) has rank
\(r(A)\) and height \(h(D_1)+h(D_2)\); the reduction is controlled by the
lexicographic order on \(\langle r(A),\,h(D_1)+h(D_2)\rangle\), with cut
the case \(m=n=1\).
(4) \emph{Preface the sketch} (after 9519) with: ``By the equipollence
lemma, first replace every cut in the given derivation by the mix instance
\(m=n=1\); the reduction below eliminates mixes, and steps 1, 2, 5, and 6
apply verbatim with `mix' for `cut' and the mix formula for the cut
formula.''  \emph{Rewrite sketch step 3} (9528--9532): if a premise of a mix ends
in a contraction of the mix formula (\textsc{CL} above the right premise,
\textsc{CR} above the left), replace that premise by the contraction's
premise and raise the corresponding multiplicity by one: rank unchanged,
height decreased by one.  Add to step 4: when the mix formula is principal
on both sides with multiplicities \((m,n)\), the principal reduction yields
mixes on proper subformulas (smaller rank) plus residual mixes on remaining
occurrences at equal rank but smaller height; the full case table remains
the omitted verification (citation at 9548 unchanged).
(5) \emph{Honest statuses.}  Append \textit{Status: proof sketch.}\ to the
theorem statement (9513--9517).  Split the corollary at 9555: a new lemma
``every cut-free propositional \(\mathsf{LK}\)- or
\(\mathsf{LJ}\)-derivation has the subformula property'' takes the existing
induction (9560--9565) as a complete proof; the corollary (existence of a
cut-free derivation with the property) is restated with \textit{Status:
proof sketch.}, its proof converted to the sketch environment with the
dependency sentence modeled on 7728--7729: ``This derivation uses the
preceding proof sketch of the cut-elimination theorem; a complete proof
supplies its reduction case table.''  The corollary at 9575 keeps its
argument verbatim, converted to the sketch environment with the same
dependency sentence and \textit{Status: proof sketch.}\ on its statement.
When cut elimination is proved in full, exactly these two status lines and
two sentences flip to complete proof.

\paragraph{Acceptance.} (1) Mix is displayed for both calculi with explicit
multiplicities \(m,n\ge1\), and the equipollence lemma is proved using only
\textsc{CL}, \textsc{CR}, and \textsc{Cut} as displayed at 9204--9224 and
9307--9375.  (2) The sketch opens by converting every cut to a mix via the
equipollence lemma, so no step manipulates an object the measure does not
cover; step 3 names only displayed rules and exhibits the measure decrease
(rank equal, height \(-1\)).  (3) The theorem and both
corollaries carry \textit{Status: proof sketch.}; the cut-free-subformula
lemma and the equipollence lemma carry \textit{Status: complete proof.}; no
complete-proof item in the chapter cites the cut-elimination theorem.

\subsection*{R-013 --- Repetition and vacuous discharge in \(\mathsf{ND}_{\mathsf{CPL}}\) (v1.0 lines 2006--2021, 2351--2379, 2443--2455; first-order model 3799--3812)}

\paragraph{Defect.} Line 2020--2021 licenses, in prose only, ``Repetition
of an already available formula is permitted,'' though the calculus is
defined (2006--2018) as a tree generated by the displayed rules and no
repetition rule is displayed; and the definition never says whether a
discharging rule may discharge zero occurrences.  Clause 2 of the closure
lemma (2360--2362) and the completeness theorem (2446--2449) apply
\(\neg I\)/\(\mathsf{RAA}\) to derivations that may contain no open
occurrence of the discharged assumption, and clause 5 (2374--2375) says
``assume \(\varphi\), repeat \(\psi\), and apply \(\to I\).''

\paragraph{Repair.} Eliminate the uses, matching 3799--3812.
(a) Delete the repetition sentence (2020--2021).  Append to the derivation
definition (after 2011): ``An application of \(\to I\), \(\neg I\), or
\(\mathsf{RAA}\) carries a label \(i\) and discharges exactly the open
occurrences marked \(i\) of the bracketed assumption above its premise; an
application of \(\lor E\) carries two labels \(i,j\), as displayed, and
discharges exactly the open occurrences of the first disjunct marked \(i\)
above its second premise and the open occurrences of the second disjunct
marked \(j\) above its third premise.  Every label carried by such an
application discharges at least one occurrence.  Marking any set of open
occurrences of one formula with a common label fresh from every label in
the derivation preserves the derivation and the occurrences discharged by
each of its rule applications, since no rule application carries the fresh
label; this merging relabeling is distinct from the injective renamings
above, and is the operation used when a newly added root inference is to
discharge several open occurrences at once.  There is no repetition step:
where an informal presentation repeats a formula, the tree reuses the node
already present.''
(b) Closure clause 2 (2358--2362), mirroring 3803--3807: if
\(\varphi\notin\Delta\), maximality yields a derivation \(D\) of \(\bot\)
with open assumptions in \(\Delta\cup\{\varphi\}\); because \(\Delta\) is
consistent, at least one open occurrence of \(\varphi\) remains in \(D\);
mark all open occurrences of \(\varphi\) in \(D\) with a common fresh label
\(i\) (the merging relabeling of (a)), and apply \(\neg I\) with label
\(i\) at the root, discharging every such occurrence, to obtain
\(\Delta\vdash\neg\varphi\); clause 1 gives \(\neg\varphi\in\Delta\).
(c) Closure clause 5 (2374--2375), replacing repetition by a conjunction
detour: if \(\psi\in\Delta\), assume \(\varphi\) with label \(i\) and take
\(\psi\) as an open assumption; \(\land I\) gives \(\varphi\land\psi\);
\(\land E_2\) gives \(\psi\); \(\to I\) discharges the single open
occurrence marked \(i\), concluding \(\varphi\to\psi\) with open
assumptions in \(\Delta\).  The discharge is of an occurrence genuinely
present, and only displayed rules occur.
(d) Completeness (2446--2449): suppose a derivation \(D\) of \(\bot\) has
open assumptions in \(\Gamma\cup\{\neg\varphi\}\).  If no open occurrence
of \(\neg\varphi\) appeared in \(D\), then \(D\) would witness
\(\Gamma\vdash\bot\), and \(\bot E\) would give \(\Gamma\vdash\varphi\),
contrary to the supposition \(\Gamma\nvdash\varphi\).  So at least one open
occurrence remains; mark all open occurrences of \(\neg\varphi\) with a
common fresh label \(i\) (the merging relabeling of (a)) and apply
\(\mathsf{RAA}\) with label \(i\) at the root, discharging every such
occurrence, to obtain \(\Gamma\vdash\varphi\) --- again contrary to
the supposition.  Hence \(\Gamma\cup\{\neg\varphi\}\) is consistent.
(e) Sweep every displayed \(\lor E\)/\(\to I\)/\(\neg I\)/\(\mathsf{RAA}\)
application in the ND chapters, examples, exercises, and solutions for
compliance (the examples at 2077--2110 already comply; the intuitionistic
calculus, defined at 2106--2109 by omitting \(\mathsf{RAA}\), inherits the
convention; 3799--3812 stands unchanged and now matches, its ``discharging
every such occurrence'' being an instance of the merging relabeling
followed by a single labeled application).

\paragraph{Acceptance.} (1) ``Repetition'' no longer occurs in the
\(\mathsf{ND}_{\mathsf{CPL}}\) section, and the derivation definition
states explicitly both the at-least-one-occurrence discharge convention
(with \(\lor E\) carrying two labels) and the fresh-label merging
relabeling used by the closure and completeness proofs.
(2) The proofs of closure clauses 2 and 5 and of completeness contain,
respectively, the non-vacuity argument and the conjunction detour, and
invoke no rule outside the display at 2022--2065.  (3) An audit of all
discharging-rule applications in displayed derivations finds none that
discharges zero occurrences.

\subsection*{R-015 --- Reduction of arbitrary embeddings to substructure inclusions (v1.0 lines 5989--5992, 6009--6012; supporting definitions 5042--5057, 5093--5103)}

\paragraph{Defect.} Model completeness is defined at 5989 as ``every
embedding between its models is elementary,'' but the \(\mathsf{DLO}\)
corollary's sketch verifies only inclusions: ``If
\(\mathcal M\subseteq\mathcal N\) are models of \(\mathsf{DLO}\)
\(\ldots\) This is model completeness'' (6009--6012).  The passage from
inclusions to arbitrary embeddings is missing, and the manuscript nowhere
proves that satisfaction is isomorphism-invariant or that an embedding is
an isomorphism onto its image.

\paragraph{Repair.} Insert Lemmas 1--2 after 5125 (end of the
Substructures section) and Lemma 3 after 5992; patch 6011--6012.
\emph{Lemma 1 (Isomorphism invariance; complete proof).}  If
\(h:\mathcal M\to\mathcal N\) is an isomorphism, then for every formula
\(\varphi(x_1,\ldots,x_k)\) and all \(a_1,\ldots,a_k\in M\):
\(\mathcal M\models\varphi(\bar a)\iff\mathcal N\models\varphi(h\bar a)\);
in particular isomorphic structures satisfy the same sentences.
\emph{Route:} induction on \(\varphi\); atomic case by the
term-preservation lemma (5060--5076), injectivity for equality atoms, and
preservation with reflection (5042--5052) for relation atoms; Boolean cases
by induction hypothesis; for \(\exists y\), a witness \(a'\in M\) maps to
\(h(a')\), and conversely a witness \(b\in N\) is \(h(a')\) by
surjectivity; \(\forall\) dually.
\emph{Lemma 2 (Image of an embedding; complete proof).}  If
\(h:\mathcal M\to\mathcal N\) is an embedding, then \(h(M)\) contains
\(c^{\mathcal N}\) for every constant and is closed under every
\(f^{\mathcal N}\); the structure \(h[\mathcal M]\) with domain \(h(M)\),
constants \(c^{\mathcal N}\), the restrictions of the \(f^{\mathcal N}\),
and relations \(R^{\mathcal N}\) restricted to \(h(M)\) is a substructure
of \(\mathcal N\) per 5093--5103; and \(h\) with codomain \(h(M)\) is an
isomorphism \(\mathcal M\to h[\mathcal M]\).  \emph{Route:}
\(c^{\mathcal N}=h(c^{\mathcal M})\in h(M)\) and
\(f^{\mathcal N}(h\bar a)=h(f^{\mathcal M}(\bar a))\in h(M)\) by the
homomorphism clauses; the inclusion preserves and reflects atomic formulas
because interpretations are restrictions, so \(h[\mathcal M]\subseteq
\mathcal N\); the corestricted \(h\) is injective (given), surjective by
construction, preserves and reflects relations since
\(R^{h[\mathcal M]}=R^{\mathcal N}\cap h(M)^{k}\) while \(h\) reflects
\(R^{\mathcal N}\); hence a surjective embedding, an isomorphism per
5055--5057.
\emph{Lemma 3 (Reduction to inclusions; complete proof).}  If
\(\mathcal M\preccurlyeq\mathcal N\) whenever
\(\mathcal M\subseteq\mathcal N\) are both models of \(T\), then \(T\) is
model complete.  \emph{Route:} an embedding \(h:\mathcal M\to\mathcal N\)
between models of \(T\) factors (Lemma 2) as the isomorphism
\(h_0:\mathcal M\to h[\mathcal M]\) followed by the inclusion
\(h[\mathcal M]\subseteq\mathcal N\); Lemma 1 gives
\(h[\mathcal M]\models T\); the hypothesis gives
\(h[\mathcal M]\preccurlyeq\mathcal N\); then for all \(\varphi,\bar a\):
\(\mathcal M\models\varphi(\bar a)\iff h[\mathcal M]\models
\varphi(h\bar a)\) (Lemma 1) \(\iff\mathcal N\models\varphi(h\bar a)\)
(elementary inclusion).  Hence \(h\) is elementary.
\emph{Patch 6011--6012:} replace ``This is model completeness.''\ by ``By
the reduction lemma, every embedding between models of \(\mathsf{DLO}\) is
elementary; this is model completeness.''  The corollary keeps
\textit{Status: proof sketch.}\ (its remaining omission is the literal-case
verification named at 6013--6014, not the embedding reduction).

\paragraph{Acceptance.} (1) The three lemmas appear with complete proofs
using only 5022--5057, 5093--5103, and the term-preservation lemma; Lemma
1's quantifier case invokes surjectivity explicitly, and Lemma 2 verifies
closure of \(h(M)\) before forming \(h[\mathcal M]\).  (2) The
\(\mathsf{DLO}\) sketch derives model completeness by citing Lemma 3.
(3) No passage concludes model completeness from inclusions alone.

\subsection*{R-016 --- Extension clause in the back-and-forth test (v1.0 lines 6068--6093; corollary 6097--6110)}

\paragraph{Defect.} The back-and-forth theorem (6068--6081) concludes only
``\(\mathcal M\cong\mathcal N\),'' but the corollary's proof asserts ``The
resulting isomorphism extends the original map'' (6108--6109) --- a
property of the proof's construction, which begins at an arbitrary
\(f_0\in\mathcal F\), not of the stated theorem.  The corollary's use is
therefore unlicensed.

\paragraph{Repair.} \emph{Restated theorem (6068--6081).}  Let
\(\mathcal M,\mathcal N\) be countable structures and \(\mathcal F\) a
nonempty set of finite partial isomorphisms from \(\mathcal M\) to
\(\mathcal N\) such that for every \(f\in\mathcal F\): (1) for every
\(a\in M\), some \(g\in\mathcal F\) extends \(f\) with \(a\) in its domain;
(2) for every \(b\in N\), some \(g\in\mathcal F\) extends \(f\) with \(b\)
in its range.  Then for every \(f_0\in\mathcal F\) there is an isomorphism
\(h:\mathcal M\to\mathcal N\) with \(f_0\subseteq h\) (extension as sets of
ordered pairs); in particular \(\mathcal M\cong\mathcal N\).
\emph{Proof edits (6083--6093):} replace ``begin with any
\(f_0\in\mathcal F\)'' by ``let \(f_0\) be the given member of
\(\mathcal F\)''; the recursion and union argument are unchanged; append
``Since \(f_0\subseteq f_1\subseteq\cdots\) and \(h=\bigcup_k f_k\), the
isomorphism extends \(f_0\).''  \textit{Status: complete proof}\ retained.
\emph{Corollary patch (6104--6110).}  Replace the final sentence by an
explicit application: the set \(\mathcal F\) of finite elementary maps from
\(\mathcal M\) to itself extending the given map \(f\) is nonempty
(\(f\in\mathcal F\)); every member is a finite partial isomorphism, since
elementary maps preserve and reflect atomic formulas and reflect equality,
hence are injective; clause (1) holds within \(\mathcal F\) because the
saturation-homogeneity theorem (5567--5571) extends any \(g\in\mathcal F\)
by any prescribed domain element to an elementary map still extending
\(f\); clause (2) holds by applying the same theorem to \(g^{-1}\) (the
inverse of a finite elementary self-map is elementary) and inverting the
result, which extends \(g\supseteq f\) and has the prescribed element in
its range.  The restated theorem applied with \(f_0:=f\) yields an
isomorphism \(h:\mathcal M\to\mathcal M\) with \(f\subseteq h\): an
automorphism extending \(f\).  \textit{Status: complete proof}\ retained.
The forward reference at 5563--5564 remains accurate.

\paragraph{Acceptance.} (1) The theorem statement contains the quantified
extension clause ``for every \(f_0\in\mathcal F\) there is an isomorphism
\(h\supseteq f_0\).''  (2) Its proof starts the recursion at the given
\(f_0\) and derives \(f_0\subseteq h\) from the chain union.  (3) The
corollary invokes the extension clause with \(f_0:=\) the given map, and
both results retain \textit{Status: complete proof.}

\endgroup

\section{Depth-gap inventory}
\label{sec:depth}

Gaps \emph{within} the eleven existing parts, each assigned a target
home in the new architecture (§\ref{sec:architecture}) and a
\textbf{treatment level}, which is a ceiling, not a floor to exceed:

\begin{center}\small
\begin{tabular}{@{}lp{4.6in}@{}}
\toprule
\textbf{E} & full reference entry (definitions, adequacy theorems with status, tables) \\
\textbf{T} & table row(s) only, in a master or per-chapter table \\
\textbf{S} & statement with exact hypotheses + status + citation; no development \\
\textbf{P} & pointer: one sentence and a source \\
\bottomrule
\end{tabular}
\end{center}

\begingroup\scriptsize\setlength{\tabcolsep}{3pt}
\begin{longtable}{@{}p{0.62in}p{2.6in}p{0.75in}p{0.22in}p{0.42in}p{1.85in}@{}}
\toprule
\textbf{ID} & \textbf{Item} & \textbf{Target} & \textbf{Tr} & \textbf{Prio} & \textbf{Notes} \\
\midrule
\endfirsthead
\toprule
\textbf{ID} & \textbf{Item} & \textbf{Target} & \textbf{Tr} & \textbf{Prio} & \textbf{Notes} \\
\midrule
\endhead
\midrule\multicolumn{6}{r}{\emph{continued}}\\\endfoot
\bottomrule\endlastfoot
D-P1-001 & Zorn's lemma: statement, proof from AC, equivalence noted & P6 & E & must & absent book-wide; felt at 3978 where it is the standard tool for Henkin--Lindenbaum; cross-cutting list (Zorn/AC-equivalents) \\
D-P1-002 & AC-equivalents theorem and master table (AC \(\Leftrightarrow\) well-ordering \(\Leftrightarrow\) Zorn \(\Leftrightarrow\) Tukey; BPI strictly weaker) & P6 & T & must & cross-cutting list; AC-equivalents table required by target architecture sec. 6 \\
D-P1-003 & Cardinal trichotomy: statement and AC-dependence & P6 & S & must & used at 672 but never stated \\
D-P1-004 & K\"onig's lemma and fan theorem, with the choice-free route to propositional compactness & P6 & E & must & absence blocks choice-free propositional compactness; ties to the broken forward promise at 546--548 (countable-language completeness) \\
D-P1-005 & BPI/Ultrafilter Lemma \(\Leftrightarrow\) propositional compactness & P2.1 & E & should & Ultrafilter Lemma present at 704--724 but never linked to compactness \\
D-P1-006 & Countable choice: statement, status, and where the book uses it & P6 & S & should & raised in an exercise, resolved nowhere \\
D-P2-001 & Hilbert-style calculus for CPL with deduction theorem and ND/LK/Hilbert equivalence and comparison & P2.1 & E & must & no Hilbert system book-wide; cross-cutting list \\
D-P2-002 & Post's functional-completeness classification (five maximal clones + negative results + Sheffer stroke/NOR) & P2.1 & E & must & one positive basis only in v1.0; cross-cutting list; Post's lattice figure in apparatus register \\
D-P2-003 & Duality theorem for propositional logic & P2.1 & E & should & cross-cutting list \\
D-P2-004 & Decidability of CPL stated as a theorem & P2.1 & S & must & never stated despite the truth-table method being developed \\
D-P2-005 & SAT NP-completeness statement in the propositional chapter with forward reference to P4 & P2.1 & S & must & Cook cited only at 12118 with no forward ref from Part II \\
D-P2-006 & Propositional interpolation theorem & P2.1 & E & should &  \\
D-P2-007 & Tseitin CNF transformation: statement and cross-reference from the CNF section & P2.1 & S & should & exists at 11945 but unreferenced from the CNF section \\
D-P2-008 & Horn-SAT and 2-SAT tractable classes & P2.1 & E & should & also feeds P9/P12 SAT material \\
D-P2-009 & Lindenbaum--Tarski algebra and Boolean-valued semantics for CPL & P10 & E & must & quotient machinery from Ch. 1 never exploited; also listed in B3 algebraic/categorical and A(d) -- coordinate with new-content register \\
D-P2-010 & Substitution and replacement congruence theorems promoted from exercise to stated results & P2.1 & S & should &  \\
D-P2-011 & Propositional logic over uncountable languages (compactness/completeness beyond the fixed countable alphabet) & P2.1 & S & could & alphabet fixed countable at 1565 \\
D-P3-001 & Hilbert-style FOL system with deduction theorem & P2.2 & E & must & cross-cutting list \\
D-P3-002 & Classical analytic tableaux for propositional and first-order logic with countermodel extraction and completeness & P2.2 & E & must & absent book-wide (only the many-valued matrix tableau at 11255); also listed in B3 automation and cross-cutting list \\
D-P3-003 & Herbrand's theorem proper: Herbrand universe/base/structure/expansion and the validity form & P2.2 & E & must & only a cited clause corollary at 4799; cross-cutting list; link to resolution \\
D-P3-004 & Craig interpolation and Beth definability: statements and forward references in the first-order chapter & P2.2 & S & must & currently only cited in Part VIII (9627/9645) with no forward reference; also flagged in B2 model theory \\
D-P3-005 & Robinson joint consistency theorem & P3 & E & should &  \\
D-P3-006 & Decidable fragments of FOL: catalog table (monadic; FO\(^2\); guarded; Bernays--Sch\"onfinkel) & P2.2 & T & must & monadic FOL one aside at 6875; master decidability table in P14 \\
D-P3-007 & Many-sorted first-order logic: syntax, semantics, reduction to one-sorted & P2.4 & E & should & only a second-order translation target at 10204 \\
D-P3-008 & Empty domains and the inclusive/free-logic convention with pointer to the free-logic entry card & P8.4 & S & should & nonempty domain imposed at 3028; full free-logic card is new-content N-b; also listed in B3 nonmonotonic/philosophical \\
D-P3-009 & Infinitary logic \(L_{\omega_1\omega}\) overview and pointer & P2.4 & S & could & one mention at 9602 \\
D-P3-010 & Definitional extensions: conservativity theorem stated and proved & P2.2 & E & must & unproved in example 3415--3418; Beth definability leans on it; coordinate with defect register \\
D-P3-011 & Equality metatheory consolidated: free-standing congruence lemma and equality-axiom theorems & P2.2 & E & must & symmetry an unlabeled example; transitivity an exercise \\
D-P3-012 & Clausal-form theorem composing prenex + Skolem + CNF & P2.2 & S & should & currently left to an exercise \\
D-P3-013 & Compactness applications suite: nonstandard models, Skolem paradox discussion, upward LS via elementary diagrams & P2.2 & E & should & one corollary only in v1.0 \\
D-P3-014 & First-order subformula definition and subformula property & P2.2 & S & must & Sub is propositional-only; subformula property is a stated learning outcome \\
D-P4-001 & Ehrenfeucht--Fra\"{\i}ss\'e games with inexpressibility applications & P3 & E & must & absent book-wide; listed in both B (first-order) and B2 (model theory); also A(a) \\
D-P4-002 & Preservation theorems ({\L}o\'s--Tarski; Chang--{\L}o\'s--Suszko) & P3 & E & should & absent despite a full ultraproduct section; also A(a) \\
D-P4-003 & Keisler--Shelah isomorphism theorem (statement) & P3 & S & should &  \\
D-P4-004 & Compactness via ultraproducts drawn as a corollary of {\L}o\'s & P3 & S & should &  \\
D-P4-005 & Prime and atomic models with the omitting-types connection & P3 & E & must & undefined though Ryll-Nardzewski is stated \\
D-P4-006 & Vaught two-cardinal theorem (statement) & P3 & S & could &  \\
D-P4-007 & Indiscernibles and Ehrenfeucht--Mostowski models (statement level) & P3 & S & should &  \\
D-P4-008 & Stability, forking, and simplicity survey & P3 & S & should & research-level; survey ceiling per treatment policy; also A(a) \\
D-P4-009 & NIP and o-minimality survey & P3 & S & should & research-level survey; also A(a) \\
D-P4-010 & Finite model theory: 0--1 laws and Fagin's theorem & P3 & S & should & only Trakhtenbrot in v1.0; descriptive-complexity link to P9; also A(a) \\
D-P4-011 & Computable model theory pointer & P3 & P & could & also A(a) \\
D-P4-012 & Lindstr\"om's theorem repositioned to model theory with abstract model theory context & P3 & S & should & currently placed in Part IX, not IV \\
D-P5-001 & \(\mu\)-recursion and partial recursive functions with the TM \(\equiv\) \(\mu\)-recursive equivalence theorem & P4 & E & must & partial computable is machine-only in v1.0 so the standard equivalence cannot be stated; also A(a) \\
D-P5-002 & Church--Turing thesis discussion & P4 & S & must &  \\
D-P5-003 & Turing degrees as a structure (poset, join, jump on degrees) & P4 & E & should & only \(\le_T\) and the jump in v1.0; also A(a) \\
D-P5-004 & Post's problem, finite-injury priority method, Friedberg--Muchnik (statement + method sketch) & P4 & S & should & also A(a) \\
D-P5-005 & C.e.-set structure theory: simple and creative sets, Myhill isomorphism, Rice--Shapiro & P4 & S & should & also A(a) \\
D-P5-006 & Kolmogorov complexity: invariance, incompressibility, halting-probability pointer & P4 & S & could & also A(a) \\
D-P5-007 & Complexity theory expansion: PSPACE, time/space hierarchy theorems, PH, complexity-class ladder table & P4 & T & must & only a short section at 12105; descriptive complexity via P9; also A(a) \\
D-P6-001 & Presburger arithmetic: decidability by quantifier elimination (statement + proof route) & P5 & S & must & absent book-wide; flagged in both B (first-order) and B2 (arithmetic); also A(a) \\
D-P6-002 & Tennenbaum's theorem (statement) & P5 & S & should & also A(a) \\
D-P6-003 & Nonstandard models of PA: order type \(\mathbb{N}+\mathbb{Z}\cdot\mathbb{Q}\), overspill lemma & P5 & E & should &  \\
D-P6-004 & Paris--Harrington and Goodstein independence results (statements) & P5 & S & should & also A(a) \\
D-P6-005 & Interpretability: relative interpretation with examples & P5 & S & could & also A(a) \\
D-P6-006 & Provable \(\Sigma_1\)-completeness stated and proved & P5 & E & must & used implicitly in the G\"odel II machinery, never stated \\
D-P6-007 & Provability logic GL: axioms, Kripke semantics, fixed-point theorem; Solovay arithmetical completeness (S) & P5 & E & must & natural companion to the proved L\"ob theorem; also flagged in B3 modal and A(b); entry card in P8.2 \\
D-P6-008 & Reverse mathematics: Z\(_2\) subsystems and calibration table (RCA\(_0\), WKL\(_0\), ACA\(_0\), ATR\(_0\), \(\Pi^1_1\)-CA\(_0\)) & P5 & T & should & also A(a) \\
D-P7-001 & Ordinal arithmetic: sums, products, exponentiation, Cantor normal form, transfinite induction vs recursion & P6 & E & must & \(\varepsilon_0\)/Cantor normal form used at 9854--9876 but never developed; listed in both B (foundations) and B2 (set theory); also A(a) \\
D-P7-002 & Cardinal arithmetic beyond \(2^\kappa\): sums, products, exponentiation, \(\aleph\)/\(\beth\) hierarchies, regular vs singular; CH/GCH statements and pointer & P6 & E & must & listed in both B (foundations) and B2 (set theory) \\
D-P7-003 & Cofinality and K\"onig's theorem & P6 & E & must &  \\
D-P7-004 & Easton's theorem (statement) & P6 & S & should &  \\
D-P7-005 & Martin's axiom: statement and sample consequences & P6 & S & should & also A(a) \\
D-P7-006 & Large-cardinal ladder: inaccessible through measurable and beyond, with consistency-strength chart & P6 & T & must & 0 occurrences of inaccessible in v1.0 -- needed at the ctm hypothesis 8767; research-level so table + statements ceiling; also A(a) \\
D-P7-007 & L\'evy hierarchy, reflection principle, absoluteness (Shoenfield statement) & P6 & E & must & needed for forcing and L discipline; also A(a) \\
D-P7-008 & Condensation and the structure of L (condensation statement; fine-structure pointer) & P6 & S & should & also A(a) \\
D-P7-009 & Clubs, stationary sets, Fodor's lemma & P6 & E & should & also A(a) \\
D-P7-010 & \(\Delta\)-system lemma and Ramsey's theorem & P6 & E & must & \(\Delta\)-system lemma needed for ccc of Cohen forcing (D-P7-012); also A(a) \\
D-P7-011 & Boolean-valued models (statement level) and product/iterated/proper forcing pointers & P6 & S & should & iterated/proper forcing NOT-IN-SCOPE beyond pointer per roadmap; also A(a) \\
D-P7-012 & ccc preservation of cardinals and ccc of the Cohen poset, proved & P6 & E & must & currently cited only; feeds the honest Con(ZFC+\(\neg\)CH) treatment \\
D-P7-013 & Descriptive set theory: Borel and projective hierarchies, regularity properties (survey) & P6 & S & should & research-level survey; also A(a) \\
D-P7-014 & Determinacy: Borel determinacy statement, AD vs AC (survey) & P6 & S & should & research-level survey; also A(a) \\
D-P8-001 & Completeness of LK for arbitrary sequents in the proof-theory part & P7 & E & must & Part VIII has only a soundness sketch \\
D-P8-002 & LK \(\leftrightarrow\) ND equivalence theorem & P7 & E & must &  \\
D-P8-003 & Midsequent theorem and Herbrand's theorem via cut elimination & P7 & S & should & also A(a) \\
D-P8-004 & Craig interpolation proved by Maehara's method & P7 & E & must & Maehara named but not carried out; cross-cutting list: Craig/Beth proved rather than cited \\
D-P8-005 & Gentzen consistency-proof components: derivation coding, ordinal assignment, decrease lemma & P7 & S & should & current treatment is an honest survey with no coding/assignment/decrease lemma \\
D-P8-006 & Ordinal-analysis calibration table beyond PA with proof-theoretic ordinal defined & P7 & T & should & v1.0 declines to define proof-theoretic ordinal; ordinal arithmetic must land before the ordinal-analysis table \\
D-P8-007 & Proof complexity: Frege systems, resolution lower bounds (pigeonhole), extended resolution & P7 & S & should & flagged in both B (propositional) and B2 (proof theory); also A(a) \\
D-P9-001 & Second-order/HOL proof calculus displayed (comprehension and choice schemas as a calculus) with Henkin completeness route & P2.3 & E & must & explicit decision in v1.0 to omit; also A(a) \\
D-P9-002 & Analytical hierarchy (\(\Sigma^1_n\)): table and basic results & P4 & T & should & also A(a) \\
D-P9-003 & Untyped \(\lambda\)-calculus with Church--Rosser/confluence & P7 & E & must & absent book-wide while \(\beta\)-reduction, subject reduction, and normalization are present; also A(a) \\
D-P9-004 & System F: definition and strong normalization (statement) & P7 & S & should & also A(a) \\
D-P10-001 & Canonical models constructed with completeness proofs for K, T, S4, S5 & P8.2 & E & must & named in sketches, never constructed; cross-cutting list \\
D-P10-002 & Filtration, finite model property, and decidability for the basic modal systems & P8.2 & E & must & filtration only an omitted obligation; no FMP/decidability result for any modal logic in v1.0; cross-cutting list \\
D-P10-003 & Frame-correspondence table with the standard translation; Sahlqvist theorem as statement & P8.2 & E & must & no correspondence theory in v1.0; axiom \(\leftrightarrow\) frame-condition \(\leftrightarrow\) algebra master table in P14 \\
D-P10-004 & Bisimulation and the van Benthem characterization theorem & P8.2 & E & should &  \\
D-P10-005 & Modal systems lattice completion: entry cards for D, B, K4, KD45, Grz & P8.2 & E & must & GL handled at D-P6-007 \\
D-P10-006 & Tense logic and multimodal fusions & P8.3 & S & should &  \\
D-P10-007 & Neighborhood semantics for non-normal modal logics & P8.2 & S & could &  \\
D-P10-008 & Quantified modal logic: Barcan formulas, constant vs varying domains & P8.2 & E & must & absent book-wide; also A(b) -- coordinate with new-content register \\
D-P10-009 & CTL, CTL*, and \(\mu\)-calculus entry cards & P8.3 & E & must & also A(c); complexity fields via D-P11-010 \\
D-P10-010 & B\"uchi automata and automata-theoretic LTL model checking & P9 & S & must & also A(c) model checking -- coordinate with new-content register \\
D-P10-011 & Dynamic epistemic logic, public announcement, distributed knowledge & P8.3 & S & should &  \\
D-P10-012 & Common-knowledge axiomatization: induction rule, completeness and non-compactness caveat & P8.3 & E & must & C\(_G\) has language + semantics but no axioms at 11179--11191; an exercise uses it; defect-register overlap (C3 finding 3) \\
D-P10-013 & Heyting algebras with algebraic completeness of IPL & P8.1 & E & must & also A(d) -- coordinate with new-content register \\
D-P10-014 & Topological semantics for IPL & P8.1 & S & should &  \\
D-P10-015 & BHK reading and Kleene realizability survey & P8.1 & S & should &  \\
D-P10-016 & Existence property for intuitionistic logic & P8.1 & S & must & cross-cutting list; disjunction property already proved in v1.0 \\
D-P10-017 & First-order intuitionistic logic: Kripke semantics with expanding domains, completeness statement, CD axiom & P8.1 & E & must & chapter is propositional-only; cross-cutting list \\
D-P10-018 & Intermediate logics KC and LC (G\"odel--Dummett as intermediate logic) & P8.1 & S & should & G\"odel t-norm semantics under D-P10-021 \\
D-P10-019 & Markov's principle: statement and status & P8.1 & S & could &  \\
D-P10-020 & {\L}ukasiewicz logics, finite- and infinite-valued: matrices, axiomatization, completeness statements & P8.5 & E & must &  \\
D-P10-021 & t-norm fuzzy logics: BL, G\"odel, product & P8.5 & S & should &  \\
D-P10-022 & MV-algebras and Chang completeness (statement) & P10 & S & should & also A(d) \\
D-P10-023 & Post logics (statement level) & P8.5 & S & could &  \\
D-P10-024 & FDE/Belnap four-valued logic entry & P8.5 & E & should &  \\
D-P10-025 & Weak Kleene and Bochvar logics & P8.5 & S & could &  \\
D-P10-026 & Supervaluations & P8.5 & S & should & also A(b) vagueness -- coordinate with new-content register \\
D-P10-027 & Registered facts: K3 has no tautologies; LP tautologies = classical; K3--LP duality & P8.5 & S & must & never registered in v1.0; cheap entry-card fields \\
D-P10-028 & Full classical linear logic: additives, multiplicatives, units, exponentials, linear negation, one-sided sequents, cut-elimination statement & P8.6 & E & must & cross-cutting list \\
D-P10-029 & Phase semantics: soundness and completeness statement & P8.6 & S & should &  \\
D-P10-030 & Coherence spaces and proof nets pointers & P8.6 & P & could & proof nets pointer-only per scope ruling \\
D-P10-031 & Display, deep-inference, hypersequent, and labelled calculi survey pointer & P7 & P & could &  \\
D-P10-032 & Relevance systems R and E as named axiom systems with the variable-sharing theorem & P8.6 & S & should & omission was deliberate in v1.0 but leaves no variable-sharing theorem \\
D-P10-033 & Lambek calculus entry & P8.6 & S & should & only named in v1.0 \\
D-P10-034 & Residuated lattices as the unifying substructural semantics & P10 & S & should & one definition only in v1.0; also A(d) \\
D-P10-035 & Nonmonotonic logics: default logic, circumscription, autoepistemic logic & P8.7 & S & should & ASP under D-P11-007; also A(b) -- coordinate with new-content register \\
D-P10-036 & AGM revision postulates table & P8.7 & T & should & also A(b) \\
D-P10-037 & Deontic logic: SDL and its paradoxes & P8.4 & S & should & also A(b) -- coordinate with new-content register \\
D-P10-038 & Counterfactuals: Stalnaker--Lewis sphere semantics & P8.4 & S & should & also A(b) -- coordinate with new-content register \\
D-P10-039 & Boolean algebras chapter: axioms, filters and ultrafilters, atomic/atomless, representation & P10 & E & must & also A(d) -- coordinate with new-content register \\
D-P10-040 & Stone duality & P10 & E & should & also A(d) \\
D-P10-041 & Cylindric and relation algebras pointer & P10 & P & could & pointer-level per A(d) \\
D-P10-042 & CCC \(\leftrightarrow\) STLC categorical semantics & P10 & S & should & also A(d) -- coordinate with new-content register \\
D-P10-043 & Hyperdoctrines and topos logic pointer & P10 & P & could & NOT-IN-SCOPE beyond CCC/topos overview per roadmap \\
D-P11-001 & CDCL: clause learning, watched literals, restarts, DRAT proofs & P12 & E & must & chapter presents 1962 DPLL only; also A(c) -- coordinate with new-content register \\
D-P11-002 & Unification: Martelli--Montanari algorithm with the unification theorem & P12 & E & must & mgu defined, occurs check in prose, no algorithm or theorem; also A(c) \\
D-P11-003 & Superposition calculus: rules, term orderings (KBO/LPO), redundancy criterion, completeness statement & P12 & S & must & named but undefined at the Part III resolution section (paramodulation/superposition); also A(c) \\
D-P11-004 & Resolution refinements: ordered/hyper/semantic resolution, set-of-support, subsumption & P12 & S & should &  \\
D-P11-005 & SMT: Nelson--Oppen combination and DPLL(T) & P12 & S & should &  \\
D-P11-006 & Logic programming: SLD resolution with soundness/completeness, negation as failure, Prolog & P9 & E & must & also A(c) -- coordinate with new-content register \\
D-P11-007 & Answer set programming and Datalog & P9 & S & should & also A(c) \\
D-P11-008 & Hoare logic rules table and weakest preconditions; separation logic pointer & P9 & E & must & also A(c) -- coordinate with new-content register \\
D-P11-009 & BDDs, QBF, \#SAT, MaxSAT & P12 & S & should & BDD absence also flagged in B (propositional) \\
D-P11-010 & Decidability and complexity results for every presented system: CPL (validity coNP-complete / SAT NP-complete), FOL (undecidable, c.e.), K/T/S4 and IPL PSPACE-complete, S5 NP-complete, LTL PSPACE-complete, QBF PSPACE-complete, plus decidable fragments & P14 & T & must & cross-cutting list; B3 automation: complexity beyond two facts absent and the coNP macro defined but never used; B3 temporal: LTL PSPACE-completeness lands here; master table harvested from entry-card fields -- depends on the P8 per-system entry cards \\
D-P11-011 & Prover and proof-assistant landscape table (Agda, HOL Light, HOL4, PVS, Mizar, ACL2, Metamath) & P12 & T & should & only Lean 4, Rocq, Isabelle named in v1.0 \\
D-P11-012 & MLTT/CIC foundations: universes, inductive types, W-types, definitional vs propositional equality, extraction & P12 & S & should & dependent types one passing mention in Part IX (B2 higher-order); also A(c) \\
D-P11-013 & HoTT and univalence pointer & P12 & P & could & NOT-IN-SCOPE beyond pointer per roadmap \\
D-P11-014 & LCF architecture and the de Bruijn criterion & P12 & S & should &  \\
D-P11-015 & Landmark formalizations table (Four Color, Odd Order, Flyspeck, CompCert, seL4) & P12 & T & could & cross-link to P13 chronology \\
\end{longtable}
\endgroup


\section{New-domain specification}
\label{sec:newdomains}

One register row per deliverable, keyed to the four domains of
decision~D1.  The \emph{sources} column seeds the bibliography intake
(§\ref{sec:apparatus}); rows with domain \textbf{x} are the
\textbf{NOT-IN-SCOPE register} --- exclusions recorded as decisions so
that absence is never again ambiguous between ``omitted deliberately''
and ``overlooked.''  The spec-first rule (§\ref{sec:campaign}) applies:
writing beyond a row's treatment ceiling requires editing the row first.

\begingroup\scriptsize\setlength{\tabcolsep}{3pt}
\begin{longtable}{@{}p{0.52in}p{0.85in}p{2.5in}p{0.22in}p{0.42in}p{1.95in}@{}}
\toprule
\textbf{ID} & \textbf{Target} & \textbf{Item} & \textbf{Tr} & \textbf{Prio} & \textbf{Sources} \\
\midrule
\endfirsthead
\toprule
\textbf{ID} & \textbf{Target} & \textbf{Item} & \textbf{Tr} & \textbf{Prio} & \textbf{Sources} \\
\midrule
\endhead
\midrule\multicolumn{6}{r}{\emph{continued}}\\\endfoot
\bottomrule\endlastfoot
N-a-001 & P3 3.5 countable-models & \(\omega\)-categoricity: Ryll-Nardzewski/Engeler/Svenonius characterization via types, DLO and random-graph examples & E & must & Ryll-Nardzewski 1959; Engeler 1959; Svenonius 1959; Hodges 1993 \\
N-a-002 & P3 3.5 countable-models & Fraisse constructions: hereditary/joint-embedding/amalgamation properties, existence and uniqueness of the Fraisse limit, random graph and extension axioms & E & must & Fraisse 1954; Rado 1964; Hodges 1993 \\
N-a-003 & P3 3.5 countable-models & prime, atomic, countably saturated, and homogeneous models: existence and uniqueness theorems over countable complete theories & E & must & Vaught 1961; Chang-Keisler 1990; Marker 2002 \\
N-a-004 & P3 3.5 countable-models & Vaught's never-two theorem (T); Vaught conjecture status report (P) & T & should & Vaught 1961; Morley 1970 \\
N-a-005 & P3 3.5 countable-models & two-cardinal theorems: Vaught's two-cardinal theorem statement and method sketch & S & could & Vaught 1965; Chang-Keisler 1990 \\
N-a-006 & P3 3.6 stability-classification & Morley rank and totally transcendental theories; outline of the proof of Morley's categoricity theorem; Baldwin-Lachlan refinement statement & S & should & Morley 1965; Baldwin-Lachlan 1971; Marker 2002 \\
N-a-007 & P3 3.6 stability-classification & stability spectrum and forking: stable/superstable/\(\omega\)-stable, nonforking extensions, symmetry and transitivity survey & S & should & Shelah 1978; Pillay 1996; Tent-Ziegler 2012 \\
N-a-008 & P3 3.6 stability-classification & indiscernible sequences and Ehrenfeucht-Mostowski models: existence via Ramsey's theorem, applications to non-isomorphic models & T & should & Ehrenfeucht-Mostowski 1956; Ramsey 1930; Hodges 1993 \\
N-a-009 & P3 3.7 tame-topology & o-minimality: definition, monotonicity and cell-decomposition statements, RCF as an o-minimal structure & S & should & Pillay-Steinhorn 1986; Knight-Pillay-Steinhorn 1986; van den Dries 1998 \\
N-a-010 & P3 3.7 tame-topology & NIP and simple theories: definitions, examples, map of the neostability landscape & S & could & Shelah 1971; Simon 2015; Tent-Ziegler 2012 \\
N-a-011 & P3 3.8 preservation-model-completeness & preservation theorems: Los-Tarski (universal), Lyndon (positivity), Chang-Los-Suszko (unions of chains) & T & should & Los 1955; Lyndon 1959; Chang-Keisler 1990; Hodges 1993 \\
N-a-012 & P3 3.8 preservation-model-completeness & model completeness and model companions: Robinson's test, ACF and RCF worked examples, existentially closed structures & E & must & Robinson 1956; Chang-Keisler 1990; Marker 2002 \\
N-a-013 & P3 3.4 interpolation-definability & Robinson joint consistency theorem and Beth definability with model-theoretic proofs; equivalence with Craig interpolation & T & should & Robinson 1956; Beth 1953; Hodges 1993 \\
N-a-014 & P3 3.2 types-saturation & ultrafilter existence from BPI, BPI versus AC status, Keisler-Shelah ultrapower isomorphism theorem statement & S & could & Keisler 1961; Shelah 1971; Bell-Slomson 1969 \\
N-a-015 & P7 7.4 lambda-calculi & untyped lambda calculus and combinatory logic as computation models: Church-Rosser theorem statement, lambda-definability equals partial recursiveness & T & must & Church 1936; Church-Rosser 1936; Kleene 1936; Curry-Feys 1958; Barendregt 1984 \\
N-a-016 & P4 4.4 degrees-and-priority & Turing degrees: partial order, join, jump operator, Kleene-Post construction of incomparable degrees below \(\mathbf 0'\) & E & must & Kleene-Post 1954; Rogers 1967; Soare 2016 \\
N-a-017 & P4 4.4 degrees-and-priority & Post's problem and the Friedberg-Muchnik theorem: finite-injury priority construction in full & T & must & Post 1944; Friedberg 1957; Muchnik 1956; Soare 2016 \\
N-a-018 & P4 4.4 degrees-and-priority & simple and creative sets, many-one and one-one reducibility, Myhill isomorphism theorem & T & should & Post 1944; Myhill 1955; Rogers 1967 \\
N-a-019 & P4 4.5 frontiers-of-computability & hyperarithmetical and analytical hierarchies: \(\Delta^1_1=\mathrm{HYP}\) statement, ordinal notations \(\mathcal O\) sketch & S & should & Kleene 1955; Rogers 1967; Sacks 1990 \\
N-a-020 & P4 4.5 frontiers-of-computability & Kolmogorov complexity: invariance theorem, incompressibility, Chaitin-style incompleteness & T & should & Kolmogorov 1965; Chaitin 1974; Li-Vitanyi 2008; Downey-Hirschfeldt 2010 \\
N-a-021 & P4 4.5 frontiers-of-computability & Hilbert's tenth problem: MRDP theorem statement, Diophantine sets equal c.e. sets, proof survey with Pell-equation core & S & must & Davis-Putnam-Robinson 1961; Matiyasevich 1970; Davis 1973 \\
N-a-022 & P4 4.5 frontiers-of-computability & undecidable word problems: Post-Markov theorem for semigroups (T), Novikov-Boone theorem for groups (S) & T & could & Post 1947; Markov 1947; Novikov 1955; Boone 1958 \\
N-a-023 & P4 4.5 frontiers-of-computability & busy beaver function: definition and Rado's noncomputability theorem & T & could & Rado 1962 \\
N-a-024 & P4 4.5 frontiers-of-computability & computable structure theory pointer: computable presentations, indices of structures & P & could & Ash-Knight 2000 \\
N-a-025 & P5 5.6 provability-logic & GL: axioms, Kripke semantics for transitive conversely-well-founded frames, completeness and FMP, fixed-point theorem, Solovay arithmetical completeness (S) & E & must & Solovay 1976; Sambin 1976; Segerberg 1971; Boolos 1993 \\
N-a-026 & P5 5.6 provability-logic & Rosser provability and witness comparison; Guaspari-Solovay \(R\) pointer & P & could & Guaspari-Solovay 1979; Lindstrom 2003 \\
N-a-027 & P5 5.5 independence-and-strength & Paris-Harrington theorem: statement, unprovability in PA via indicator/model-theoretic argument (S), provability in full-induction ACA and nonprovability in ACA_0 & S & must & Paris-Harrington 1977; Kanamori-McAloon 1987; Kaye 1991 \\
N-a-028 & P5 5.5 independence-and-strength & Goodstein sequences and the Kirby-Paris hydra: termination via \(\varepsilon_0\) (T), unprovability in PA statement & T & must & Goodstein 1944; Kirby-Paris 1982 \\
N-a-029 & P5 5.5 independence-and-strength & interpretability and relative interpretation: Tarski-Mostowski-Robinson method, essential undecidability transfer, interpretability logic pointer & S & should & Tarski-Mostowski-Robinson 1953; Visser 1997; Lindstrom 2003 \\
N-a-030 & P5 5.5 independence-and-strength & subsystems of second-order arithmetic and reverse mathematics: \(\mathsf{RCA}_0\), \(\mathsf{WKL}_0\), \(\mathsf{ACA}_0\), \(\mathsf{ATR}_0\), \(\Pi^1_1\text{-}\mathsf{CA}_0\); Big Five phenomenon with sample equivalences & S & must & Friedman 1975; Simpson 2009 \\
N-a-031 & P5 5.5 independence-and-strength & intensionality of consistency statements: Feferman's analysis, dependence on the provability predicate & S & should & Feferman 1960 \\
N-a-032 & P5 5.4 models-of-pa & nonstandard models of PA: existence and order type \(\mathbb N+\mathbb Z\cdot\mathbb Q\), overspill and its uses, Tennenbaum's theorem & E & must & Skolem 1934; Tennenbaum 1959; Kaye 1991 \\
N-a-033 & P5 5.7 decidable-arithmetic & Presburger arithmetic: quantifier elimination with congruence predicates, decidability, Fischer-Rabin lower bound statement & T & must & Presburger 1929; Fischer-Rabin 1974; Enderton 2001 \\
N-a-034 & P5 5.7 decidable-arithmetic & Skolem arithmetic of multiplication: decidability via Feferman-Vaught statement & S & could & Skolem 1930; Feferman-Vaught 1959 \\
N-a-035 & P6 6.5 cardinal-arithmetic & cofinality, regular and singular cardinals, Konig's theorem, Hausdorff recursion, Easton's theorem statement & E & must & Konig 1905; Easton 1970; Jech 2003; Kunen 2011 \\
N-a-036 & P6 6.5 cardinal-arithmetic & singular cardinal behavior: Silver's theorem statement, PCF pointer & S & could & Silver 1975; Shelah 1994 \\
N-a-037 & P6 6.4 infinitary-combinatorics & closed unbounded and stationary sets, Fodor's pressing-down lemma, diagonal intersection & E & must & Fodor 1956; Jech 2003; Kunen 2011 \\
N-a-038 & P6 6.4 infinitary-combinatorics & Delta-system lemma with proof and the ccc verification pattern for Cohen forcing & T & must & Shanin 1946; Kunen 2011 \\
N-a-039 & P6 6.4 infinitary-combinatorics & infinite Ramsey theorem (T); Erdos-Rado theorem and partition calculus notation (S) & T & must & Ramsey 1930; Erdos-Rado 1956; Jech 2003 \\
N-a-040 & P6 6.4 infinitary-combinatorics & trees: Konig infinity lemma, Aronszajn tree construction, Suslin hypothesis and its independence status survey & S & should & Konig 1927; Kurepa 1935; Todorcevic 1984; Kunen 2011 \\
N-a-041 & P6 6.4 infinitary-combinatorics & combinatorial principles: diamond and square statements with one application each; Martin's axiom with ccc applications & S & should & Jensen 1972; Martin-Solovay 1970; Kunen 2011 \\
N-a-042 & P6 6.3 constructibility-forcing & Cohen forcing of \(\neg\mathrm{CH}\): full ccc cardinal-preservation verification, nice names counting & T & must & Cohen 1963; Kunen 2011 \\
N-a-043 & P6 6.3 constructibility-forcing & independence of AC: symmetric extensions method survey; permutation models of ZFA & S & should & Fraenkel 1922; Cohen 1963; Jech 1973; Halbeisen 2017 \\
N-a-044 & P6 6.6 large-cardinals & large cardinals: inaccessible, Mahlo, weakly compact, measurable via ultrapowers and elementary embeddings, Scott's theorem \(V\neq L\), consistency-strength hierarchy chart to supercompact & S & must & Ulam 1930; Scott 1961; Keisler 1962; Kanamori 2003 \\
N-a-045 & P6 6.6 large-cardinals & Levy hierarchy, absoluteness for transitive models, Shoenfield absoluteness statement with applications & T & should & Levy 1965; Shoenfield 1961; Jech 2003 \\
N-a-046 & P6 6.7 descriptive-set-theory & descriptive set theory: Polish spaces, Borel hierarchy, analytic sets, Suslin's theorem, perfect-set property, regularity properties & E & must & Suslin 1917; Kechris 1995; Moschovakis 2009 \\
N-a-047 & P6 6.7 descriptive-set-theory & effective descriptive set theory bridge to computability: lightface pointclasses pointer & P & could & Moschovakis 2009; Sacks 1990 \\
N-a-048 & P6 6.8 determinacy & infinite games: Gale-Stewart theorem (T), Borel determinacy statement, AD versus AC, projective determinacy and Woodin cardinals pointer & S & should & Gale-Stewart 1953; Mycielski-Steinhaus 1962; Martin 1975; Martin-Steel 1989 \\
N-a-049 & P7 7.5 arithmetic-proof-theory & Gentzen's consistency proof for PA: reduction procedure on derivations, transfinite induction up to \(\varepsilon_0\), ordinal assignment in full for the propositional core & T & must & Gentzen 1936; Gentzen 1938; Takeuti 1987; Schwichtenberg 1977 \\
N-a-050 & P7 7.5 arithmetic-proof-theory & ordinal notation systems: Cantor normal form arithmetic, Veblen functions, \(\Gamma_0\) and Feferman-Schutte statement & S & should & Veblen 1908; Feferman 1964; Schutte 1977; Pohlers 2009 \\
N-a-051 & P7 7.5 arithmetic-proof-theory & PRA and finitism: Skolem's arithmetic, Tait's identification of finitism with PRA & S & should & Skolem 1923; Tait 1981 \\
N-a-052 & P7 7.5 arithmetic-proof-theory & omega-logic: the \(\omega\)-rule, cut elimination for PA with the \(\omega\)-rule, bound \(\varepsilon_0\) & S & could & Schutte 1951; Schwichtenberg 1977; Pohlers 2009 \\
N-a-053 & P7 7.6 functional-interpretations & Godel's Dialectica interpretation: System T, interpretation of HA, consistency corollary & S & should & Godel 1958; Troelstra 1973; Avigad-Feferman 1998 \\
N-a-054 & P7 7.6 functional-interpretations & Kleene number realizability: soundness for HA, ECT and Markov's principle independence applications & T & must & Kleene 1945; Troelstra 1998; van Oosten 2008 \\
N-a-055 & P7 7.6 functional-interpretations & disjunction and numerical existence properties for IPC and HA with proofs & T & must & Godel 1932; Gentzen 1935; Kleene 1945; Troelstra-van Dalen 1988 \\
N-a-056 & P7 7.7 structural-extensions & epsilon calculus: epsilon terms, first and second epsilon theorems statements, Herbrand's theorem in general form & S & should & Hilbert-Bernays 1939; Herbrand 1930; Buss 1998 \\
N-a-057 & P7 7.7 structural-extensions & generalized sequent frameworks survey: hypersequents, labelled calculi, display calculi, deep inference, with one worked modal example each & S & should & Avron 1996; Negri 2005; Belnap 1982; Guglielmi 2007; Negri-von Plato 2001 \\
N-a-058 & P7 7.7 structural-extensions & proof mining pointer: extraction of bounds from classical proofs & P & could & Kohlenbach 2008 \\
N-a-059 & P3 3.9 abstract-model-theory & infinitary logic \(L_{\omega_1\omega}\): syntax, Scott isomorphism theorem, Scott sentences, Karp completeness statement, compactness failure & S & should & Karp 1964; Scott 1965; Keisler 1971; Marker 2016 \\
N-a-060 & P3 3.9 abstract-model-theory & generalized quantifiers: Mostowski and Lindstrom quantifiers, completeness for \(L(Q_1)\) statement & S & could & Mostowski 1957; Lindstrom 1966; Keisler 1970 \\
N-a-061 & P3 3.9 abstract-model-theory & abstract logics and the full proof of Lindstrom's first theorem & T & should & Lindstrom 1969; Barwise 1974; Ebbinghaus-Flum-Thomas 1994 \\
N-a-062 & P3 3.9 abstract-model-theory & independence-friendly and dependence logic pointer & P & could & Hintikka-Sandu 1989; Vaananen 2007 \\
N-a-063 & P4 4.1 machines-and-models & \(\mu\)-recursion and partial recursive functions: normal form and s-m-n recap, equivalence theorem Turing-computable \(\equiv\) \(\mu\)-recursive & E & must & Kleene 1936; Turing 1936; Kleene 1952; Odifreddi 1989 \\
N-a-064 & P4 4.6 complexity & complexity theory as a subject: class ladder P, NP, coNP, PH, PSPACE, EXPTIME, time and space hierarchy theorems, logic-relevant completeness results table & T & must & Hartmanis-Stearns 1965; Cook 1971; Karp 1972; Stockmeyer 1976; Papadimitriou 1994; Arora-Barak 2009 \\
N-a-065 & P3 3.1 elementary-maps-ef-games & Ehrenfeucht-Fraisse games: game characterization of elementary equivalence, quantifier rank, inexpressibility applications & E & must & Fraisse 1954; Ehrenfeucht 1961; Hodges 1993; Ebbinghaus-Flum-Thomas 1994 \\
N-a-066 & P6 6.1 ordinal-arithmetic & ordinal arithmetic developed as theory: sums, products, exponentiation, Cantor normal form, transfinite induction and recursion theorems & E & must & Cantor 1897; Sierpinski 1965; Jech 2003; Kunen 2011 \\
N-a-067 & P6 6.3 constructibility-forcing & constructible hierarchy \(L\): condensation lemma statement, \(\mathrm{GCH}\) in \(L\), reflection principle, fine-structure pointer & S & should & Godel 1938; Levy 1960; Devlin 1984; Jech 2003 \\
N-a-068 & P6 6.3 constructibility-forcing & Boolean-valued models at statement level; product and iterated forcing pointers & S & should & Scott-Solovay 1967; Bell 2005; Kunen 2011 \\
N-a-069 & P7 7.2 cut-elimination & midsequent theorem via cut elimination and Herbrand's theorem as corollary & S & should & Gentzen 1935; Herbrand 1930; Takeuti 1987; Buss 1998 \\
N-a-070 & P7 7.2 cut-elimination & Craig interpolation proved by Maehara's method with Beth definability corollary & E & must & Craig 1957; Maehara 1960; Takeuti 1987 \\
N-a-071 & P2 2.3 second-higher-order & second-order and HOL proof calculi displayed: comprehension and choice schemas as a calculus, Henkin completeness route & E & must & Church 1940; Henkin 1950; Andrews 2002; Shapiro 1991 \\
N-b-001 & P8 8.2 modal & canonical models and completeness for K, T, D, B, K4, S4, S5: Lindenbaum construction for normal modal logics, canonicity discussion & E & must & Lemmon-Scott 1977; Chellas 1980; Blackburn-de Rijke-Venema 2001 \\
N-b-002 & P8 8.2 modal & frame correspondence: definability of the T/4/5/B/D axioms, Sahlqvist theorem statement with algorithmic examples (proof S) & T & must & Sahlqvist 1975; van Benthem 1983; Blackburn-de Rijke-Venema 2001 \\
N-b-003 & P8 8.2 modal & bisimulation: definition, invariance, van Benthem characterization theorem & T & must & van Benthem 1976; Blackburn-de Rijke-Venema 2001 \\
N-b-004 & P8 8.2 modal & filtration, finite model property, and decidability for K, T, S4, S5 & T & must & Segerberg 1971; Chagrov-Zakharyaschev 1997; Blackburn-de Rijke-Venema 2001 \\
N-b-005 & P8 8.2 modal & Kripke frame incompleteness phenomena: general frames, Thomason and Fine counterexamples survey & S & should & Thomason 1974; Fine 1974; Chagrov-Zakharyaschev 1997 \\
N-b-006 & P8 8.2 modal & non-normal modal logics and neighborhood semantics: systems E, M, C, N, regular logics, completeness & E & should & Scott 1970; Montague 1970; Chellas 1980; Pacuit 2017 \\
N-b-007 & P8 8.2 modal & per-system entry cards for the modal cube: K, D, T, B, K4, S4, S4.2, S5, GL, Grz -- axioms, frame class, FMP, decidability, characteristic theorems & E & must & Chellas 1980; Hughes-Cresswell 1996; Chagrov-Zakharyaschev 1997 \\
N-b-008 & P8 8.2 modal & two-dimensional semantics: actuality operator, fixedly, Kaplan character/content distinction & S & should & Crossley-Humberstone 1977; Davies-Humberstone 1980; Kaplan 1989 \\
N-b-009 & P8 8.2 modal & quantified modal logic: constant versus varying domains, Barcan and converse Barcan formulas, rigid designators, existence predicate, actualism/possibilism & E & must & Barcan 1946; Kripke 1963; Fitting-Mendelsohn 1998; Garson 2001; Williamson 2013 \\
N-b-010 & P8 8.4 conditionals-deontic-free & free logic: positive/negative/neutral systems, inner and outer domains, definite descriptions (Russell versus free-logic treatments) & E & should & Russell 1905; Lambert 1967; Bencivenga 1986; Lehmann 2002 \\
N-b-011 & P8 8.3 temporal-dynamic-epistemic & multi-agent epistemic logic: S5\(_n\)/KD45\(_n\), common knowledge as fixed point, muddy children analysis & E & must & Hintikka 1962; Aumann 1976; Fagin-Halpern-Moses-Vardi 1995 \\
N-b-012 & P8 8.3 temporal-dynamic-epistemic & dynamic epistemic logic: public announcement logic axioms, reduction laws, unsuccessful updates & S & should & Plaza 1989; Baltag-Moss-Solecki 1998; van Ditmarsch-van der Hoek-Kooi 2008 \\
N-b-013 & P8 8.4 conditionals-deontic-free & deontic logic: SDL as KD, Ross and Good Samaritan and Chisholm contrary-to-duty paradoxes, dyadic conditional obligation & E & must & von Wright 1951; Chisholm 1963; Hansson 1969; McNamara 2006 \\
N-b-014 & P8 8.3 temporal-dynamic-epistemic & philosophical tense logic: Prior's systems, since/until and Kamp's expressive completeness statement, Ockhamist versus Peircean branching time & S & should & Prior 1967; Kamp 1968; Burgess 1984 \\
N-b-015 & P8 8.3 temporal-dynamic-epistemic & hybrid logic: nominals, satisfaction operators, axiomatization survey & S & could & Prior 1967; Blackburn 2000; Areces-ten Cate 2007 \\
N-b-016 & P8 8.4 conditionals-deontic-free & counterfactuals: Stalnaker selection functions versus Lewis sphere systems, conditional excluded middle debate, limit assumption, completeness statements for C1 and VC & E & must & Stalnaker 1968; Lewis 1973; Nute 1980 \\
N-b-017 & P8 8.4 conditionals-deontic-free & indicative conditionals: paradoxes of material implication, Adams's thesis, Lewis triviality results, suppositional theories & S & must & Adams 1975; Lewis 1976; Jackson 1987; Edgington 1995; Bennett 2003 \\
N-b-018 & P8 8.4 conditionals-deontic-free & conditional logics as a family: CK and neighbors, selection-function and premise semantics, correspondence table & T & should & Chellas 1975; Nute 1980 \\
N-b-019 & P11 11.2 truth-theories & Tarskian truth as a philosophical unit: convention T, object/metalanguage discipline, hierarchy of languages & S & should & Tarski 1935; Tarski 1944 \\
N-b-020 & P11 11.2 truth-theories & Kripke's fixed-point theory of truth: partial Kleene valuation schemes, least and maximal fixed points, groundedness, ghost of the Tarski hierarchy & T & must & Kripke 1975; Martin 1984; Halbach 2011 \\
N-b-021 & P11 11.2 truth-theories & revision theory of truth survey & S & could & Herzberger 1982; Gupta-Belnap 1993 \\
N-b-022 & P11 11.2 truth-theories & axiomatic theories of truth: TB/CT conservativity, KF, FS, comparative survey & S & should & Friedman-Sheard 1987; Feferman 1991; Halbach 2011; Horsten 2011 \\
N-b-023 & P11 11.1 paradoxes & paradox catalog with formal diagnoses: strengthened liar, Curry, Yablo, Grelling, Berry & E & must & Curry 1942; Yablo 1993; Priest 2006; Beall-Glanzberg-Ripley 2018 \\
N-b-024 & P11 11.4 vagueness & vagueness: sorites structure, supervaluationism with global versus local validity, degree-theoretic approaches, epistemicism, contextualism pointer & S & must & Fine 1975; Williamson 1994; Keefe 2000; Smith 2008 \\
N-b-025 & P8 8.5 many-valued-fuzzy & first-degree entailment: Belnap-Dunn four-valued semantics, bilattice reading, relation to K3 and LP & E & must & Dunn 1976; Belnap 1977; Omori-Wansing 2017 \\
N-b-026 & P8 8.7 paraconsistent-nonmonotonic & dialetheism and the choice of paraconsistent consequence: motivation from paradox, inclosure schema, replies & S & should & Priest 1979; Priest 2006 \\
N-b-027 & P8 8.6 substructural-linear & relevance logic landscape: E and R axiomatic systems, variable-sharing proof, ternary-frame conditions per axiom, admissibility of gamma & S & should & Anderson-Belnap 1975; Meyer-Dunn 1969; Routley-Meyer 1973; Dunn-Restall 2002; Mares 2004 \\
N-b-028 & P8 8.5 many-valued-fuzzy & Lukasiewicz logics: L3, finite Ln, infinite-valued L, axiomatization, McNaughton theorem statement, real-valued semantics & E & must & Lukasiewicz 1920; Lukasiewicz-Tarski 1930; McNaughton 1951; Cignoli-D'Ottaviano-Mundici 2000 \\
N-b-029 & P8 8.5 many-valued-fuzzy & t-norm fuzzy logics: continuous t-norms, BL and MTL, Godel and product logics, standard completeness survey & S & must & Hajek 1998; Esteva-Godo 2001; Cintula-Hajek-Noguera 2011 \\
N-b-030 & P8 8.1 intuitionistic-intermediate & intuitionism beyond the calculus: BHK reading, Brouwerian counterexamples, choice sequences, bar induction, continuity principles, Markov's principle status & S & should & Brouwer 1908; Heyting 1956; Dummett 1977; Troelstra-van Dalen 1988 \\
N-b-031 & P8 8.1 intuitionistic-intermediate & intermediate logics: Godel-Dummett LC, KC, the lattice of superintuitionistic logics survey & S & should & Godel 1932; Dummett 1959; Jankov 1968; Chagrov-Zakharyaschev 1997 \\
N-b-032 & P8 8.7 paraconsistent-nonmonotonic & nonmonotonic consequence: preferential models, KLM system P, rational monotony; AGM belief revision postulates, Levi and Harper identities & E & must & Alchourron-Gardenfors-Makinson 1985; Kraus-Lehmann-Magidor 1990; Lehmann-Magidor 1992; Makinson 2005 \\
N-b-033 & P8 8.7 paraconsistent-nonmonotonic & default logic and circumscription entry cards & S & should & Reiter 1980; McCarthy 1980 \\
N-b-034 & P11 11.3 consequence-and-constants & philosophy of logical consequence: Tarski's analysis, model-theoretic versus proof-theoretic accounts, demarcation of logical constants (invariance), harmony & S & should & Tarski 1936; Prawitz 1974; Etchemendy 1990; Sher 1991; Dummett 1991 \\
N-b-035 & P11 11.5 ontological-commitment & plural quantification and the status of second-order logic: Boolos's plural reading, ontological commitment, Quine's 'set theory in disguise' charge & S & should & Quine 1970; Boolos 1984; Boolos 1985; Shapiro 1991 \\
N-c-001 & P9 9.1 finite-model-theory & finite model theory foundations: failure of compactness and completeness in the finite, Craig/Beth failure, EF-game method on finite structures, Trakhtenbrot recap & E & must & Ebbinghaus-Flum 1995; Libkin 2004 \\
N-c-002 & P9 9.1 finite-model-theory & zero-one laws for first-order logic: Glebskii and Fagin theorems, extension axioms, the random-graph link & T & must & Glebskii et al. 1969; Fagin 1976; Libkin 2004 \\
N-c-003 & P9 9.1 finite-model-theory & locality: Hanf and Gaifman theorems with inexpressibility applications (connectivity, acyclicity) & T & should & Hanf 1965; Gaifman 1982; Libkin 2004 \\
N-c-004 & P9 9.1 finite-model-theory & descriptive complexity: Fagin's theorem (existential SO captures NP) with proof sketch, Immerman-Vardi (LFP captures P on ordered structures) statement & T & must & Fagin 1974; Vardi 1982; Immerman 1986; Immerman 1999 \\
N-c-005 & P9 9.2 decidable-fragments & the classical decision problem: prefix-class classification, Bernays-Schonfinkel-Ramsey, Ackermann, Godel-Kalmar-Schutte classes, reduction classes & S & should & Lewis 1979; Borger-Gradel-Gurevich 1997 \\
N-c-006 & P9 9.2 decidable-fragments & two-variable fragment: finite model property and NEXPTIME decidability statement & S & should & Mortimer 1975; Gradel-Kolaitis-Vardi 1997 \\
N-c-007 & P9 9.2 decidable-fragments & guarded fragment: decidability and the modal connection & S & should & Andreka-van Benthem-Nemeti 1998; Gradel 1999 \\
N-c-008 & P9 9.2 decidable-fragments & monadic first-order logic: Lowenheim-Behmann decidability & T & should & Lowenheim 1915; Behmann 1922; Borger-Gradel-Gurevich 1997 \\
N-c-009 & P9 9.3 description-logics & ALC: syntax, model semantics, TBox/ABox, tableau decision procedure with blocking, ExpTime complexity statement & E & must & Schmidt-Schauss-Smolka 1991; Baader et al. 2003; Baader-Horrocks-Lutz-Sattler 2017 \\
N-c-010 & P9 9.3 description-logics & lightweight and expressive DLs: EL family with polynomial classification, DL-Lite and query answering, SHOIQ/OWL profiles map & S & should & Baader-Brandt-Lutz 2005; Calvanese et al. 2007; Baader-Horrocks-Lutz-Sattler 2017 \\
N-c-011 & P9 9.3 description-logics & ALC as notational variant of multimodal K (Schild correspondence) & T & should & Schild 1991 \\
N-c-012 & P9 9.4 temporal-logics-automata & LTL deepened: past operators, Kamp and Gabbay-Pnueli-Shelah-Stavi expressive completeness statements, safety/liveness, PSPACE satisfiability & T & must & Kamp 1968; Pnueli 1977; Gabbay-Pnueli-Shelah-Stavi 1980; Sistla-Clarke 1985 \\
N-c-013 & P9 9.4 temporal-logics-automata & omega-automata: Buchi/Rabin/parity acceptance, closure properties, complementation statement, Vardi-Wolper LTL-to-automata translation & S & must & Buchi 1962; McNaughton 1966; Vardi-Wolper 1986; Safra 1988; Thomas 1997 \\
N-c-014 & P9 9.4 temporal-logics-automata & MSO and automata: Buchi-Elgot-Trakhtenbrot theorem on finite words (T), S1S decidability, Rabin's S2S statement & T & should & Buchi 1960; Elgot 1961; Trakhtenbrot 1961; Rabin 1969; Thomas 1997 \\
N-c-015 & P9 9.4 temporal-logics-automata & tree automata pointer and applications to decidability transfer & P & could & Thatcher-Wright 1968; Comon et al. 2007 \\
N-c-016 & P9 9.5 model-checking & CTL and CTL*: syntax, semantics, fixpoint characterizations of EG/EU, labeling algorithm with complexity, LTL versus CTL expressiveness & E & must & Clarke-Emerson 1981; Queille-Sifakis 1982; Emerson-Halpern 1986; Clarke-Grumberg-Peled 1999; Baier-Katoen 2008 \\
N-c-017 & P9 9.5 model-checking & symbolic model checking: BDDs and Bryant canonicity, fixpoint computation, bounded model checking via SAT, CEGAR pointer & S & must & Bryant 1986; Burch et al. 1992; McMillan 1993; Biere-Cimatti-Clarke-Zhu 1999; Clarke et al. 2003 \\
N-c-018 & P9 9.5 model-checking & modal mu-calculus: syntax, Knaster-Tarski semantics, relation to parity games, alternation hierarchy statement & S & should & Kozen 1983; Emerson-Jutla 1991; Bradfield-Stirling 2007 \\
N-c-019 & P9 9.6 program-logics & Hoare logic: partial and total correctness, axioms and rules, soundness proof, relative completeness (Cook), weakest preconditions (Dijkstra) & E & must & Floyd 1967; Hoare 1969; Dijkstra 1975; Cook 1978; Apt 1981 \\
N-c-020 & P9 9.6 program-logics & separation logic: heap assertions, frame rule, symbolic execution, bi-abduction pointer & S & should & O'Hearn-Reynolds-Yang 2001; Reynolds 2002; Calcagno et al. 2011; O'Hearn 2019 \\
N-c-021 & P9 9.6 program-logics & propositional dynamic logic: syntax, regular programs, Segerberg axioms with Kozen-Parikh completeness statement, Fischer-Ladner closure, ExpTime decidability & E & should & Pratt 1976; Fischer-Ladner 1979; Kozen-Parikh 1981; Harel-Kozen-Tiuryn 2000 \\
N-c-022 & P9 9.7 logic-programming-databases & logic programming: definite programs, least Herbrand model, immediate-consequence fixpoint, SLD-resolution soundness and completeness, negation as failure with Clark completion, SLDNF & E & must & Kowalski 1974; van Emden-Kowalski 1976; Clark 1978; Apt-van Emden 1982; Lloyd 1987 \\
N-c-023 & P9 9.7 logic-programming-databases & answer set programming: stable-model semantics via Gelfond-Lifschitz reduct, complexity statement & S & should & Gelfond-Lifschitz 1988; Marek-Truszczynski 1999; Brewka-Eiter-Truszczynski 2011 \\
N-c-024 & P9 9.7 logic-programming-databases & Datalog and conjunctive queries: Chandra-Merlin homomorphism theorem, containment decidability, data/combined complexity & S & should & Chandra-Merlin 1977; Vardi 1982; Abiteboul-Hull-Vianu 1995 \\
N-c-025 & P9 9.7 logic-programming-databases & constraint satisfaction and logic pointer: CSP as homomorphism problem, dichotomy theorem statement & P & could & Feder-Vardi 1998; Bulatov 2017; Zhuk 2017 \\
N-c-026 & P9 9.8 complexity-of-logic & complexity of logical tasks: Cook-Levin recap, QBF PSPACE-completeness with proof, Ladner-style table for modal logics (K/S4 PSPACE, S5 NP), FO validity undecidable recap & T & must & Cook 1971; Stockmeyer-Meyer 1973; Ladner 1977; Papadimitriou 1994 \\
N-c-027 & P7 7.8 proof-complexity & proof complexity: Cook-Reckhow framework, resolution lower bound for the pigeonhole principle (Haken) proof sketch, Frege systems open-problems map & S & should & Cook-Reckhow 1979; Haken 1985; Beame-Pitassi 1998; Krajicek 2019 \\
N-c-028 & P9 9.8 complexity-of-logic & interpolation in verification pointer: McMillan-style interpolant-based model checking & P & could & McMillan 2003 \\
N-c-029 & P7 7.4 lambda-calculi & System F: syntax, impredicative encodings, Girard's strong normalization statement, Curry-Howard for second-order, parametricity pointer & S & must & Girard 1972; Reynolds 1974; Reynolds 1983; Girard-Lafont-Taylor 1989 \\
N-c-030 & P7 7.4 lambda-calculi & Martin-Lof type theory: judgments, Pi/Sigma/identity types, universes, propositions-as-types with proof relevance, normalization and decidable type checking statements & S & must & Martin-Lof 1984; Nordstrom-Petersson-Smith 1990; Hofmann 1997 \\
N-c-031 & P12 12.5 interactive-proof & lambda cube and pure type systems; calculus of (inductive) constructions as the Coq/Lean foundation & S & should & Coquand-Huet 1988; Barendregt 1992; Paulin-Mohring 2015 \\
N-c-032 & P12 12.1 sat & CDCL: clause learning with UIP, watched literals, restarts, VSIDS; GRASP/Chaff/MiniSat lineage; DRAT proofs and certified UNSAT & E & must & Marques-Silva-Sakallah 1999; Moskewicz et al. 2001; Een-Sorensson 2003; Marques-Silva-Lynce-Malik 2009; Wetzler-Heule-Hunt 2014 \\
N-c-033 & P12 12.1 sat & tractable cases and hardness structure: 2-SAT via implication graphs, Horn-SAT unit propagation, random k-SAT phase transition pointer & T & should & Aspvall-Plass-Tarjan 1979; Dowling-Gallier 1984; Mitchell-Selman-Levesque 1992 \\
N-c-034 & P12 12.2 smt & SMT internals: DPLL(T) architecture, congruence closure for EUF, simplex for linear arithmetic, Nelson-Oppen combination with proof for convex stably-infinite theories & E & must & Nelson-Oppen 1979; Dutertre-de Moura 2006; Nieuwenhuis-Oliveras-Tinelli 2006; de Moura-Bjorner 2008; Barrett-Sebastiani-Seshia-Tinelli 2009 \\
N-c-035 & P12 12.3 first-order-atp & unification theory deepened: Martelli-Montanari rule-based algorithm with termination proof, mgu theorem, occurs-check, equational and AC-unification survey, undecidability of higher-order unification & E & must & Robinson 1965; Huet 1973; Goldfarb 1981; Martelli-Montanari 1982; Baader-Snyder 2001 \\
N-c-036 & P12 12.3 first-order-atp & ordered resolution and the superposition calculus: reduction orders, redundancy and saturation, model-construction completeness proof sketch & S & must & Bachmair-Ganzinger 1994; Nieuwenhuis-Rubio 2001 \\
N-c-037 & P12 12.4 rewriting & term rewriting: abstract reduction systems, Newman's lemma, critical pairs and Knuth-Bendix completion, termination orders (LPO/KBO), word-problem applications & E & must & Newman 1942; Knuth-Bendix 1970; Huet 1980; Dershowitz 1982; Baader-Nipkow 1998 \\
N-c-038 & P12 12.5 interactive-proof & LCF architecture and kernels: de Bruijn criterion, proof terms versus proof scripts, trusted computing bases of Coq/Lean/Isabelle/HOL, reflection & S & should & de Bruijn 1970; Gordon-Milner-Wadsworth 1979; Pollack 1998; Harrison 2009 \\
N-c-039 & P12 12.5 interactive-proof & program extraction from constructive proofs: realizability-based extraction & S & could & Paulin-Mohring 1989; Letouzey 2003 \\
N-c-040 & P12 12.5 interactive-proof & ATP-ITP integration pointer: hammers and proof reconstruction & P & could & Paulson-Blanchette 2010; Blanchette et al. 2016 \\
N-d-001 & P10 10.1 universal-algebra & universal algebra toolkit: algebras, homomorphisms, congruences and quotients, subdirect representation, Birkhoff HSP variety theorem & T & must & Birkhoff 1935; Birkhoff 1944; Burris-Sankappanavar 1981 \\
N-d-002 & P10 10.1 universal-algebra & equational logic: term algebras, free algebras, Birkhoff completeness for equational consequence & T & must & Birkhoff 1935; Burris-Sankappanavar 1981 \\
N-d-003 & P10 10.2 boolean-stone & Boolean algebras: filters, ideals, ultrafilters, atomic versus atomless, BPI status relative to AC & E & must & Halmos 1963; Bell-Slomson 1969; Koppelberg 1989 \\
N-d-004 & P10 10.2 boolean-stone & Stone representation and duality: proof of the representation theorem, Stone spaces, clopen-algebra inverse, uniqueness of the countable atomless BA & T & must & Stone 1936; Halmos 1963; Johnstone 1982; Givant-Halmos 2009 \\
N-d-005 & P10 10.2 boolean-stone & Lindenbaum-Tarski algebras: algebraic completeness proof for CPL, free Boolean algebras & E & must & Tarski 1935; Rasiowa-Sikorski 1963 \\
N-d-006 & P10 10.2 boolean-stone & Priestley duality for distributive lattices & S & should & Priestley 1970; Davey-Priestley 2002 \\
N-d-007 & P10 10.2 boolean-stone & Stone-duality dividends for logic: compactness as compactness of the Stone space, type spaces of theories as Stone spaces & S & should & Johnstone 1982; Marker 2002 \\
N-d-008 & P10 10.3 heyting-algebras & Heyting algebras: definition, opens of a topological space, algebraic completeness of IPC, Rasiowa-Sikorski lemma & T & must & Heyting 1930; McKinsey-Tarski 1946; Rasiowa-Sikorski 1963; Balbes-Dwinger 1974 \\
N-d-009 & P10 10.3 heyting-algebras & Esakia duality and the lattice of intermediate logics survey & S & should & Esakia 1974; Chagrov-Zakharyaschev 1997 \\
N-d-010 & P10 10.3 heyting-algebras & locales and pointfree topology pointer & P & could & Johnstone 1982; Picado-Pultr 2012 \\
N-d-011 & P10 10.4 modal-algebras & Boolean algebras with operators: Jonsson-Tarski representation, canonical extensions, duality with general frames, algebraic completeness of normal modal logics & T & must & Jonsson-Tarski 1951; Goldblatt 1989; Blackburn-de Rijke-Venema 2001; Venema 2007 \\
N-d-012 & P10 10.4 modal-algebras & topological semantics: McKinsey-Tarski completeness of S4 for topological interior, completeness for the real line statement & T & should & McKinsey-Tarski 1944; Rasiowa-Sikorski 1963; van Benthem-Bezhanishvili 2007 \\
N-d-013 & P10 10.5 residuated-mv & MV-algebras and Chang completeness for infinite-valued Lukasiewicz logic; Mundici's l-group correspondence statement & T & must & Chang 1958; Chang 1959; Mundici 1986; Cignoli-D'Ottaviano-Mundici 2000 \\
N-d-014 & P10 10.5 residuated-mv & residuated lattices and FL-algebras: algebraization of substructural logics, subvariety lattice survey, cut elimination correspondences & S & must & Ono-Komori 1985; Galatos-Jipsen-Kowalski-Ono 2007 \\
N-d-015 & P10 10.5 residuated-mv & BL- and MTL-algebras and standard completeness statements & S & should & Hajek 1998; Cignoli-Esteva-Godo-Torrens 2000; Jenei-Montagna 2002 \\
N-d-016 & P10 10.6 algebraized-fol & cylindric algebras: algebraization of first-order logic, representability problems survey & S & should & Henkin-Monk-Tarski 1971; Henkin-Monk-Tarski 1985; Monk 2000 \\
N-d-017 & P10 10.6 algebraized-fol & relation algebras: Tarski's calculus of relations, representable relation algebras, undecidability statement, set theory without variables pointer & S & should & Tarski 1941; Tarski-Givant 1987; Maddux 2006 \\
N-d-018 & P10 10.6 algebraized-fol & Blok-Pigozzi algebraizability pointer & P & could & Blok-Pigozzi 1989; Font 2016 \\
N-d-019 & P10 10.7 category-toolkit & category theory for logicians: categories, functors, natural transformations, adjunctions, limits and colimits, Yoneda lemma statement & E & must & Mac Lane 1998; Awodey 2010 \\
N-d-020 & P10 10.7 category-toolkit & Lawvere algebraic theories and functorial semantics & S & should & Lawvere 1963; Adamek-Rosicky-Vitale 2011 \\
N-d-021 & P10 10.8 ccc-lambda & cartesian closed categories and typed lambda calculus: Curry-Howard-Lambek three-way correspondence with the equivalence theorem & T & must & Lambek 1980; Lambek-Scott 1986 \\
N-d-022 & P10 10.9 topos-overview & quantifiers as adjoints; hyperdoctrine view of first-order logic overview & S & should & Lawvere 1969; Pitts 2000 \\
N-d-023 & P10 10.9 topos-overview & elementary topos overview: subobject classifier, internal logic is intuitionistic HOL, presheaf examples, Kripke-Joyal semantics & S & must & Lawvere 1970; Lambek-Scott 1986; Mac Lane-Moerdijk 1992; Johnstone 2002 \\
N-d-024 & P10 10.9 topos-overview & topos-theoretic independence pointer: Tierney's sheaf-theoretic forcing and CH & P & could & Tierney 1972; Mac Lane-Moerdijk 1992 \\
N-x-001 & P7 -- & homotopy type theory and univalent foundations beyond a one-page pointer & P & -- & Univalent Foundations Program 2013 \\
N-x-002 & P7 -- & higher inductive types and cubical type theory internals & P & -- & Cohen-Coquand-Huber-Mortberg 2016 \\
N-x-003 & P6 -- & proper forcing, iterated forcing, and the forcing-axiom landscape (PFA/MM) beyond a pointer & P & -- & Shelah 1982; Foreman-Magidor-Shelah 1988 \\
N-x-004 & P6 -- & inner model theory, fine structure, and core models & P & -- & Jensen 1972; Mitchell-Steel 1994 \\
N-x-005 & P6 -- & PCF theory and advanced cardinal arithmetic beyond Silver's theorem statement & P & -- & Shelah 1994 \\
N-x-006 & P8 -- & quantum logic (orthomodular lattices as logic) & P & -- & Birkhoff-von Neumann 1936 \\
N-x-007 & P8 -- & adaptive logics & P & -- & Batens 2007 \\
N-x-008 & P11 -- & inductive logic and confirmation theory & P & -- & Carnap 1950 \\
N-x-009 & P8 -- & formal argumentation frameworks & P & -- & Dung 1995 \\
N-x-010 & P10 -- & abstract algebraic logic beyond the Blok-Pigozzi pointer & P & -- & Font-Jansana-Pigozzi 2003; Font 2016 \\
N-x-011 & P7 -- & geometry of interaction & P & -- & Girard 1989 \\
N-x-012 & P7 -- & proof nets beyond a pointer (correctness criteria, MLL net theory) & P & -- & Girard 1987; Danos-Regnier 1989 \\
N-x-013 & P10 -- & categorical logic beyond the CCC and topos overview (fibrations, sketches, infinity-topoi) & P & -- & Jacobs 1999; Lurie 2009 \\
N-x-014 & P4 -- & computable analysis and computability over the reals (TTE, Weihrauch degrees) & P & -- & Weihrauch 2000; Pour-El-Richards 1989 \\
N-x-015 & P9 -- & probabilistic and stochastic model checking, probability logics beyond a pointer & P & -- & Halpern 2003; Baier-Katoen 2008 \\
\end{longtable}
\endgroup


\section{Target handbook architecture}
\label{sec:architecture}

\subsection{Design rules}
\begin{enumerate}
  \item The handbook has \textbf{14 parts plus a proof annex}; the body
        is lookup material (entries, tables, statements with status and
        citation) and the annex holds compressed proofs, organized
        parallel to the body.
  \item Per-chapter proof sections and a separate companion volume are
        both rejected: the former re-interleaves proofs with lookup
        material; the latter is premature (the annex can be split out
        later without rework).
  \item The v1.0 manuscript is \textbf{not printed}.  After Phase-2
        repairs it lives at \texttt{source/} as the verification
        substrate; its exercises, feedback chapter, and qualifying
        problems are retired from print but retained as verification
        assets, and solutions that currently do theorems' work are
        promoted into annex proofs during migration.
  \item Every numbered body item carries a status field with an address:
        \emph{proved} (Annex~P$n$.$x$) / \emph{sketched}
        (Annex~\ldots) / \emph{cited} (source).  Compression may only
        \emph{downgrade} status; upgrading requires re-verification.
\end{enumerate}

\subsection{The fourteen parts}
\begin{description}[style=nextline]
  \item[Part 1 --- Notation, Symbols, and Conventions.]
    Master symbol registry (generated from
    \texttt{registry/symbols.csv}); binding and precedence for
    \emph{all} systems; notation variants across the literature; the
    16 binary truth functions; mathematical preliminaries as
    definition/theorem tables.
  \item[Part 2 --- Classical Logic.]
    2.1 propositional (semantics, equivalences, normal forms incl.\
    Tseitin, functional completeness with Post's lattice, duality,
    calculi side by side: Hilbert / ND / LK / tableaux / resolution,
    metatheorems, complexity); 2.2 first-order (syntax, semantics,
    calculi, metatheorem catalog incl.\ the countable/choice-free
    completeness case, prenex--Skolem--clausal pipeline,
    decidable-fragments table, equality metatheory); 2.3 second- and
    higher-order (full vs.\ Henkin, a displayed SOL calculus, simple
    type theory); 2.4 variants (many-sorted, free, infinitary,
    generalized quantifiers).
  \item[Part 3 --- Model Theory.]
    Elementary maps and diagrams; Ehrenfeucht--Fra\"iss\'e games;
    types, saturation, prime/atomic models; ultraproducts and
    \L o\'s; preservation theorems; the QE catalog per theory;
    categoricity and Morley; stability/NIP/o-minimality overview
    tables; finite model theory; Lindstr\"om.
  \item[Part 4 --- Computability and Complexity.]
    Machine models and $\mu$-recursion with the equivalence theorem;
    Church--Turing; s-m-n, recursion, Rice; undecidability catalog;
    c.e.-set structure; degrees, jump, Post's problem, priority
    method; arithmetical and analytical hierarchies; Kolmogorov
    complexity; complexity-class ladder and logic-relevant
    completeness results; descriptive complexity.
  \item[Part 5 --- Arithmetic and Incompleteness.]
    Q and PA; representability and coding (repaired); the diagonal
    lemma; the G\"odel~I and~II, Rosser, L\"ob, and Tarski entries;
    provability logic GL with Solovay completeness; decidable
    arithmetics; nonstandard models and Tennenbaum; independence
    catalog (Paris--Harrington, Goodstein); interpretability;
    second-order arithmetic and the reverse-mathematics calibration
    table.
  \item[Part 6 --- Set Theory.]
    ZFC axiom table and system variants; ordinal arithmetic and
    Cantor normal form; cardinal arithmetic, cofinality, K\"onig; AC
    equivalents; clubs/stationary/Fodor, $\Delta$-system, Ramsey;
    combinatorial principles ($\Diamond$, $\square$, MA); $V$, $L$,
    HOD; absoluteness and reflection; forcing entries with the honest
    $\neg$CH treatment; independence catalog; the large-cardinal
    chart; descriptive set theory and determinacy.
  \item[Part 7 --- Proof Theory and $\lambda$-Calculi.]
    Calculi catalog with identity cards; structural-rules table; cut
    elimination with mix displayed; subformula property, midsequent,
    Herbrand; interpolation (Maehara); normalization and strong
    normalization; Curry--Howard correspondence table; untyped
    $\lambda$-calculus and Church--Rosser; System F; MLTT/CIC
    overview; ordinal-analysis calibration; Gentzen's consistency
    proof (survey, honestly labeled); proof complexity.
  \item[Part 8 --- Nonclassical Logics.]
    The flagship part: uniform per-system \textbf{entry cards}
    (name/aliases, language, semantics, axioms/rules with a pointer to
    the rule sheet, adequacy theorems with status, characteristic
    theorems and failures, decidability and complexity, translations,
    algebraic counterpart, sources).  8.1 intuitionistic and
    intermediate; 8.2 modal (systems lattice, correspondence table,
    canonical models, filtration, decidability, quantified modal
    logic, GL); 8.3 temporal, dynamic, epistemic (LTL, CTL/CTL*,
    $\mu$-calculus, PDL, S5$_n$, common knowledge with its induction
    axiom, public announcement); 8.4 conditionals, deontic, free
    logic; 8.5 many-valued and fuzzy; 8.6 substructural (family map,
    relevance, full classical linear logic, Lambek); 8.7
    paraconsistent and nonmonotonic (with the AGM postulates table).
  \item[Part 9 --- Logic in Computer Science.]
    Description-logic family table with complexity; logic programming
    (SLD resolution, negation as failure, Prolog, ASP, Datalog);
    Hoare logic and weakest preconditions, separation-logic pointer;
    model checking (automata-theoretic LTL, fixpoint CTL, BDDs);
    database and finite-model connections; type systems.
  \item[Part 10 --- Algebraic and Categorical Logic.]
    Boolean algebras and Stone duality; Lindenbaum--Tarski algebras;
    Heyting, MV, and modal algebras with J\'onsson--Tarski;
    residuated lattices; cylindric/relation-algebra pointers;
    categorical semantics (CCC $\leftrightarrow$ STLC, hyperdoctrine
    pointer, topos overview).
  \item[Part 11 --- Philosophy of Logic.]
    Paradox catalog with resolution strategies; theories of truth
    (Tarski hierarchy, Kripke fixed points, revision); consequence
    relations; vagueness; identity; ontological commitment; schools
    capsules.
  \item[Part 12 --- Automation and Mechanization.]
    SAT: DPLL and CDCL, 2-SAT/Horn; resolution refinements;
    unification (Martelli--Montanari algorithm and theorem);
    superposition as a calculus; SMT (theories table, Nelson--Oppen,
    DPLL(T)); tableau and connection methods; QBF/\#SAT/MaxSAT;
    prover and proof-assistant landscape; kernel foundations (LCF,
    de Bruijn criterion, MLTT/CIC in kernels); landmark
    formalizations; typesetting and reproducibility reference.
  \item[Part 13 --- Chronology and Biography.]
    Chronology table; landmark-papers table; $\sim$40 capsule
    biographies; notation-history table.
  \item[Part 14 --- Apparatus.]
    Complete rule sheets ($\sim$20, generated from the calculus
    registry); master cross-reference tables (decidability and
    complexity per logic; consistency-strength ladder; expressive
    power; axiom$\,\leftrightarrow\,$frame$\,\leftrightarrow\,$algebra
    correspondence) --- \emph{harvested} from entry-card fields, not
    authored; glossary; per-part further reading; subject, symbol,
    person, and named-theorem indexes.
  \item[Appendix P --- Proof Annex.]
    Parallel to Parts 2--12.  Tier 1: method entries (Henkin,
    diagonalization, \L o\'s, back-and-forth, canonical models,
    filtration, cut elimination, normalization, forcing, EF, priority)
    --- compressed proof plus a body method box.  Tier 2: results the
    manuscript proves completely --- compressed to 40--60\% of source
    length, \texttt{proofboundary} markers retained.  Tier 3: cited
    theorems --- statement only, no annex entry.
\end{description}

\subsection{Source-to-target mapping}
\begin{center}\small
\begin{tabular}{@{}llp{2.9in}@{}}
\toprule
\textbf{v1.0 part (lines)} & \textbf{Target} & \textbf{Disposition} \\
\midrule
I Foundations (363--1546) & P1, P6 & compress $\sim$4:1 to registry tables; proof-writing chapter's core into front matter \\
II Propositional (1547--2695) & P2.1 & entries + tables; proofs to annex \\
III First-order (2696--4994) & P2.2, P2.4 & entries + tables; Henkin to annex Tier 1 \\
IV Model theory (4995--6203) & P3 & entries; \L o\'s, back-and-forth to annex Tier 1 \\
V Computability (6204--7190) & P4 & entries + hierarchy tables \\
VI Arithmetic (7191--8068) & P5 & entries; audit table institutionalized \\
VII Set theory (8069--9134) & P6 & entries; forcing to annex Tier 1 (repaired) \\
VIII Proof theory (9135--9948) & P7 & entries; cut elimination to annex Tier 1 \\
IX Second/higher-order (9949--10621) & P2.3 & entries; many-sorted translation to annex Tier 2 \\
X Nonclassical (10622--11683) & P8 & exploded into entry cards \\
XI Practice/automation (11684--12675) & P12 & compressed + heavy additions \\
Back matter (12676--14780) & P14 / retired & registers and sheets rebuilt from data; exercises retired to verification assets \\
\bottomrule
\end{tabular}
\end{center}

\subsection{Templates}
\label{sec:templates}

The entry-card and theorem environments carry \emph{mandatory} keys;
compilation fails without them (ruling~\ref{ruling:status}).  Sketch of
the planned interface:

\begin{lstlisting}
\begin{systementry}{name=S5, aliases={KT5},
    language=modal-prop, sheet=sheet:S5}
  \entrysemantics{Kripke frames with equivalence relations}
  \entryadequacy{sound+strongly complete for local consequence,
    status=proved, annex=P8.2.4}
  \entrydecidability{decidable; SAT is NP-complete,
    status=cited, source=Ladner1977}
  \entrytranslations{into FOL via standard translation}
  \entryalgebra{monadic algebras}
  \entrysources{BlackburnDeRijkeVenema2001}
\end{systementry}

\begin{hbtheorem}{title=Compactness, status=proved,
    annex=P2.2.7, sources={Godel1930,Malcev1936}}
  ...statement only; the proof lives at the annex address...
\end{hbtheorem}
\end{lstlisting}

Label taxonomy for the new tree: \texttt{entry:}, \texttt{thm:},
\texttt{tab:}, \texttt{fig:}, \texttt{sheet:}, \texttt{annex:},
\texttt{sec:} --- replacing the v1.0 label population of which
$\sim$85\% is never referenced.

\section{Apparatus program}
\label{sec:apparatus}

\paragraph{Indexes.}
Four indexes, hand-authored at definition sites with inverted headings
and sort keys (\texttt{Godel@G\"odel}); \emph{see}/\emph{see also}
entries; target 2{,}500--4{,}000 subject entries.  The v1.0 auto-index
wrapper (theorem titles as entries) is retired.

\paragraph{Registries as data.}
\texttt{registry/symbols.csv} (glyph, macro, role, scope, precedence
rank, first use) generates the Part-1 tables and drives the notation
lint; \texttt{registry/calculi/*.tex} holds every inference rule ---
body chapters, entry cards, and Part-14 rule sheets all \verb|\input|
from it, making body--sheet divergence structurally impossible.

\paragraph{Bibliography.}
Manual \texttt{thebibliography} (64 entries) converts to biblatex in
Phase~3; the five missing primaries (Glivenko 1929, G\"odel 1933,
Gentzen 1936, Kleene 1938, Priest 1979) enter immediately; the
per-domain source lists in \texttt{new-content.csv} are the intake
queue.  Projected final size 250--350 entries.

\paragraph{The build list.}
\begingroup\scriptsize\setlength{\tabcolsep}{3pt}
\begin{longtable}{@{}p{0.45in}p{0.78in}p{2.02in}p{2.6in}p{0.42in}p{0.16in}@{}}
\toprule
\textbf{ID} & \textbf{Category} & \textbf{Item} & \textbf{Specification} & \textbf{Prio} & \textbf{Ph} \\
\midrule
\endfirsthead
\toprule
\textbf{ID} & \textbf{Category} & \textbf{Item} & \textbf{Specification} & \textbf{Prio} & \textbf{Ph} \\
\midrule
\endhead
\midrule\multicolumn{6}{r}{\emph{continued}}\\\endfoot
\bottomrule\endlastfoot
A-001 & rule-sheet & Rule sheet: ND-CPL & All twelve natural-deduction rules plus the repetition rule and the vacuous-discharge convention, displayed, named, and labeled, generated from the calculus-registry file. & must & 6 \\
A-002 & rule-sheet & Rule sheet: ND-FOL & Quantifier and equality rules with the generalization-variable and eigenvariable side conditions stated on the sheet itself. & must & 6 \\
A-003 & rule-sheet & Rule sheet: unified LK & Canonical \(\mathsf{LK}\) (primitive negation, contexted initial sequents) with the Part VIII presentation documented as a named variant and an equivalence-theorem cross-reference. & must & 6 \\
A-004 & rule-sheet & Rule sheet: LK with equality & \(\mathsf{LK}^{=}\) as a separately named calculus displaying the equality rules and stating that cut elimination for it is a distinct theorem with its own status. & must & 6 \\
A-005 & rule-sheet & Rule sheet: LJ & Single-succedent intuitionistic sequent calculus, load-bearing for the disjunction-property proof, with the restriction stated explicitly. & must & 6 \\
A-006 & rule-sheet & Rule sheet: Hilbert CPL & Axiom schemas and modus ponens with the deduction theorem stated as a named result on the sheet. & must & 6 \\
A-007 & rule-sheet & Rule sheet: Hilbert FOL & Axiom schemas, generalization rule, and the deduction theorem with its free-variable side condition. & must & 6 \\
A-008 & rule-sheet & Rule sheet: Hilbert K & Axiom \(\mathsf K\), necessitation, and uniform substitution as the base normal modal system, with the substitution terminology matching the registered term. & must & 6 \\
A-009 & rule-sheet & Rule sheet: T & System \(\mathsf T\) as \(\mathsf K\) plus the reflexivity axiom, with its frame condition cross-referenced to the correspondence table. & must & 6 \\
A-010 & rule-sheet & Rule sheet: S4 & System \(\mathsf{S4}\) as \(\mathsf T\) plus axiom 4, with its frame condition cross-referenced to the correspondence table. & must & 6 \\
A-011 & rule-sheet & Rule sheet: S5 & System \(\mathsf{S5}\) with axiom 5 (or B plus 4), its equivalence-relation frame condition, and the NP satisfiability note. & must & 6 \\
A-012 & rule-sheet & Rule sheet: GL & L\"ob axiom system with necessitation, the arithmetical interpretation, and the Solovay completeness statement cross-referenced. & must & 6 \\
A-013 & rule-sheet & Rule sheet: LTL & Axioms including \(\mathsf U\)-unfolding and the induction rule, with operator precedence licensed by the completed precedence table. & must & 6 \\
A-014 & rule-sheet & Rule sheet: S5n epistemic & Multi-agent \(\mathsf{S5}_n\) with \(K_i\), \(E_G\), and \(C_G\) axioms, the \(C_G\) induction rule, and the common-knowledge weak-completeness caveat. & must & 6 \\
A-015 & rule-sheet & Rule sheet: classical propositional tableaux & Expansion and closure rules with the systematic-tableau completeness statement and countermodel read-off. & must & 6 \\
A-016 & rule-sheet & Rule sheet: classical first-order tableaux & Gamma and delta rules with freshness side conditions displayed on the sheet. & must & 6 \\
A-017 & rule-sheet & Rule sheet: many-valued matrix tableau & Value-labelled expansion rules with the notion of a saturated branch defined on the sheet, curing the undefined-term defect. & must & 6 \\
A-018 & rule-sheet & Rule sheet: resolution, factoring, superposition & Rules with the soundness-critical side conditions (renaming apart, most general unifier) displayed, plus the term ordering (KBO/LPO) and redundancy criteria for superposition. & must & 6 \\
A-019 & rule-sheet & Rule sheet: IMLL/ILL/CLL & Sequent rules for intuitionistic and classical linear logic including the four exponential rules, each named, with the one-sided classical presentation. & must & 6 \\
A-020 & rule-sheet & Rule sheet: Lambek calculus & Directional implication introduction and elimination rules in sequent form with the no-structural-rules statement. & should & 6 \\
A-021 & rule-sheet & Rule sheet: STLC / Curry-Howard & Typing rules annotated with proof terms beside the corresponding natural-deduction rules. & must & 6 \\
A-022 & rule-sheet & Rule sheet: Hoare logic & Axiom schemas and inference rules including the consequence rule, cross-referenced to the weakest-precondition table. & must & 6 \\
A-023 & rule-sheet & Rule sheet: SLD resolution & Definite-clause resolution with the selection rule stated and the negation-as-failure caveat. & should & 6 \\
A-024 & table & Table: the sixteen binary connectives & All sixteen binary truth functions with standard names, truth tables, and definability notes including Sheffer stroke and NOR. & must & 4 \\
A-025 & table & Table: notation variants across the literature & Peano/Russell, Polish, Hilbert, and Gentzen notations for connectives and quantifiers keyed to the handbook's registered symbols. & should & 4 \\
A-026 & table & Table: Greek, Fraktur, and blackboard letter conventions & Registered uses of Greek, Fraktur, and blackboard alphabets by domain with scope column matching the symbol registry. & should & 4 \\
A-027 & table & Table: standard propositional equivalence laws & De Morgan, distribution, absorption, double negation, material implication, and contraposition laws with labels citable from body chapters. & must & 4 \\
A-028 & table & Table: duality & Dual connectives and the duality theorem for propositional formulas with the negation-normal-form connection. & should & 4 \\
A-029 & table & Table: Post functional-completeness criteria & The five maximal clones with membership tests and the completeness criterion for a set of connectives, supporting the Post-lattice figure. & must & 4 \\
A-030 & table & Table: tractable propositional classes & Horn, dual-Horn, 2-SAT, and XOR-SAT with their satisfiability algorithms and complexity. & should & 4 \\
A-031 & table & Table: propositional proof systems side-by-side & Hilbert, ND, LK, tableau, and resolution for CPL compared on rules, adequacy theorems, and proof-search behavior. & must & 4 \\
A-032 & table & Table: natural-language translation guide & Canonical translations of English connective idioms (only, unless, neither, only if) correcting the v1.0 translation-table errors. & should & 4 \\
A-033 & table & Table: quantifier equivalence laws & Prenexing, distribution, vacuous-quantifier, and duality equivalences with side conditions on free variables. & must & 4 \\
A-034 & table & Table: normal-form comparison & NNF, prenex, Skolem, CNF, DNF, and clausal form compared on existence, equivalence versus equisatisfiability, and size blowup. & must & 4 \\
A-035 & table & Table: first-order metatheorem summary & Soundness, completeness, compactness, both L\"owenheim--Skolem theorems, Herbrand, interpolation, Beth, and Lindstr\"om with exact hypotheses including choice-principle usage. & must & 4 \\
A-036 & table & Table: decidable fragments of first-order logic & Monadic FOL, \(\mathrm{FO}^2\), guarded fragment, and Bernays--Sch\"onfinkel class with decision complexity and undecidable boundary cases. & must & 4 \\
A-037 & table & Table: first-order proof systems side-by-side & Hilbert, ND, LK, and tableaux for FOL compared on rules, side conditions, and adequacy theorems. & should & 4 \\
A-038 & table & Table: extensions of first-order logic & Many-sorted, free, infinitary, generalized-quantifier, and second-order extensions with the Lindstr\"om-theorem calibration of which metatheorems survive. & should & 4 \\
A-039 & table & Table: full versus Henkin second-order semantics & Comparison of full and Henkin semantics on completeness, compactness, L\"owenheim--Skolem, and categoricity, resolving the D2 gap. & must & 4 \\
A-040 & table & Table: first-order theories master table & DLO, ACF, RCF, Presburger, PA, true arithmetic, and ZFC classified by completeness, decidability, quantifier elimination, and categoricity spectrum. & must & 4 \\
A-041 & table & Table: quantifier-elimination catalog & Per-theory QE status with eliminating extensions (definable Skolem functions, model completeness) and the DLO/ACF/RCF witnesses. & must & 4 \\
A-042 & table & Table: preservation theorems & {\L}o\'s--Tarski, Chang--{\L}o\'s--Suszko, and Lyndon theorems pairing syntactic classes with preservation properties. & must & 4 \\
A-043 & table & Table: stability, NIP, and o-minimality overview & Classification-theoretic dividing lines with defining conditions and canonical example theories. & should & 4 \\
A-044 & table & Table: finite model theory results & Zero-one laws, Fagin's theorem, and Trakhtenbrot's theorem with the finite-versus-unrestricted contrast. & should & 4 \\
A-045 & table & Table: machine models and equivalences & Turing machines, register machines, mu-recursive functions, and lambda-definability with the pairwise equivalence theorems and the Church--Turing discussion pointer. & must & 4 \\
A-046 & table & Table: undecidability catalog & Halting problem, Rice, Hilbert's tenth, word problems, tiling, and FOL validity with reduction sources and citations. & must & 4 \\
A-047 & table & Table: arithmetical hierarchy & \(\Sigma^0_n\)/\(\Pi^0_n\)/\(\Delta^0_n\) levels with closure properties, complete problems, Post's theorem, and the jump correspondence. & must & 4 \\
A-048 & table & Table: analytical hierarchy & \(\Sigma^1_n\)/\(\Pi^1_n\) levels with complete sets and the relation to second-order arithmetic. & must & 4 \\
A-049 & table & Table: Turing degrees and jump facts & Degree-structure facts (jump, join, Post's problem, Friedberg--Muchnik) at statement level with citations. & should & 4 \\
A-050 & table & Table: complexity-class ladder & P, NP, coNP, PH, PSPACE, EXPTIME, and beyond with known inclusions, separations, open problems, and canonical complete problems. & must & 4 \\
A-051 & table & Table: logic decision-problem complexity (master) & For every logic in the handbook, satisfiability/validity complexity or undecidability with completeness degree, harvested from entry-card fields. & must & 6 \\
A-052 & table & Table: descriptive-complexity correspondences & FO, existential SO (Fagin), least-fixed-point, and their capture results over ordered finite structures. & should & 4 \\
A-053 & table & Table: Q and PA axioms & Robinson arithmetic and PA axiom schemas in the fixed signature, with the induction schema's side conditions. & must & 4 \\
A-054 & table & Table: arithmetic independence catalog & G\"odel and Rosser sentences, \(\mathsf{Con}\) statements, Paris--Harrington, and Goodstein with the theory each is independent of. & must & 4 \\
A-055 & table & Table: decidable arithmetic theories & Presburger and Skolem arithmetic with decision procedures, complexity bounds, and the contrast with full arithmetic. & should & 4 \\
A-056 & table & Table: reverse-mathematics calibration & The Big Five subsystems of \(\mathsf Z_2\) with axioms and representative equivalent theorems. & must & 4 \\
A-057 & table & Table: provability-logic dictionary & GL modal principles paired with their arithmetical readings and the Solovay arithmetical-completeness statement. & must & 4 \\
A-058 & table & Table: ZFC axioms & Formal statement of every ZFC axiom and schema with side conditions, replacing the prose-only presentation. & must & 4 \\
A-059 & table & Table: set-theory system variants & ZF, Z, KP, NBG, MK, and NF compared on classes, foundation, choice, and relative strength. & must & 4 \\
A-060 & table & Table: ordinal arithmetic laws & Addition, multiplication, and exponentiation laws with continuity and commutativity failures and Cantor normal form. & must & 4 \\
A-061 & table & Table: cardinal arithmetic laws & Sum, product, and exponentiation rules, cofinality, regular/singular classification, K\"onig's theorem, and the \(\aleph\)/\(\beth\) hierarchies. & must & 4 \\
A-062 & table & Table: AC equivalents & Well-ordering theorem, Zorn's lemma, Tychonoff, and further equivalents with ZF-provability status of each implication. & must & 4 \\
A-063 & table & Table: choice-principle dependency & Which handbook theorems require AC, DC, countable choice, or BPI, including the choice-free countable-language completeness theorem. & must & 4 \\
A-064 & table & Table: combinatorial principles & Diamond, square, MA, club and stationary sets, Fodor, Delta-system lemma, and Ramsey's theorem with implication structure. & must & 4 \\
A-065 & table & Table: large-cardinal ladder & Inaccessible through supercompact in implication order with defining properties and first failures of \(V=L\); stronger hypotheses remain an explicit pointer. & must & 4 \\
A-066 & table & Table: independence catalog & Statements independent of ZFC (CH, GCH patterns, Suslin hypothesis, AC over ZF) with the method (forcing, inner models) and sources. & must & 4 \\
A-067 & table & Table: forcing-poset catalog & Cohen, random, and collapse forcings with what each adds and preserves and its chain conditions. & should & 4 \\
A-068 & table & Table: Borel and projective hierarchies & Pointclass levels with closure properties, regularity properties, and the determinacy calibration (Borel determinacy in ZFC; PD from large cardinals; AD versus AC). & should & 4 \\
A-069 & table & Table: consistency-strength ladder (master) & PRA through PA, subsystems of \(\mathsf Z_2\), Z, ZF, ZFC, and large-cardinal extensions in one ordered ladder harvested from entry cards, paired with the ladder figure. & must & 6 \\
A-070 & table & Table: structural rules per calculus & Weakening, contraction, exchange, and cut across LK, LJ, linear, relevance, affine, and Lambek systems, keying the substructural family map. & must & 4 \\
A-071 & table & Table: calculi catalog & LK, LJ, NK, NJ with hypersequent, display, and deep-inference pointers, each with its adequacy status. & should & 4 \\
A-072 & table & Table: Curry-Howard correspondence & Logic to type theory to category theory dictionary (formulas/types/objects, proofs/terms/morphisms, normalization/reduction). & must & 4 \\
A-073 & table & Table: lambda-calculi comparison & Untyped, simply typed, System F, and MLTT/CIC compared on Church--Rosser, normalization, and typability decidability. & must & 4 \\
A-074 & table & Table: ordinal-analysis calibration & Proof-theoretic ordinals for Q-range fragments, PA (\(\varepsilon_0\)), and stronger systems at statement level with citations. & must & 6 \\
A-075 & table & Table: proof-complexity systems & Resolution, Frege, and extended Frege with known lower bounds (pigeonhole in resolution) and simulation relations. & should & 4 \\
A-076 & table & Table: modal systems and axioms & K, D, T, B, K4, KD45, S4, S4.2, S4.3, S5, GL, and Grz with axioms and containments, backing the lattice figure. & must & 4 \\
A-077 & table & Table: frame correspondence (axiom, condition, algebra) & Modal axioms paired with first-order frame conditions and algebraic conditions, with Sahlqvist coverage noted, harvested from entry cards. & must & 6 \\
A-078 & table & Table: Barcan and domain variation & Barcan and converse Barcan formulas against constant, increasing, and varying domains with validity status in each regime. & must & 4 \\
A-079 & table & Table: temporal logics comparison & LTL, CTL, CTL*, and the modal \(\mu\)-calculus compared on expressiveness and satisfiability and model-checking complexity. & must & 4 \\
A-080 & table & Table: epistemic systems & \(\mathsf{S5}_n\), group knowledge \(E_G\)/\(C_G\)/\(D_G\), and public-announcement axioms with the \(C_G\) induction rule included. & must & 4 \\
A-081 & table & Table: intermediate logics & KC and LC (G\"odel--Dummett) with axioms, frame classes, disjunction property, and finite-model-property status. & should & 4 \\
A-082 & table & Table: negative translations & G\"odel--Gentzen, Glivenko, Kolmogorov, and Kuroda translations with exact clauses (settling the \(\bot\)-atomic case) and preservation theorems. & must & 4 \\
A-083 & table & Table: modal embeddings and standard translation & The FO standard translation, the G\"odel--McKinsey--Tarski embedding of IPL into S4, and the GL arithmetical interpretation as one translation catalog. & should & 4 \\
A-084 & table & Table: many-valued matrices comparison & K3, LP, weak Kleene, FDE, and finite {\L}ukasiewicz matrices with designated values, tautology sets, and the K3/LP duality registered. & must & 4 \\
A-085 & table & Table: t-norm fuzzy logics & BL, {\L}ukasiewicz, G\"odel, and product logics with t-norms, axiomatizations, and MV-algebra completeness statements. & should & 4 \\
A-086 & table & Table: substructural family map & Which structural rules are dropped to obtain relevance, linear, affine, and Lambek systems, paired with the family-map figure. & must & 4 \\
A-087 & table & Table: relevance systems & R, E, and RM axiom systems with the variable-sharing property and its status per system. & should & 4 \\
A-088 & table & Table: linear-logic connectives & Multiplicative, additive, and exponential connectives with units, polarities, and De Morgan duals in one registry-aligned table. & must & 4 \\
A-089 & table & Table: conditional and counterfactual systems & Stalnaker and Lewis sphere-semantics systems with characteristic validities and failures (antecedent strengthening, transitivity). & must & 4 \\
A-090 & table & Table: deontic systems & Standard deontic logic with its axioms and the classical paradoxes (Ross, gentle murder) at statement level. & should & 4 \\
A-091 & table & Table: free-logic variants & Positive, negative, and neutral free logics with existence predicate, quantifier rules, and definite-description treatment. & should & 4 \\
A-092 & table & Table: nonmonotonic formalisms comparison & Default logic, circumscription, autoepistemic logic, and ASP compared on consequence definition and complexity. & should & 4 \\
A-093 & table & Table: AGM postulates & Contraction and revision postulates with the Levi and Harper identities. & must & 4 \\
A-094 & table & Table: paraconsistent systems & LP and da Costa C-systems with explosion failure and designated-value or consistency-operator mechanisms. & should & 4 \\
A-095 & table & Table: description-logic family complexity & ALC through SROIQ with OWL correspondences and reasoning-problem complexity per logic. & must & 4 \\
A-096 & table & Table: Hoare logic rules & All Hoare rules with side conditions plus the weakest-precondition calculus and a separation-logic pointer. & must & 4 \\
A-097 & table & Table: logic-programming semantics & SLD resolution, negation as failure, completion, well-founded, and stable-model semantics compared, with Datalog noted. & should & 4 \\
A-098 & table & Table: Boolean algebra laws & Axiomatizations of Boolean algebras with the lattice-theoretic and ring-theoretic presentations. & must & 4 \\
A-099 & table & Table: Stone duality dictionary & The algebra-to-space correspondences (algebras/Stone spaces, homomorphisms/continuous maps, ideals/open sets), paired with the duality figure. & must & 4 \\
A-100 & table & Table: logic-algebra correspondences & CPL/Boolean, IPL/Heyting, S4/interior, {\L}ukasiewicz/MV, and substructural/residuated-lattice pairs via the Lindenbaum--Tarski construction. & must & 4 \\
A-101 & table & Table: paradox catalog & Liar, Russell, Curry, Berry, Yablo, sorites, and Skolem paradoxes with formulation and principal resolution strategies. & must & 4 \\
A-102 & table & Table: truth-theories comparison & Tarski hierarchy, Kripke fixed points, and revision theory compared on object-language truth predicates and paradox handling. & must & 4 \\
A-103 & table & Table: consequence-relation properties & Tarskian closure conditions (reflexivity, monotonicity, cut) and which substructural and nonmonotonic systems drop each. & should & 4 \\
A-104 & table & Table: vagueness approaches & Supervaluationist, degree-theoretic, and epistemicist treatments of the sorites at capsule level. & could & 4 \\
A-105 & table & Table: SAT-solving techniques & DPLL versus CDCL with clause learning, watched literals, restarts, and DRAT proof logging. & must & 4 \\
A-106 & table & Table: unification algorithm & Martelli--Montanari transformation rules with the occurs check and the mgu existence/uniqueness theorem. & must & 4 \\
A-107 & table & Table: SMT theory catalog & EUF, arrays, bit-vectors, and arithmetic theories with decidability, complexity, and Nelson--Oppen combinability conditions. & must & 4 \\
A-108 & table & Table: prover and proof-assistant landscape & Lean, Rocq, Isabelle, Agda, HOL Light/HOL4, PVS, Mizar, and ACL2 classified by foundation, kernel architecture, and de Bruijn criterion. & must & 4 \\
A-109 & table & Table: chronology of logic & Dated milestone table from Aristotle to the present keyed to handbook entries. & should & 6 \\
A-110 & table & Table: landmark papers & The landmark-papers register with full citations and one-line significance statements linked to the bibliography. & should & 6 \\
A-111 & table & Table: capsule biographies & Short biographical capsules for the figures in the person index. & could & 6 \\
A-112 & table & Table: expressive-power hierarchy (master) & Propositional, modal (via bisimulation invariance), FO, LFP, and SO ordered by expressive power with witnessing separations. & must & 6 \\
A-113 & table & Table: per-part audit registers & One generated audit table per part on the model of the v1.0 incompleteness audit table (line 8005), emitted by the audit-table generator, never hand-authored. & must & 6 \\
A-114 & figure & Figure: formula syntax tree & Parse tree of a worked formula supporting the unique-readability material and the syntax-tree exercises. & must & 4 \\
A-115 & figure & Figure: ND derivation tree example & A complete natural-deduction derivation displayed as a tree with discharge labels. & must & 4 \\
A-116 & figure & Figure: LK derivation tree example & A cut-free LK derivation displayed as a tree with rule names at each inference. & must & 4 \\
A-117 & figure & Figure: sequent-versus-ND correspondence & Side-by-side translation of one derivation between ND and LK showing the correspondence of rules. & should & 4 \\
A-118 & figure & Figure: normal-form conversion pipeline & Flowchart from formula through NNF, prenex, Skolem, and clausal form marking equivalence-preserving versus equisatisfiability-preserving steps. & should & 4 \\
A-119 & figure & Figure: Post's lattice & The lattice of clones with the five maximal clones highlighted, supporting the functional-completeness criterion table. & must & 4 \\
A-120 & figure & Figure: closed propositional tableau example & A closed classical tableau tree for a worked unsatisfiable formula. & must & 4 \\
A-121 & figure & Figure: resolution refutation example & A resolution refutation displayed as a DAG from input clauses to the empty clause. & should & 4 \\
A-122 & figure & Figure: Ehrenfeucht--Fra\"iss\'e game & A game board for two small structures illustrating a winning duplicator strategy over a fixed number of rounds. & should & 4 \\
A-123 & figure & Figure: bisimulation zigzag & The zigzag diagram for bisimulation between two Kripke models supporting the modal-invariance material. & should & 4 \\
A-124 & figure & Figure: arithmetical hierarchy & Inclusion diagram of \(\Sigma^0_n\)/\(\Pi^0_n\)/\(\Delta^0_n\) with strictness arrows and complete-set labels. & must & 4 \\
A-125 & figure & Figure: analytical hierarchy & Inclusion diagram of \(\Sigma^1_n\)/\(\Pi^1_n\) atop the arithmetical hierarchy. & should & 4 \\
A-126 & figure & Figure: Turing degrees and the jump & Schematic of the degree ordering with \(\mathbf 0\), \(\mathbf 0'\), incomparable c.e. degrees, and the jump operation. & should & 4 \\
A-127 & figure & Figure: complexity-class inclusion diagram & Hasse-style diagram of P through EXPTIME with known strict inclusions marked, paired with the ladder table. & must & 4 \\
A-128 & figure & Figure: reverse-mathematics Big Five & Inclusion diagram of the five subsystems of \(\mathsf Z_2\) with representative theorems attached. & should & 4 \\
A-129 & figure & Figure: ordinals to epsilon-zero & Number-line schematic of ordinals through \(\omega\), \(\omega^\omega\), and \(\varepsilon_0\) with Cantor-normal-form annotations. & should & 4 \\
A-130 & figure & Figure: cumulative hierarchy & The \(V_\alpha\) picture with rank axis and the position of small large cardinals indicated. & must & 4 \\
A-131 & figure & Figure: constructible hierarchy & The \(L_\alpha\) hierarchy inside \(V\) illustrating condensation and the \(V=L\) contrast. & should & 4 \\
A-132 & figure & Figure: forcing extension & Ground model, generic filter, and extension \(M[G]\) with names and their evaluations schematically indicated. & must & 4 \\
A-133 & figure & Figure: choice-principle implication diagram & AC, DC, countable choice, and BPI with known non-reversals marked, paired with the dependency table. & should & 4 \\
A-134 & figure & Figure: consistency-strength ladder & Vertical ladder from PRA to large cardinals matching the master consistency-strength table. & must & 6 \\
A-135 & figure & Figure: Borel and projective hierarchies & Pointclass diagram with the Borel levels and projective classes and their closure arrows. & should & 4 \\
A-136 & figure & Figure: modal-systems lattice & Containment lattice of K, D, T, B, K4, S4, S4.2, S4.3, S5, GL, and Grz, paired with the axiom table. & must & 4 \\
A-137 & figure & Figure: Kripke countermodel for excluded middle & Two-world intuitionistic model refuting \(p\lor\neg p\) with forcing annotations at each world. & must & 4 \\
A-138 & figure & Figure: Kripke countermodel for Peirce's law & Intuitionistic countermodel to \(((p\to q)\to p)\to p\) with forcing annotations. & must & 4 \\
A-139 & figure & Figure: Kripke countermodel refuting T & Non-reflexive frame refuting \(\Box p\to p\) demonstrating the correspondence direction. & must & 4 \\
A-140 & figure & Figure: Kripke countermodel refuting 4 & Non-transitive frame refuting \(\Box p\to\Box\Box p\). & must & 4 \\
A-141 & figure & Figure: Kripke countermodel refuting 5 & Non-Euclidean frame refuting \(\Diamond p\to\Box\Diamond p\). & must & 4 \\
A-142 & figure & Figure: Kripke countermodel for the Barcan formula & Increasing-domain model refuting the Barcan formula, paired with the domain-variation table. & must & 4 \\
A-143 & figure & Figure: Routley--Meyer ternary relation & Diagram of the ternary accessibility relation evaluating a relevant conditional at a world. & should & 4 \\
A-144 & figure & Figure: lattice of intermediate logics & IPL to CPL interval with KC and LC placed, matching the intermediate-logics table. & could & 4 \\
A-145 & figure & Figure: LTL evaluation on a path & A timeline diagram evaluating \(\mathsf X\), \(\mathsf U\), \(\mathsf F\), and \(\mathsf G\) formulas at positions of one path. & must & 4 \\
A-146 & figure & Figure: CTL computation tree & Branching-time tree contrasting path quantifiers with the linear LTL picture. & should & 4 \\
A-147 & figure & Figure: B\"uchi automaton example & A B\"uchi automaton for a worked LTL formula with its acceptance condition, supporting the model-checking entry. & should & 4 \\
A-148 & figure & Figure: SLD tree & An SLD derivation tree for a worked definite program and goal showing success and failure branches. & should & 4 \\
A-149 & figure & Figure: CDCL implication graph & Implication graph at a conflict with the learned clause and the backjump level indicated. & must & 4 \\
A-150 & figure & Figure: Stone duality picture & A finite Boolean algebra beside its Stone space with the correspondence arrows drawn. & should & 4 \\
A-151 & figure & Figure: Curry-Howard triangle & Triangle diagram connecting logic, type theory, and category theory with the translations labeled on the edges. & should & 4 \\
A-152 & figure & Figure: type-derivation tree & An STLC typing derivation displayed as a tree beside its ND counterpart. & must & 4 \\
A-153 & figure & Figure: substructural family map & Map from dropped structural rules to the resulting logics, paired with the family table. & must & 4 \\
A-154 & index & Subject index: inverted headings policy & Handbook-grade subject index of 2500--4000 entries with inverted noun-first headings replacing the sentence-shaped auto-titles, tagged inline from birth during migration. & must & 6 \\
A-155 & index & Subject index: sort keys and escaping & Every entry carries an explicit sort key; accented names (G\"odel, L\"ob, L\"owenheim, Lindstr\"om, Schr\"oder) and math-leading heads are escaped so makeindex sorts and renders them correctly. & must & 6 \\
A-156 & index & Subject index: see and see-also layer & Cross-reference layer of see and see-also entries connecting synonyms and neighboring topics across parts. & must & 6 \\
A-157 & index & Symbol index & Index of every registered symbol with page of registration and principal uses, generated from the symbol-registry CSV. & must & 6 \\
A-158 & index & Person index & Index of persons with accented names correctly sorted, linked to the biography capsules where present. & must & 6 \\
A-159 & index & Named-theorem index & Index of all named theorems and lemmas (including the currently unindexable lemmas such as Lindenbaum and both truth lemmas) pointing to statement and annex proof. & must & 6 \\
A-160 & symbol-registry & Symbol-registry CSV as single source & One machine-readable registry (symbol, meaning, scope, precedence, first-use anchor) that generates the Part 1 tables and drives the notation lint. & must & 3 \\
A-161 & symbol-registry & Registry sub-table: classical logic symbols & Connectives, quantifiers, turnstiles, and metavariable conventions for Parts 1--2 with scope column. & must & 4 \\
A-162 & symbol-registry & Registry sub-table: model-theory symbols & Adds the normalized Part III notation: \(\preccurlyeq\), elementary \(\equiv\), Diag/ElDiag, \(S_n(T)\), ultraproducts \(\prod\mathcal M_i/\mathcal U\), and isomorphism \(\cong\). & must & 4 \\
A-163 & symbol-registry & Registry sub-table: computability and arithmetic symbols & Adds \(\leq_m\), \(\leq_T\), \(\varphi_e\), \(\emptyset^{(n)}\), \(\ulcorner\cdot\urcorner\), \(\bar n\), \(\mathsf{Prov}_T\), \(\mathsf{Con}(T)\), and the object-language \(\leq\). & must & 4 \\
A-164 & symbol-registry & Registry sub-table: set-theory symbols & Adds \(V_\alpha\), rank, \(L_\alpha\), \(\mathbb P\), names and evaluations, \(M[G]\), and the forcing relation. & must & 4 \\
A-165 & symbol-registry & Registry sub-table: nonclassical symbols & Adds \(\Vdash\), \(\Box\), \(\Diamond\), \(\mathsf X\)/\(\mathsf U\)/\(\mathsf F\)/\(\mathsf G\), \(K_i\)/\(E_G\)/\(C_G\), linear connectives, matrix data, and the per-system turnstiles the style chapter promises. & must & 4 \\
A-166 & symbol-registry & Precedence completion & Precedence and associativity ranks for every Part X operator (modal, temporal, epistemic, linear) extending the seven classical ranks, licensing the \(\mathsf U\)-unfolding axiom. & must & 3 \\
A-167 & symbol-registry & Collision ruling: absolute-value bars & \(\lvert\cdot\rvert\) means cardinality only; structure domains are written \(\operatorname{dom}\mathcal M\), closing the Part III/IV seam and the nested-bars defect. & must & 3 \\
A-168 & symbol-registry & Collision ruling: equivalence signs & \(\equiv\) means elementary equivalence only; formula equivalence is \(\Leftrightarrow\); the definitional sign is \(:\equiv\). & must & 3 \\
A-169 & symbol-registry & Collision ruling: forcing macro & One registered forcing macro rendering \(\Vdash\) everywhere, ending the body-versus-feedback drift between the two current macros. & must & 3 \\
A-170 & symbol-registry & Collision ruling: letter overloads & The overloaded letters (formula metavariable versus carrier \(A\), worlds versus valuation \(v\), \(V\), \(D\), \(M\)) are separated by scope entries or renamed. & must & 3 \\
A-171 & symbol-registry & Registry row corrections & Remove the over-scoped \(R^*\) common-knowledge claim, fix the valuation-row category error, and clarify the boldface \(\Sigma^0_n\) pointclass convention. & must & 3 \\
A-172 & symbol-registry & Term-denotation bridging convention & Register once the bridge between the Part III denotation-bracket notation and the Part IV superscript form for term values. & should & 3 \\
A-173 & environment & Entry-card environment with mandatory status key & Theorem-class entry-card environment whose status key (complete proof / proof sketch / cited theorem) is compile-mandatory, with sources field, erroring on omission. & must & 3 \\
A-174 & environment & Numbered warning environment & Numbered, labeled, indexable warning environment replacing prose warnings. & must & 3 \\
A-175 & environment & Numbered proofboundary environment & Numbered internal-proof-boundary environment replacing the unreferenceable paragraph device. & must & 3 \\
A-176 & environment & Audit-table generator & Script emitting the per-part audit longtables (with proper continuation heads and feet) from entry-card metadata; audit tables are generated, never authored. & must & 3 \\
A-177 & environment & Calculus-registry rule display macro & Named, labeled rule-display macro used only inside registry files; body, entry cards, and rule sheets all input the registry, making body-sheet divergence impossible. & must & 3 \\
A-178 & environment & Numbered convention environment book-wide & The convention environment numbered and used from Part 1 onward, promoting the precedence and quantifier-scope conventions out of prose. & must & 3 \\
A-179 & environment & Uniform solution environment & All exercise solutions use the defined solution environment (replacing the 112 hand-rolled italic headers) with graded solution/sketch variants. & must & 3 \\
A-180 & bibliography & Biblatex conversion & Convert the 64 manual bibitems to a biblatex database with normalized DOI/URL practice and stub keys for planned content. & must & 3 \\
A-181 & bibliography & Five missing primaries & Add Glivenko 1929, G\"odel 1933, Gentzen 1936, Kleene's K3 source, and Priest 1979, cited at their points of use. & must & 3 \\
A-182 & bibliography & Per-domain shopping-list intake & Canonical-source lists per new domain (including Martelli--Montanari, CDCL sources, Sahlqvist, Solovay, Paris--Harrington) entered as the domains are written. & must & 4 \\
A-183 & navigation & tocdepth and numbered rule-sheet sections & Set tocdepth to show subsections and number the rule-sheet sections so every sheet is referenceable from the body. & must & 3 \\
A-184 & navigation & List of tables & A list of tables covering every register and comparison table in the handbook. & must & 6 \\
A-185 & navigation & List of figures and listings & Lists of figures and of captioned listings, absent entirely from v1.0. & must & 6 \\
A-186 & navigation & Per-part further reading & A further-reading section per part drawn from the per-domain shopping lists. & must & 4 \\
A-187 & navigation & Glossary & Alphabetical glossary of registered terms with one-sentence definitions and pointers to the defining entries. & must & 6 \\
A-188 & navigation & Exercise-solution bidirectional links & Every exercise hyperlinks to its solution and every solution back to its exercise, preserving the 112=112 completeness discipline. & must & 4 \\
A-189 & navigation & Body-to-appendix cross-references & Body chapters cite the rule sheets, precedence table, and notation registry, and forward references connect LK to cut elimination and resolution to automated reasoning. & should & 4 \\
A-190 & check & Check: status coverage & Every theorem-class environment carries a status key; coverage must be 100 percent and the build fails otherwise. & must & 3 \\
A-191 & check & Check: unlabeled-environment scan & No numbered result environment without a label; no result-bearing sectioning paragraphs. & must & 3 \\
A-192 & check & Check: rule-display ban & No inference-rule display outside the calculus-registry files. & must & 3 \\
A-193 & check & Check: symbol-registry coverage & Notation lint verifying every registered glyph has one role in scope and no body glyph is missing from the registry. & must & 3 \\
A-194 & check & Check: body-annex bidirectional completeness & Every body entry has its annex proof or citation and every annex item has a body entry, generalizing the exercise-feedback discipline. & must & 3 \\
A-195 & check & Check: zero-citation chapters & No chapter with zero citations, enforcing the primary-source rule. & must & 3 \\
A-196 & check & Check: index health & Every index entry has a sort key, no raw-math or sentence-shaped headings, and accented heads escape correctly under makeindex. & must & 3 \\
A-197 & check & Check: dead labels and dangling references & No dangling cross-reference or citation targets, unreferenced labels reported, and prose-relative pointers (above, below, the cited theorem) flagged for conversion to references. & must & 3 \\
A-198 & check & Check: bibliography hygiene & Cited keys and bibliography entries match bidirectionally with consistent DOI/URL fields. & must & 3 \\
A-199 & check & Check: compile warnings as errors & latexmk build treating the declared warning subset (overfull tables, missing characters, undefined references) as errors. & must & 3 \\
A-200 & check & Check: audit-table freshness & Per-part audit tables regenerate cleanly and match entry-card metadata at every session close. & must & 3 \\
A-201 & table & Notation-history table (Part 13) & One row per major notation lineage (Peano--Russell dots, Polish, Gentzen sequent, Hilbert school) with period, exemplar glyphs, and modern descendants; feeds the Part-1 notation-variants table & could & 6 \\
\end{longtable}
\endgroup


\section{Standards and rulings}
\label{sec:rulings}

Binding for the new tree; enforced by the check suite
(\texttt{tools/check.py --strict}).

\begin{enumerate}[label=\textbf{Ruling \arabic*.},ref=\arabic*,leftmargin=2.2em]
  \item \label{ruling:status}\textbf{Status is compile-mandatory.}
        Every theorem-class environment takes a status key
        (\texttt{proved}/\texttt{sketched}/\texttt{cited}) with an annex
        address or source; omission is a compile error.  The body
        contains no \texttt{proof} environments at all --- proofs live
        only in the annex.
  \item \textbf{LK nomenclature.}  Canonical \textsf{LK} is the
        Parts~II--III presentation (primitive negation rules, contexted
        initial sequents) --- the one the rule sheets and completeness
        proofs already use.  Part~VIII's defined-negation variant is
        documented as a named variant with an equivalence theorem.
        \textsf{LK}$^=$ (with equality rules) is a separately named
        calculus whose cut-elimination status is its own row.  No
        theorem cites a calculus without its exact registry name.
  \item \textbf{Glyph roles.}  $|\cdot|$ means cardinality only;
        structure domains are written $\mathrm{dom}\,\mathcal M$
        (closing the Part III/IV seam).  $\equiv$ means elementary
        equivalence only; formula equivalence is $\Leftrightarrow$; the
        definitional sign is $:\equiv$.  Forcing is one registered
        macro ($\Vdash$ with a negated partner), ending the
        $\models$/$\Vdash$ drift.  Worlds, matrices, and carriers get
        registered letters distinct from valuations, structures, and
        formulas.
  \item \textbf{Rules live in the registry.}  Every inference rule is
        named, labeled, and displayed only in
        \texttt{registry/calculi/}; every display shows premises,
        conclusion, name, discharge, and side conditions (the book's
        own requirement at v1.0 lines 12518--12525, mechanized).
  \item \textbf{Numbered semantic environments.}  \texttt{warning} and
        \texttt{proofboundary} are numbered, indexed environments;
        \verb|\paragraph| never again does semantic work.
  \item \textbf{Generated audit tables.}  Every part ends with a
        generated register of its numbered results with status and
        dependencies (institutionalizing the v1.0 incompleteness audit
        table at lines 8002--8035).
  \item \label{ruling:annex}\textbf{Bidirectional completeness.}  Every
        \emph{proved} status resolves to exactly one annex item and
        every annex item is pointed to exactly once --- the Exercise
        Feedback chapter's 112=112 discipline, mechanized as a check.
  \item \textbf{Primary sources.}  Historical results cite primary
        sources (with modern expositions alongside); the
        zero-citation-chapter check enforces the floor.
\end{enumerate}

\section{Campaign plan and session protocol}
\label{sec:campaign}

\subsection{Phases}
\begin{center}\small
\begin{tabular}{@{}clp{2.75in}l@{}}
\toprule
\textbf{Ph} & \textbf{Name} & \textbf{Output} & \textbf{Scale} \\
\midrule
0 & Infrastructure & repo, \texttt{v1.0-audit} tag, baseline build, check suite & done \\
1 & Gap report & this document + four registers & done at sign-off \\
2 & Source math repairs & FIX-SOURCE rows executed; tag \texttt{v1.1-verified} & $\sim$25--40 proof sites \\
3 & Skeleton + registries & new tree, templates, calculus/symbol registries, biblatex & +5\% \\
4 & Migration + depth fill & Parts 1--7 body+annex, Part 8/12 cores; one part = one visit & $0.9$--$1.1\times$ \\
5 & New domains & Parts 8--13 completed; deepenings in 3--7 & $+0.6$--$0.9\times$ \\
6 & Master apparatus & Part 14, four indexes, master tables & $+0.1\times$ \\
7 & Release audit & the book passes its own checklist; v2.0 & --- \\
\bottomrule
\end{tabular}
\end{center}

\noindent Projected final size $1.7$--$2.1\times$ the current volume; a
hard ceiling of $2.2\times$ triggers re-scoping review.

\subsection{Dependency register}
(1)~Phase-2 repairs strictly before migration of the affected parts.
(2)~Rulings before registries before any prose in the new tree.
(3)~Calculus registry before any calculus-bearing chapter migrates.
(4)~biblatex before migration.  (5)~Inline \verb|\index| tagging from
each chapter's birth; the index \emph{pass} only in Phase~6.
(6)~Part~8 entry cards before Parts~9--11 (they cite into them).
(7)~Part-14 master tables after Parts 2--12 land (harvested from entry
fields).  (8)~Part~13 last.

\subsection{Session protocol}
Every later session: (1)~run \texttt{tools/check.py}; (2)~pick open
register rows for the current phase; (3)~\textbf{spec-first rule}: any
scope deviation edits the register row before prose is written;
(4)~work; (5)~flip row status with the commit hash; (6)~regenerate
tables (\texttt{tools/csv2tex.py}); (7)~run checks; (8)~commit.  The
registers alone make any session resumable.

\subsection{Standing controls}
Treatment-level ceilings and the NOT-IN-SCOPE register against scope
explosion; the calculus registry against body--sheet divergence; the
compile-mandatory status key against policy drift (base rate: the v1.0
policy collapsed twice); the mechanical check suite every session;
\texttt{source/} immutable outside Phase~2 so the old book remains
intact and compilable throughout.

\section{Traceability}
\label{sec:traceability}

Finding-to-ID mapping lives in the registers themselves (each row cites
its v1.0 lines); the full audit reports are archived in the campaign
plan file.  During migration each executed row gains the new tree's
label in its status cell, so the specification remains navigable in both
directions for the whole campaign.

\subsection{Frozen v1.0 table of contents}
Parts and chapters of the audited manuscript with their line anchors
(generated from the tagged source):

\begingroup\footnotesize\setlength{\tabcolsep}{3pt}
\begin{longtable}{@{}p{2.1in}p{3.6in}p{0.6in}@{}}
\toprule
\textbf{Part} & \textbf{Chapter} & \textbf{Line} \\
\midrule
\endfirsthead
\toprule
\textbf{Part} & \textbf{Chapter} & \textbf{Line} \\
\midrule
\endhead
\midrule\multicolumn{3}{r}{\emph{continued}}\\\endfoot
\bottomrule\endlastfoot
\textbf{Foundations of Formal Work} &  & 363 \\
 & Sets, Relations, Functions, and Cardinality & 365 \\
 & Levels, Objects, and Consequence & 798 \\
 & Finite Strings, Induction, and Recursion & 1101 \\
 & Writing Formal Definitions and Proofs & 1411 \\
\textbf{Classical Propositional Logic} &  & 1649 \\
 & Propositional Syntax and Semantics & 1658 \\
 & Natural Deduction and Sequents & 2187 \\
 & Soundness, Completeness, Compactness, and Resolution & 2451 \\
\textbf{Classical First-Order Logic} &  & 3022 \\
 & First-Order Syntax & 3032 \\
 & Structures, Satisfaction, and Translation & 3429 \\
 & Theories and Axiomatization & 3729 \\
 & First-Order Deduction & 3927 \\
 & Henkin Completeness and Compactness & 4306 \\
 & Cardinality, Normal Forms, and Resolution & 4968 \\
\textbf{Model Theory} &  & 6049 \\
 & Maps, Substructures, and Elementarity & 6052 \\
 & Types, Saturation, and Ultraproducts & 6510 \\
 & Quantifier Elimination and Categoricity & 7150 \\
\textbf{Computability} &  & 7494 \\
 & Machines and Effective Procedures & 7497 \\
 & Universal Computation, Reductions, and Fixed Points & 7851 \\
 & Logical Decision Problems and the Arithmetical Hierarchy & 8248 \\
\textbf{Formal Arithmetic and Incompleteness} &  & 9080 \\
 & Arithmetic, Representation, and Coding & 9083 \\
 & Diagonalization and the First Incompleteness Theorem & 9439 \\
 & Provability, the Second Incompleteness Theorem, and Truth & 9750 \\
\textbf{Axiomatic Set Theory} &  & 10060 \\
 & The Axioms of ZF and ZFC & 10063 \\
 & Hierarchy, Ordinals, and Cardinals & 10353 \\
 & Constructibility and Forcing & 10675 \\
\textbf{Proof Theory} &  & 11354 \\
 & Classical and Intuitionistic Calculi & 11357 \\
 & Cut Elimination and Analytic Derivations & 11711 \\
 & Normalization and Ordinal Analysis & 11961 \\
\textbf{Second-Order and Higher-Order Logic} &  & 12249 \\
 & Second-Order Syntax and Two Semantics & 12252 \\
 & Simple Type Theory & 12721 \\
\textbf{Nonclassical Logic} &  & 12970 \\
 & Intuitionistic Logic & 12973 \\
 & Modal, Temporal, and Epistemic Logic & 13201 \\
 & Many-Valued and Paraconsistent Logic & 13584 \\
 & Relevance and Linear Logic & 13789 \\
\textbf{Formal Practice, Automation, and Verification} &  & 14095 \\
 & Constructing Derivations, Models, and Countermodels & 14097 \\
 & Automated Reasoning & 14271 \\
 & Mechanized Proof Checking & 14591 \\
 & Writing and Typesetting Formal Logic & 14729 \\
 & Research Workflow and Reproducibility & 15030 \\
 & Registered Notation & 15095 \\
 & Binding and Precedence & 15207 \\
 & Rule Sheets & 15240 \\
 & Qualifying-Examination Problems and Selected Solutions & 15489 \\
 & Exercise Feedback & 15638 \\
 & Editorial and Formal-Logical Release Checklist & 16743 \\
\end{longtable}
\endgroup


\end{document}
